%<*driver>
\def\nameofplainTeX{plain}
\ifx\fmtname\nameofplainTeX
\input l3docstrip.tex
\keepsilent
\askforoverwritefalse
\preamble

========================================================================
 osglecture - profile-driven lecture documents

 (C) 2022-2026 Matthias Werner
 E-mail: matthias.werner@informatik.tu-chemnitz.de

 Released under the LaTeX Project Public License v1.3c or later
 See http://www.latex-project.org/lppl.txt
========================================================================

\endpreamble
\postamble

 This work is "maintained" (as per LPPL maintenance status) by
 Matthias Werner.
\endpostamble

\generate{
  \file{osglecture.cls}{\from{osglecture.dtx}{class}}
  \file{osglecture-config.sty}{\from{osglecture.dtx}{config}}
  \file{osglecture-project.sty}{\from{osglecture.dtx}{project}}
  \file{osglecture-metadata.sty}{\from{osglecture.dtx}{metadata}}
  \file{osglecture-references.sty}{\from{osglecture.dtx}{references}}
  \file{osglecture-continuation.sty}{\from{osglecture.dtx}{continuation}}
  \file{osglecture-integration.sty}{\from{osglecture.dtx}{integration}}
  \file{osglecture-structure.sty}{\from{osglecture.dtx}{structure}}
  \file{osglecture-adapters.sty}{\from{osglecture.dtx}{adapters}}
  \file{osglecture-profiles.sty}{\from{osglecture.dtx}{profiles}}
  \file{osglecture-presitemize.sty}{\from{osglecture.dtx}{presitemize}}
  \file{osglecture-twocolumns.sty}{\from{osglecture.dtx}{twocolumns}}
  \file{osglecture-series.sty}{\from{osglecture.dtx}{series}}
  \file{osglecture-osgbeamer.code.tex}{\from{osglecture.dtx}{osgbeamer}}
  \file{osglecture-adapter-beamer.def}{\from{osglecture-adapters.dtx}{beamer}}
  \file{osglecture-adapter-book.def}{\from{osglecture-adapters.dtx}{book}}
  \file{osglecture-adapter-ltx-talk.def}{\from{osglecture-adapters.dtx}{ltx-talk}}
  \file{osglecture-adapter-scrbook.def}{\from{osglecture-adapters.dtx}{scrbook}}
  \file{osglecture-profile-beamer.def}{\from{osglecture-profiles.dtx}{beamer}}
  \file{osglecture-profile-book.def}{\from{osglecture-profiles.dtx}{book}}
  \file{osglecture-profile-ltx-talk.def}{\from{osglecture-profiles.dtx}{ltx-talk}}
  \file{osglecture-profile-scrbook.def}{\from{osglecture-profiles.dtx}{scrbook}}
}
\expandafter\endbatchfile
\fi
%</driver>
%<*class>
\NeedsTeXFormat{LaTeX2e}[2022/06/01]
\ProvidesExplClass
  {osglecture}
  {2026/08/21}
  {v0.13.0}
  {Profile-driven lecture documents}
%</class>
%<*driver>
\DocumentMetadata{}
\documentclass[
load-preamble+
]{osgdoc}

\usepackage[
  languages={de,en},
  map={de=german,en=english},
  targetlang={job=2},
  load babel
]{langselect}
\usepackage{microtype}
\usepackage{booktabs}
\usepackage{enumitem}

\makeindex
\setcnltx{
  name    = osglecture,
  class   = osglecture,
  version = \UseVersionOf{osglecture.cls},
  date    = \UseDateOf{osglecture.cls},
  title   = {\ldeen{Die \pkg*{osglecture}-Klasse}{The \pkg*{osglecture}\ class}},
  info    = {\ldeen{Metaklasse für Vorlesungsdokumente}{Meta class for lecture documents}},
  authors = {Matthias Werner[matthias.werner@informatik.tu-chemnitz.de]},
  abstract = {
    \ldeen*{
      Die Klasse @1 ist die \LuaLaTeX-Kernklasse des \pkg*{osglecture}-Bundles. Sie
      projiziert ein gemeinsame Lehrquellendokument auf 
      auf Präsentationen und Langformen (Skript, Buch).
    }{
      THe class @1 is the \LuaLaTeX core class of the \pkg*{osglecture} bundle. It
      projects a shared teaching source onto
      presentations and long-form documents (script or book)
    }{\cls*{osglecture}}
  },
  url = https://github.com/werner-matthias/osglecture,
  build-title
}

%\OnlyDescription

\begin{document}
\ldeen{
  Diese Dokumentation enthält nur eine Beschreibung der Klassenoptionen und die Implementation.
  Für die Nutzung von \pkg*{osglecture} siehe \code{osglecture-manual-de.pdf}.}{
  This documentation contains only a description of the class options and its implementation.
  For information on using \pkg*{osglecture}, see \code{osglecture-manual-en.pdf}.
}
\section{\ldeen{Klassenoptionen}{Class options}}

\begin{options}
  \keybool{standalone}\Default{false}
  \ldeen{Nutzung der Klasse außerhalb einer (Veranstaltungs-)Serie.}{Using the class outside of a (course) series}.
  \keyval{doctype}{\ldeen{Typ}{type}}
  \ldeen{Setzt den Zieltyp des Dokuments, wodurch u.\,a.\ die tatsächliche Basisklasse bestimmt wird.}{Sets the document's target type, which, among other things, determines its actual base class.}
  (\ldeen{Dies ist in der Regel nur bei \option{standalone} nötig, sonst sollte es über \code{ollm} gesteuert werden.}{This is usually only necessary with \option{standalone}; otherwise, it should be controlled using \code{ollm}.})
  \keyval{profile}{\ldeen{Profil}{profile}}
  \ldeen{Überschreibt das zum Doctype gehörende Dokumentprofil (legt u.\,a.\ Basisklasse und Backend fest); wird sonst über Konfiguration geregelt.}{Overrides the document profile tied to the doctype (sets, among other things, the base class and backend); otherwise controlled via configuration.}
  \keyval{class-options}{\ldeen{Ausdrücke}{expressions}}
  \ldeen{Setzt Optionen in der Basisklasse}{Sets options of the base class}.
  \keyval{lang}{\ldeen{Sprachcode}{language code}}
  \ldeen{Wählt die Sprache des Zieldokuments aus.}{Selects the language of the target document.}
  (\ldeen{Dies ist in der Regel nur bei \option{standalone} nötig, sonst sollte es über \code{ollm} gesteuert werden.}{This is usually only necessary with \option{standalone}; otherwise, it should be controlled using \code{ollm}.})
  \keybool{patch-overlays}\Default{true}
  \ldeen{Erweitert die Overlay-Syntax von \pkg*{ltx-talk} oder \pkg*{beamer} um portable Modusausdrücke.}{Extends the overlay syntax of \pkg*{ltx-talk} or \pkg*{beamer} to include portable mode expressions.}
\end{options}
\section{Implementation}
\MaybeStop{\ldeen*{Um die Implementationsbeschreibung zu erhalten, kommentieren Sie bitte \cs{OnlyDescription} in \code{osglecture.dtx} aus.}{To view the implementation description, please uncomment \cs{OnlyDescription} in \code{osglecture.dtx}.}}

\subsection{\code{osglecture.cls}}
\AutoDocInput{osglecture.dtx}{class}
\subsection{\code{osglecture-config.sty}}
\AutoDocInput{osglecture.dtx}{config}
\subsection{\code{osglecture-project.sty}}
\AutoDocInput{osglecture.dtx}{project}
\subsection{\code{osglecture-metadata.sty}}
\AutoDocInput{osglecture.dtx}{metadata}
\subsection{\code{osglecture-references.sty}}
\AutoDocInput{osglecture.dtx}{references}
\subsection{\code{osglecture-continuation.sty}}
\AutoDocInput{osglecture.dtx}{continuation}
\subsection{\code{osglecture-integration.sty}}
\AutoDocInput{osglecture.dtx}{integration}
\subsection{\code{osglecture-structure.sty}}
\AutoDocInput{osglecture.dtx}{structure}
\subsection{\code{osglecture-adapters.sty}}
\AutoDocInput{osglecture.dtx}{adapters}
\subsection{\code{osglecture-profiles.sty}}
\AutoDocInput{osglecture.dtx}{profiles}
\subsection{\code{osglecture-presitemize.sty}}
\AutoDocInput{osglecture.dtx}{presitemize}
\subsection{\code{osglecture-twocolumns.sty}}
\AutoDocInput{osglecture.dtx}{twocolumns}
\subsection{\code{osglecture-series.sty}}
\AutoDocInput{osglecture.dtx}{series}
\subsection{\code{osglecture-osgbeamer.code.tex}}
\ldeen{Unpublizierte Legacy-Implementierung; keine automatische
  Kommentar-Dokumentation.}{Unpublished legacy implementation; not
  auto-documented.}

\subsection{\code{osglecture-adapter-beamer.def}}
\AutoDocInput{osglecture-adapters.dtx}{beamer}
\subsection{\code{osglecture-adapter-book.def}}
\AutoDocInput{osglecture-adapters.dtx}{book}
\subsection{\code{osglecture-adapter-ltx-talk.def}}
\AutoDocInput{osglecture-adapters.dtx}{ltx-talk}
\subsection{\code{osglecture-adapter-scrbook.def}}
\AutoDocInput{osglecture-adapters.dtx}{scrbook}

\subsection{\code{osglecture-profile-beamer.def}}
\AutoDocInput{osglecture-profiles.dtx}{beamer}
\subsection{\code{osglecture-profile-book.def}}
\AutoDocInput{osglecture-profiles.dtx}{book}
\subsection{\code{osglecture-profile-ltx-talk.def}}
\AutoDocInput{osglecture-profiles.dtx}{ltx-talk}
\subsection{\code{osglecture-profile-scrbook.def}}
\AutoDocInput{osglecture-profiles.dtx}{scrbook}

\end{document}
%</driver>

%<*config>
\NeedsTeXFormat{LaTeX2e}[2022/06/01]
\ProvidesExplPackage
  {osglecture-config}
  {2026/08/21}
  {v0.13.0}
  {Build-request reader for osglecture}
%</config>
%<*project>
\NeedsTeXFormat{LaTeX2e}[2022/06/01]
\ProvidesExplPackage
  {osglecture-project}
  {2026/08/21}
  {v0.13.0}
  {Deferred project configuration for osglecture}
%</project>
%<*metadata>
\NeedsTeXFormat{LaTeX2e}[2022/06/01]
\ProvidesExplPackage
  {osglecture-metadata}
  {2026/08/21}
  {v0.13.0}
  {Semantic and extensible metadata for osglecture}
%</metadata>
%<*references>
\NeedsTeXFormat{LaTeX2e}[2024/06/01]
\ProvidesExplPackage
  {osglecture-references}
  {2026/08/21}
  {v0.13.0}
  {Logical cross-unit references for osglecture}
%</references>
%<*continuation>
\NeedsTeXFormat{LaTeX2e}[2024/06/01]
\ProvidesExplPackage
  {osglecture-continuation}
  {2026/08/21}
  {v0.13.0}
  {Mode-aware counter continuation between series units}
%</continuation>
%<*integration>
\NeedsTeXFormat{LaTeX2e}[2024/06/01]
\ProvidesExplPackage
  {osglecture-integration}
  {2026/08/21}
  {v0.13.0}
  {Logical-unit PDF integration through tagpax}
%</integration>
%<*structure>
\NeedsTeXFormat{LaTeX2e}[2024/06/01]
\ProvidesExplPackage
  {osglecture-structure}
  {2026/08/21}
  {v0.13.0}
  {Logical unit structure and build-result export for osglecture}
%</structure>
%<*adapters>
\NeedsTeXFormat{LaTeX2e}[2024/06/01]
\ProvidesExplPackage
  {osglecture-adapters}
  {2026/08/21}
  {v0.13.0}
  {Backend adapter loader for osglecture}
%</adapters>
%<*profiles>
\NeedsTeXFormat{LaTeX2e}[2022/06/01]
\ProvidesExplPackage
  {osglecture-profiles}
  {2026/08/21}
  {v0.13.0}
  {Profile registry for osglecture}
%</profiles>
%<*presitemize>
\ProvidesExplPackage
  {osglecture-presitemize}
  {2026/08/21}
  {v0.13.0}
  {Prose-aware presentation lists for osglecture}
%</presitemize>
%<*twocolumns>
\ProvidesExplPackage
  {osglecture-twocolumns}
  {2026/08/21}
  {v0.13.0}
  {Mode-aware two-column layouts for osglecture}
%</twocolumns>
%<*series>
\NeedsTeXFormat{LaTeX2e}[2022/06/01]
\ProvidesExplPackage
  {osglecture-series}
  {2026/08/21}
  {v0.13.0}
  {Series index bridge for osglecture}
%</series>
%<*class>

% \ldeen*{Die legacy Implementierung bleibt vorerst über die
% Klassenoption @1 weiterhin erreichbar.}{
% The legacy implementation remains reachable via the @1 class
% option.}{\option{osgbeamer}}
\@ifclasswith{osglecture}{osgbeamer}
  {
    \input{osglecture-osgbeamer.code.tex}
    \endinput
  }
  { }

\RequirePackage{iftex}
\ifLuaTeX\else
  \ClassError{osglecture}
    {LuaLaTeX~is~required}
    {Run~this~document~with~LuaLaTeX.}
\fi

\RequirePackage{osglecture-config}
\RequirePackage{osglecture-profiles}
\RequirePackage{osglecture-project}

\ExplSyntaxOn

\tl_new:N \l__osglecture_doctype_tl
\tl_new:N \l__osglecture_document_profile_tl
\tl_new:N \l__osglecture_target_profile_class_tl
\tl_new:N \l__osglecture_backend_tl
\tl_new:N \l__osglecture_adapter_tl
\tl_new:N \l__osglecture_base_class_tl
\tl_new:N \l__osglecture_language_tl
\tl_new:N \l__osglecture_base_options_tl
\tl_new:N \l__osglecture_selected_options_tl
\tl_new:N \l__osglecture_setup_file_tl
\tl_new:N \l__osglecture_mode_setup_file_tl
\tl_new:N \l__osglecture_additional_modes_tl
\tl_new:N \l__osglecture_project_config_tl
\bool_new:N \l__osglecture_doctype_explicit_bool
\bool_new:N \l__osglecture_profile_explicit_bool
\bool_new:N \l__osglecture_language_explicit_bool
\bool_new:N \l__osglecture_standalone_bool
\bool_new:N \l__osglecture_patch_overlays_bool
\bool_set_true:N \l__osglecture_patch_overlays_bool

\msg_new:nnnn { osglecture } { profile-required }
  { Document~type~'#1'~has~no~default~document~profile. }
  { Select~a~profile~which~declares~this~document~type. }
\msg_new:nnnn { osglecture } { profile-conflict }
  { Class~profile~'#1'~conflicts~with~build~profile~'#2'. }
  { Change~the~project~configuration~and~regenerate~the~build~file. }
\msg_new:nnnn { osglecture } { metadata-required }
  { Document~profile~'#1'~requires~\token_to_str:N\DocumentMetadata. }
  { Put~\token_to_str:N\DocumentMetadata\c_left_brace_str\c_right_brace_str~
    before~\token_to_str:N\documentclass~in~the~document~source. }
\msg_new:nnnn { osglecture } { metadata-forbidden }
  { Document~profile~'#1'~does~not~support~\token_to_str:N\DocumentMetadata. }
  { Disable~document~metadata~for~this~target~or~select~another~profile. }
\msg_new:nnnn { osglecture } { metadata-policy-conflict }
  { Build~metadata~policy~'#1'~conflicts~with~profile~'#2'~('#3'). }
  { Change~the~target-specific~document_metadata~policy~or~the~target~profile. }
\msg_new:nnnn { osglecture } { standalone-job-conflict }
  { Standalone~mode~cannot~be~combined~with~an~OLLM~job. }
  { Remove~the~standalone~class~option~or~build~without~OLLM. }
\msg_new:nnnn { osglecture } { build-override }
  { Explicit~class~#1~'#2'~overrides~build~value~'#3'. }
  { This~is~intended~for~debugging;~the~job-bound~BuildSpec~remains~unchanged. }

\keys_define:nn { osglecture }
  {
    osgbeamer .code:n = { },
    standalone .bool_set:N = \l__osglecture_standalone_bool,
    doctype .code:n =
      {
        \tl_set:Nn \l__osglecture_doctype_tl {#1}
        \bool_set_true:N \l__osglecture_doctype_explicit_bool
      },
    profile .code:n =
      {
        \tl_set:Nn \l__osglecture_document_profile_tl {#1}
        \bool_set_true:N \l__osglecture_profile_explicit_bool
      },
    lang .code:n =
      {
        \tl_set:Nn \l__osglecture_language_tl {#1}
        \bool_set_true:N \l__osglecture_language_explicit_bool
      },
    class-options .tl_set:N = \l__osglecture_base_options_tl,
    patch-overlays .bool_set:N =
      \l__osglecture_patch_overlays_bool,
  }

\ProcessKeyOptions [ osglecture ]

\bool_if:NTF \l__osglecture_standalone_bool
  {
    \bool_lazy_or:nnT
      { \cs_if_exist_p:N \OSGLectureJobFile }
      { \cs_if_exist_p:N \OSGLectureProjectManifestFile }
      { \msg_fatal:nn {osglecture}{standalone-job-conflict} }
  }
  { \OsgLectureLoadBuildFileForJob }

\OsgLectureBuildLoadedTF
  {
    \bool_if:NTF \l__osglecture_doctype_explicit_bool
      {
        \tl_set:Ne \l_tmpa_tl { \OsgLectureBuildValue{doctype} }
        \str_if_eq:VVF \l__osglecture_doctype_tl \l_tmpa_tl
          {
            \msg_warning:nnxxx {osglecture}{build-override}
              {doctype}{\l__osglecture_doctype_tl}{\l_tmpa_tl}
          }
      }
      {
        \tl_set:Ne \l__osglecture_doctype_tl
          { \OsgLectureBuildValue{doctype} }
      }
    \bool_if:NTF \l__osglecture_language_explicit_bool
      {
        \tl_set:Ne \l_tmpa_tl { \OsgLectureBuildValue{language} }
        \str_if_eq:VVF \l__osglecture_language_tl \l_tmpa_tl
          {
            \msg_warning:nnxxx {osglecture}{build-override}
              {language}{\l__osglecture_language_tl}{\l_tmpa_tl}
          }
      }
      {
        \tl_set:Ne \l__osglecture_language_tl
          { \OsgLectureBuildValue{language} }
      }
    \tl_set:Ne \l__osglecture_target_profile_class_tl
      { \OsgLectureBuildValue{profile-class} }
  }
  { }

\tl_if_empty:NT \l__osglecture_doctype_tl
  { \tl_set:Nn \l__osglecture_doctype_tl {slides} }
\tl_if_empty:NT \l__osglecture_target_profile_class_tl
  {
    \str_case:VnF \l__osglecture_doctype_tl
      {{slides}{\tl_set:Nn \l__osglecture_target_profile_class_tl {presentation}}
       {handout}{\tl_set:Nn \l__osglecture_target_profile_class_tl {presentation}}}
      {\tl_set:Nn \l__osglecture_target_profile_class_tl {longform}}
  }

% \ldeen{Ort und Name der Projektkonfiguration stammen aus dem
% gemeinsamen Manifest-Modul (siehe ARCHITECTURE.md Abschnitt 12): Es
% liest \code{[project.tex]} direkt aus dem Projektmanifest, das es
% selbst ausgehend vom Arbeitsverzeichnis findet. Ein fehlgeschlagenes
% Laden bleibt still -- die Klasse lädt dann einfach keine
% Projektkonfiguration.}{Location and name of the project configuration
% come from the shared manifest module (see ARCHITECTURE.md section 12):
% it reads \code{[project.tex]} directly from the project manifest,
% which it finds itself starting from the working directory. A failed
% load stays silent -- the class then simply loads no project
% configuration.}
\OsgLectureBuildLoadedTF
  {
    \directlua{ require("osglecture-manifest").load_tex() }
    \str_if_eq:eeT { \OsgLectureManifestAvailableFlag } { 1 }
      {
        \tl_set:Ne \l__osglecture_project_config_tl
          { \OsgLectureManifestProjectConfigPath }
        \file_if_exist:VTF \l__osglecture_project_config_tl
          {
            \osglecture_project_intercept_package_commands:
            \exp_args:NV \file_input:n \l__osglecture_project_config_tl
            \osglecture_project_restore_package_commands:
          }
          { }
      }
  }
  { }
\osglecture_project_suspend_metadata_commands:

\str_if_eq:VnTF \l__osglecture_target_profile_class_tl {presentation}
  { \tl_set:Nn \l_tmpa_tl {presentation-profile} }
  { \tl_set:Nn \l_tmpa_tl {longform-profile} }
\exp_args:NNV \prop_get:NnN
  \g__osglecture_project_policy_prop \l_tmpa_tl \l_tmpb_tl
\bool_set_false:N \l_tmpa_bool
\exp_args:NNV \prop_get:NnNT
  \g__osglecture_project_policy_enforced_prop \l_tmpa_tl \l_tmpd_tl
  { \str_if_eq:VnT \l_tmpd_tl {true} {\bool_set_true:N \l_tmpa_bool} }
\OsgLectureBuildLoadedTF
  {
    \tl_set:Ne \l_tmpc_tl {\OsgLectureBuildValue{target}}
    \exp_args:NNV \prop_get:NnNT
      \g__osglecture_project_target_profile_prop \l_tmpc_tl \l_tmpd_tl
      {
        \bool_if:NTF \l_tmpa_bool
          { \str_if_eq:VVF \l_tmpb_tl \l_tmpd_tl
              { \msg_fatal:nnxxx {osglecture-project}{enforced-conflict}
                  {\l_tmpa_tl}{\l_tmpb_tl}{\l_tmpd_tl} } }
          { \tl_set_eq:NN \l_tmpb_tl \l_tmpd_tl }
      }
  }
  { }
\bool_if:NTF \l__osglecture_profile_explicit_bool
  {
    \exp_args:NNV \prop_get:NnNT
      \g__osglecture_project_policy_enforced_prop \l_tmpa_tl \l_tmpc_tl
      { \str_if_eq:VnT \l_tmpc_tl {true}
          { \str_if_eq:VVF \l__osglecture_document_profile_tl \l_tmpb_tl
              { \msg_fatal:nnxx {osglecture}{profile-conflict}
                  {\l__osglecture_document_profile_tl}{\l_tmpb_tl} } } }
  }
  { \tl_set_eq:NN \l__osglecture_document_profile_tl \l_tmpb_tl }
\tl_if_empty:NT \l__osglecture_document_profile_tl
  { \msg_fatal:nnV {osglecture}{profile-required}\l__osglecture_doctype_tl }

\osglecture_profile_select:VV
  \l__osglecture_document_profile_tl \l__osglecture_doctype_tl
\tl_set:Ne \l__osglecture_backend_tl
  {
    \osglecture_profile_value:Vn
      \l__osglecture_document_profile_tl {backend}
  }
\tl_set:Ne \l__osglecture_adapter_tl
  {
    \osglecture_profile_value:Vn
      \l__osglecture_document_profile_tl {adapter}
  }
\tl_set:Ne \l__osglecture_base_class_tl
  {
    \osglecture_profile_value:Vn
      \l__osglecture_document_profile_tl {class}
  }
\tl_set:Ne \l_tmpa_tl
  {
    \osglecture_profile_value:Vn
      \l__osglecture_document_profile_tl {document-metadata}
  }
\bool_set_false:N \l_tmpa_bool
\cs_if_exist:NT \IfDocumentMetadataTF
  { \IfDocumentMetadataTF {\bool_set_true:N \l_tmpa_bool} { } }
\str_if_eq:VnT \l_tmpa_tl {required}
  {
    \bool_if:NF \l_tmpa_bool
      { \msg_fatal:nnV {osglecture}{metadata-required}
          \l__osglecture_document_profile_tl }
  }
\str_if_eq:VnT \l_tmpa_tl {forbidden}
  { \bool_if:NT \l_tmpa_bool
      { \msg_fatal:nnV {osglecture}{metadata-forbidden}
          \l__osglecture_document_profile_tl } }
\OsgLectureBuildLoadedTF
  {
    \tl_set:Ne \l_tmpb_tl
      {\OsgLectureBuildValue{document-metadata-policy}}
    \bool_lazy_or:nnT
      { \bool_lazy_and_p:nn
          {\str_if_eq_p:Vn \l_tmpa_tl {required}}
          {\str_if_eq_p:Vn \l_tmpb_tl {disabled}} }
      { \bool_lazy_and_p:nn
          {\str_if_eq_p:Vn \l_tmpa_tl {forbidden}}
          {\str_if_eq_p:Vn \l_tmpb_tl {required}} }
      { \msg_fatal:nnxxx {osglecture}{metadata-policy-conflict}
          {\l_tmpb_tl}{\l__osglecture_document_profile_tl}{\l_tmpa_tl} }
  }
  { }
\tl_set:Ne \l__osglecture_selected_options_tl
  {
    \osglecture_profile_class_options:VV
      \l__osglecture_document_profile_tl \l__osglecture_doctype_tl
  }
\tl_set:Ne \l_tmpa_tl { \OsgLectureProjectClassOptions }
\tl_if_empty:NF \l_tmpa_tl
  {
    \tl_if_empty:NF \l__osglecture_selected_options_tl
      { \tl_put_right:Nn \l__osglecture_selected_options_tl {,} }
    \tl_put_right:NV \l__osglecture_selected_options_tl \l_tmpa_tl
      }
\tl_set:Ne \l_tmpa_tl
  { \exp_args:NV \OsgLectureSetupClassOptions \l__osglecture_base_class_tl }
\tl_if_empty:NF \l_tmpa_tl
  {
    \tl_if_empty:NF \l__osglecture_selected_options_tl
      { \tl_put_right:Nn \l__osglecture_selected_options_tl {,} }
    \tl_put_right:NV \l__osglecture_selected_options_tl \l_tmpa_tl
  }
\tl_if_empty:NF \l__osglecture_base_options_tl
  {
    \tl_if_empty:NF \l__osglecture_selected_options_tl
      { \tl_put_right:Nn \l__osglecture_selected_options_tl {,} }
    \tl_put_right:NV \l__osglecture_selected_options_tl
      \l__osglecture_base_options_tl
  }
\tl_if_empty:NF \l__osglecture_selected_options_tl
  {
    \exp_args:NVV \PassOptionsToClass
      \l__osglecture_selected_options_tl \l__osglecture_base_class_tl
  }
% \ldeen*{@1 setzt Leerzeichen (Katcode~32) auf Katcode~9 (ignoriert), nicht
% Katcode~10 wie normales \LaTeX. Legacy-Basisklassen wie KOMA-Skripts @2
% verwenden intern eigene, nicht-@3-basierte Makros (etwa
% \cs{@pr@videpackage}), deren Parametermuster ein Leerzeichen als
% Trenntoken erwarten -- unter Katcode~9 verschwindet dieses Trenntoken beim
% Tokenisieren spurlos, das Muster passt nicht mehr, und die Klasse bricht
% mit einer Kaskade von Folgefehlern ab (dasselbe Symptombild wie bei einer
% früheren \code{kvoptions.sty}-Korruption in dieser Datei). Dieselbe
% \cs{@pr@videpackage}-Prüfung läuft nicht nur bei @4, sondern bei jedem
% \cs{RequirePackage} weiterer legacy-artiger Pakete -- die
% Katcode-10-Phase muss deshalb bis nach dem letzten solchen Aufruf
% andauern; die Wiederherstellung auf Katcode~9 erfolgt daher erst kurz vor
% der erneuten @1-Aktivierung, nicht unmittelbar nach @4. Die
% Katcode-Änderung darf hier nicht in einer \TeX-Gruppe stehen: @4
% verbietet das Laden einer Klasse innerhalb einer Gruppe explizit
% ("Loading a class or package in a group").}
% {@1 sets space (catcode~32) to catcode~9 (ignored), not catcode~10 as
% normal \LaTeX\ does. Legacy base classes such as KOMA-Script's @2 use
% their own, non-@3-based macros internally (e.g. \cs{@pr@videpackage}),
% whose parameter patterns expect a space as a delimiter token -- under
% catcode~9 that delimiter vanishes without a trace during tokenization,
% the pattern no longer matches, and the class aborts with a cascade of
% follow-on errors (the same symptom picture as an earlier
% \code{kvoptions.sty} corruption in this file). The same
% \cs{@pr@videpackage} check runs not only for @4 but for every
% \cs{RequirePackage} of further legacy-style packages -- the catcode-10
% window therefore has to last until after the last such call; restoring
% catcode~9 happens only shortly before @1 is reactivated, not right after
% @4. The catcode change must not be inside a \TeX\ group here: @4
% explicitly forbids loading a class inside a group ("Loading a class or
% package in a group").}
% {\cs{ExplSyntaxOn}}{\code{scrbook}}{expl3}{\cs{LoadClass}}
\char_set_catcode:nn {32}{10}
\exp_args:NV \LoadClass \l__osglecture_base_class_tl

\bool_if:NF \l__osglecture_patch_overlays_bool
  { \PassOptionsToPackage{patch-overlays=false}{osglecture-modes} }
\RequirePackage{osglecture-modes}
\DeclareModeBehaviorClass{presentation-dynamic}
\DeclareModeBehaviorClass{presentation-static}
\DeclareModeBehaviorClass{longform}
\AssignLectureModeBehavior{slides}{presentation-dynamic}
\AssignLectureModeBehavior{handout}{presentation-static}
\AssignLectureModeBehavior{script,article}{longform}
\tl_set:Ne \l__osglecture_mode_setup_file_tl
  {
    \osglecture_profile_value:Vn
      \l__osglecture_document_profile_tl {mode-setup-file}
  }
\tl_if_empty:NF \l__osglecture_mode_setup_file_tl
  { \exp_args:NV \file_input:n \l__osglecture_mode_setup_file_tl }
\str_case:VnF \l__osglecture_doctype_tl
  {
    {script}  { \SetLectureModes{script,article} }
    {article} { \SetLectureModes{article,script} }
  }
  { \exp_args:NV \SetLectureModes \l__osglecture_doctype_tl }
\str_if_eq:VnT \l__osglecture_backend_tl {beamer}
  { \AddLectureModes{beamer} }
\tl_set:Ne \l_tmpa_tl
  {
    \osglecture_profile_value:Vn
      \l__osglecture_document_profile_tl {modes}
  }
\tl_if_empty:NF \l_tmpa_tl { \exp_args:NV \AddLectureModes \l_tmpa_tl }
\tl_if_empty:NF \l__osglecture_additional_modes_tl
  { \exp_args:NV \AddLectureModes \l__osglecture_additional_modes_tl }
\FinalizeLectureModes

\NewExpandableDocumentCommand \OsgLectureDocumentType { }
  { \tl_use:N \l__osglecture_doctype_tl }
\NewExpandableDocumentCommand \OsgLectureDocumentProfile { }
  { \tl_use:N \l__osglecture_document_profile_tl }
\NewExpandableDocumentCommand \OsgLectureBackend { }
  { \tl_use:N \l__osglecture_backend_tl }
\NewExpandableDocumentCommand \OsgLectureAdapter { }
  { \tl_use:N \l__osglecture_adapter_tl }
\NewExpandableDocumentCommand \OsgLectureLanguage { }
  { \tl_use:N \l__osglecture_language_tl }
\NewDocumentCommand \OsgLectureStandaloneTF { m m }
  {
    \bool_if:NTF \l__osglecture_standalone_bool
      {#1}
      { \OsgLectureBuildLoadedTF {#2} {#1} }
  }

% \ldeen*{Ohne \pkg{unicode-math} oder \pkg{lua-unicode-math} überspringt
% @1s eigene Auto-Erkennung (Default \code{auto} von
% \code{math/mathml/luamml/load}) die MathML-Erzeugung für Mathe-Inhalt
% stillschweigend und warnt nur -- leicht zu übersehen, und \pkg{amsmath}
% ist verbreitet genug, dass dies sonst in den meisten getaggten
% Dokumenten die Mathe-Barrierefreiheit unbemerkt beeinträchtigen würde.
% \pkg{luamml} selbst benötigt kein Unicode-Mathe-Font-Paket; das
% erzwungene Einschalten ist der dokumentierte Umweg genau für diesen
% Fall, ohne selbst ein Font-Auswahlpaket (samt dessen optischen
% Nebenwirkungen) zu laden.}{Without \pkg{unicode-math} or
% \pkg{lua-unicode-math}, @1's own auto-detection (default \code{auto} of
% \code{math/mathml/luamml/load}) silently skips MathML generation for
% math content and only warns about it -- easy to miss, and \pkg{amsmath}
% is common enough in practice that this would otherwise quietly degrade
% math accessibility in most tagged documents that use it. \pkg{luamml}
% itself does not require a Unicode math font package; forcing it on is
% the documented bypass for exactly this situation, without loading
% either font-selection package (and their visual side effects)
% ourselves.}{latex-lab}
\cs_if_exist:NT \tagstructbegin
  { \tagpdfsetup{math/mathml/luamml/load=true} }

\RequirePackage{osglecture-series}
\RequirePackage{osglecture-metadata}
\char_set_catcode:nn {32}{9}
\ExplSyntaxOn
\exp_args:Nx \OsgLectureConfigureMetadata
  {
    course-target =
      \osglecture_profile_value:Vn
        \l__osglecture_document_profile_tl {course-target},
    event-target =
      \osglecture_profile_value:Vn
        \l__osglecture_document_profile_tl {event-target}
  }

\RequirePackage{osglecture-presitemize}
\RequirePackage{osglecture-twocolumns}
\ExplSyntaxOn
\RequirePackage{osglecture-adapters}
\exp_args:NV \OsgLectureLoadAdapter \l__osglecture_adapter_tl
\RequirePackage{osglecture-structure}
\RequirePackage{osglecture-references}
\RequirePackage{osglecture-continuation}
\ExplSyntaxOn
\OsgLectureBuildLoadedTF
  {
    \directlua{ require("osglecture-series-index").current_unit_tex() }
    \str_if_eq:eeT { \OsgLectureUnitRole } { i }
      { \RequirePackage{osglecture-integration} }
  }
  { }
\ExplSyntaxOn

\tl_set:Ne \l__osglecture_setup_file_tl
  {
    \osglecture_profile_value:Vn
      \l__osglecture_document_profile_tl {setup-file}
  }
\osglecture_configuration_priority_set:n {10}
\tl_if_empty:NF \l__osglecture_setup_file_tl
  { \exp_args:NV \file_input:n \l__osglecture_setup_file_tl }
\osglecture_configuration_priority_set:n {30}
\osglecture_project_apply_mode_setup:
\exp_args:NV \osglecture_project_include_preambles:n
  \l__osglecture_base_class_tl

\ExplSyntaxOff
%</class>
%<*config>

\ExplSyntaxOn

\prop_new:N \g__osglecture_build_prop
\bool_new:N \g__osglecture_build_loaded_bool

% \ldeen{\code{l3kernel} deklariert nur \cs{l\_tmpa\_tl}/\cs{l\_tmpb\_tl}
% als generische Scratch-\code{tl}-Variablen; \code{c}/\code{d} sind kein
% garantierter Teil des Kernels. \code{osglecture.cls},
% \code{osglecture-project.sty} und \code{osglecture-metadata.sty}
% benötigen wegen verschachtelter Nachschlagevorgänge dennoch vier
% gleichzeitige Scratch-Variablen und verwenden dafür (in Anlehnung an die
% Kernelkonvention) \code{c}/\code{d} zusätzlich. Die Existenzprüfung
% verhindert einen Fehler auf Systemen, auf denen diese Namen bereits durch
% ein anderes geladenes Paket belegt sind.}{\code{l3kernel} only declares
% \cs{l\_tmpa\_tl}/\cs{l\_tmpb\_tl} as generic scratch \code{tl}
% variables; \code{c}/\code{d} are not a guaranteed part of the kernel.
% \code{osglecture.cls}, \code{osglecture-project.sty} and
% \code{osglecture-metadata.sty} nonetheless need four simultaneous scratch
% variables because of nested lookups, and use \code{c}/\code{d} for that
% (following the kernel's own naming convention). The existence check
% avoids an error on systems where these names are already taken by some
% other loaded package.}
\cs_if_exist:NF \l_tmpc_tl { \tl_new:N \l_tmpc_tl }
\cs_if_exist:NF \l_tmpd_tl { \tl_new:N \l_tmpd_tl }

% \ldeen{Defaults für die von \code{osglecture-manifest.lua}s
% \code{manifest.load\_tex} global gesetzten Makros: Ohne diese
% Vorabdeklaration wäre eine Referenz vor dem ersten (erfolgreichen oder
% fehlgeschlagenen) Aufruf eine undefinierte Kontrollsequenz statt eines
% leeren Werts.}{Defaults for the macros globally set by
% \code{osglecture-manifest.lua}'s \code{manifest.load\_tex}: without this
% pre-declaration, a reference before the first call (successful or
% failed) would be an undefined control sequence rather than an empty
% value.}
\cs_new:Npn \OsgLectureManifestAvailableFlag { 0 }
\cs_new:Npn \OsgLectureManifestProjectRoot { }
\cs_new:Npn \OsgLectureManifestSharedTexDirectory { }
\cs_new:Npn \OsgLectureManifestProjectConfigFile { }
\cs_new:Npn \OsgLectureManifestProjectConfigPath { }
\cs_new:Npn \OsgLectureManifestAvailableLanguages { }
\cs_new:Npn \OsgLectureManifestBundlePreset { }

\msg_new:nnnn { osglecture-config } { unsupported-schema }
  { Unsupported~osglecture~build-file~schema~'#1'. }
  { This~version~supports~build-file~schema~'1'. }
\msg_new:nnnn { osglecture-config } { missing-key }
  { The~osglecture~build~file~does~not~define~'#1'. }
  { Regenerate~the~build~file~with~OLLM. }
\msg_new:nnnn { osglecture-config } { job-mismatch }
  { Build~file~job-id~'#1'~does~not~match~TeX~jobname~'#2'. }
  { Use~the~build~file~generated~for~this~jobname. }
\msg_new:nnnn { osglecture-config } { file-not-found }
  { OSG~lecture~build~file~'#1'~was~not~found. }
  { Generate~it~with~OLLM~or~use~standalone~class~options. }
\msg_new:nnnn { osglecture-config } { job-symbol-missing }
  { No~OLLM~job~was~specified~for~this~series~build. }
  { Start~the~build~with~OLLM~or~select~the~class~option~'standalone'. }
\msg_new:nnnn { osglecture-config } { manifest-symbol-missing }
  { No~project~manifest~was~specified~for~this~series~build. }
  { Start~the~build~with~OLLM~or~select~the~class~option~'standalone'. }
\msg_new:nnnn { osglecture-config } { manifest-not-found }
  { The~project~manifest~'#1'~was~not~found. }
  { Restore~the~series~manifest~or~select~the~class~option~'standalone'. }

\cs_new_protected:Npn \__osglecture_build_put:nn #1#2
  { \prop_gput:Nnn \g__osglecture_build_prop {#1} {#2} }

\keys_define:nn { osglecture/build }
  {
    schema .code:n =
      {
        \int_compare:nNnTF {#1} = {1}
          { \__osglecture_build_put:nn {schema} {#1} }
          { \msg_error:nnn {osglecture-config}{unsupported-schema}{#1} }
      },
    job-id              .code:n = \__osglecture_build_put:nn {job-id} {#1},
    context             .code:n = \__osglecture_build_put:nn {context} {#1},
    series-id           .code:n = \__osglecture_build_put:nn {series-id} {#1},
    unit-id             .code:n = \__osglecture_build_put:nn {unit-id} {#1},
    generation-id       .code:n =
      \__osglecture_build_put:nn {generation-id} {#1},
    target              .code:n = \__osglecture_build_put:nn {target} {#1},
    profile-class       .code:n = \__osglecture_build_put:nn {profile-class} {#1},
    document-metadata-policy .choices:nn = { required, disabled }
      { \__osglecture_build_put:nn {document-metadata-policy} {#1} },
    doctype             .code:n = \__osglecture_build_put:nn {doctype} {#1},
    applicable-unit-scopes .code:n =
      \__osglecture_build_put:nn {applicable-unit-scopes} {#1},
    language            .code:n = \__osglecture_build_put:nn {language} {#1},
    shell-escape        .code:n =
      \__osglecture_build_put:nn {shell-escape} {#1},
    structure-signature .code:n =
      \__osglecture_build_put:nn {structure-signature} {#1},
    config-signature    .code:n =
      \__osglecture_build_put:nn {config-signature} {#1},
  }

% \ldeen{Die Pflichtschlüssel entsprechen genau der Klassifikation
% "OLLM"/"Grenzfall" in \file{ARCHITECTURE.md} Abschnitt 12. Als
% "gemeinsames Modul" klassifizierte Schlüssel haben hier bewusst
% keinen \code{.code:n}-Eintrag: Eine Build-Datei, die einen solchen
% Schlüssel trotzdem sendet, erhält dadurch einen klaren
% "unknown key"-Fehler, statt stillschweigend ignoriert zu
% werden.}{The required keys match exactly the "OLLM"/"Grenzfall"
% classification in \file{ARCHITECTURE.md} section 12. Keys classified
% as "gemeinsames Modul" deliberately have no \code{.code:n} entry here:
% a build file that sends such a key regardless gets a clear
% "unknown key" error instead of being silently ignored.}
\cs_new_protected:Npn \__osglecture_build_validate:
  {
    \clist_map_inline:nn
      {schema,job-id,context,series-id,target,profile-class,
       document-metadata-policy,
       doctype,applicable-unit-scopes,language,
       structure-signature,
       config-signature}
      {
        \prop_if_in:NnF \g__osglecture_build_prop {##1}
          { \msg_error:nnn {osglecture-config}{missing-key}{##1} }
      }
    \prop_get:NnN \g__osglecture_build_prop {job-id} \l_tmpa_tl
    \str_if_eq:eeF {\l_tmpa_tl} {\c_sys_jobname_str}
      {
        \msg_error:nnee {osglecture-config}{job-mismatch}
          {\l_tmpa_tl}{\c_sys_jobname_str}
      }
    \bool_gset_true:N \g__osglecture_build_loaded_bool
  }

\NewDocumentCommand \OsgLectureBuildSetup { m }
  {
    \prop_gclear:N \g__osglecture_build_prop
    \bool_gset_false:N \g__osglecture_build_loaded_bool
    \keys_set:nn {osglecture/build} {#1}
    \__osglecture_build_validate:
  }

\NewExpandableDocumentCommand \OsgLectureBuildValue { m }
  { \prop_item:Nn \g__osglecture_build_prop {#1} }

\NewDocumentCommand \OsgLectureBuildLoadedTF { m m }
  { \bool_if:NTF \g__osglecture_build_loaded_bool {#1} {#2} }

% ExplSyntaxOn sets catcode 126 (the tilde) to 10 (space), since expl3
% source uses a tilde as a visible-space marker. The build file is read
% while this class stays in expl3 syntax, so a tilde in a path value (e.g.
% a Windows short/8.3 directory name such as "RUNNER~1") would silently
% become a plain space at tokenization time -- unrecoverable by any later
% escaping, since the character identity is lost, not just its catcode.
% Setting it back to 12 (other) for the read makes it an inert character,
% safe to store and pass on to TeX or Lua code afterwards.
\NewDocumentCommand \OsgLectureLoadBuildFile { m }
  {
    \file_if_exist:nTF {#1}
      {
        \group_begin:
        \char_set_catcode:nn {126} {12}
        \file_input:n {#1}
        \group_end:
      }
      { \msg_error:nnn {osglecture-config}{file-not-found}{#1} }
  }

\NewDocumentCommand \OsgLectureLoadBuildFileForJob { }
  {
    \cs_if_exist:NTF \OSGLectureProjectManifestFile
      {
        \exp_args:Ne \file_if_exist:nTF
          { \OSGLectureProjectManifestFile }
          { }
          {
            \msg_fatal:nnx {osglecture-config}{manifest-not-found}
              { \OSGLectureProjectManifestFile }
          }
      }
      { \msg_fatal:nn {osglecture-config}{manifest-symbol-missing} }
    \cs_if_exist:NTF \OSGLectureJobFile
      { \exp_args:Ne \OsgLectureLoadBuildFile { \OSGLectureJobFile } }
      { \msg_fatal:nn {osglecture-config}{job-symbol-missing} }
  }

\ExplSyntaxOff
%</config>
%<*project>

\ExplSyntaxOn

\clist_new:N \g__osglecture_metadata_fields_clist
\seq_new:N \g__osglecture_metadata_pending_seq
\int_new:N \g__osglecture_configuration_priority_int
\bool_new:N \g__osglecture_configuration_enforced_bool
\tl_new:N \g__osglecture_project_class_options_tl
\prop_new:N \g__osglecture_project_policy_prop
\prop_new:N \g__osglecture_project_policy_priority_prop
\prop_new:N \g__osglecture_project_policy_enforced_prop
\prop_new:N \g__osglecture_project_target_profile_prop
\prop_new:N \g__osglecture_project_target_filecolor_prop
\prop_new:N \g__osglecture_setup_class_prop
\prop_new:N \g__osglecture_setup_mode_prop
\seq_new:N \g__osglecture_project_preamble_seq
\tl_new:N \l__osglecture_project_target_tl
\tl_new:N \l__osglecture_preamble_class_tl
\tl_new:N \l__osglecture_preamble_current_class_tl
\bool_new:N \l__osglecture_preamble_include_bool
\msg_new:nnnn {osglecture-project}{enforced-conflict}
  { Configuration~'#1'~is~enforced~as~'#2';~cannot~set~'#3'. }
  { Change~the~project~enforcement~instead~of~overriding~it~locally. }
\msg_new:nnnn {osglecture-project}{preamble-not-found}
  { Preamble~fragment~'#1'~was~not~found. }
  { Put~'#1.tex'~in~the~shared~TeX~directory~or~install~
    'osglecture-preamble-#1.tex'~in~a~TeX~tree. }
\msg_new:nnnn {osglecture-project}{invalid-setup-key}
  { Unknown~osglecture~setup~key~'#1'. }
  { Use~an~explicit~'class/<name>'~or~'mode/<name>'~key. }
\msg_new:nnnn {osglecture-project}{preamble-star-condition}
  {
    \string\IncludeOsgLecturePreamble*~loads~'#1'~immediately~and~cannot~
    evaluate~a~mode~or~class~condition~at~that~point.
  }
  {
    Drop~the~condition,~or~drop~the~star~to~defer~the~fragment~until~after~
    the~base~class,~mode~graph,~and~adapter~are~known.
  }
\msg_new:nnnn {osglecture-project}{unknown-setup-class}
  { Unknown~base~class~'#1'~in~osglecture~setup. }
  { Load~or~declare~a~document~profile~for~that~class~before~using~its~setup. }
\cs_new_protected:Npn \__osglecture_metadata_receiver:nnnnn #1#2#3#4#5
  {
    \seq_gput_right:Nx \g__osglecture_metadata_pending_seq
      {
        \exp_not:N \__osglecture_metadata_record:nnnnn
        {\exp_not:n{#1}}{\exp_not:n{#2}}{\exp_not:n{#3}}{\exp_not:n{#4}}
        {#5}
      }
  }

% \ldeen*{@1 wird während des frühen Lesens von \file{projectconfig.tex}
% installiert (siehe @2 in \file{osglecture.cls}), \emph{bevor} die
% Zielbasisklasse geladen ist. Ein Paket, das dort direkt geladen wird und
% selbst prüft, ob die Basisklasse bereits etwas definiert hat, sieht
% einen falschen Zustand -- @3 dokumentiert einen konkreten Fall
% (\pkg{langselect}s \pkg{csquotes}-Integration kollidiert mit
% \pkg*{ltx-talk}s eigener \env{quote}-Umgebung). @4 und @5 fangen
% \cs{usepackage}/\cs{RequirePackage}-Aufrufe während dieses Fensters
% deshalb genauso ab wie die vorhandenen Metadatenfelder: Sie werden in
% dieselbe, reihenfolgeerhaltende Warteschlange @6 eingereiht und erst
% abgespielt, nachdem \cs{LoadClass} gelaufen ist (siehe
% \file{osglecture-metadata.sty}) -- ein Aufruf wie \cs{usepackage} gefolgt
% von \cs{title} mit einem vom Paket bereitgestellten Befehl funktioniert
% dadurch in normaler Leserichtung, ohne dass Nutzer \cs{AtEndOfClass}
% kennen müssen. Das Sicherheitsnetz gilt gezielt für direkt in
% \file{projectconfig.tex} geschriebene Paketladungen; die empfohlene
% Konvention für umfangreicheren Präambelcode bleibt eine ausgelagerte
% Datei über @7.}{@1 is installed while \file{projectconfig.tex} is read
% early (see @2 in \file{osglecture.cls}), \emph{before} the target base
% class has loaded. A package loaded directly there that itself checks
% whether the base class has already defined something sees a false
% state -- @3 documents a concrete case (\pkg{langselect}'s
% \pkg{csquotes} integration collides with \pkg*{ltx-talk}'s own
% \env{quote} environment). @4 and @5 therefore intercept
% \cs{usepackage}/\cs{RequirePackage} calls during this window the same
% way the existing metadata fields already are: they are enqueued into
% the same order-preserving queue @6 and only replayed after
% \cs{LoadClass} has run (see \file{osglecture-metadata.sty}) -- a call
% like \cs{usepackage} followed by \cs{title} using a command the package
% provides then works in normal reading order, with no need for users to
% know about \cs{AtEndOfClass}. This safety net targets package loads
% written directly in \file{projectconfig.tex}; the recommended
% convention for more substantial preamble code remains an outsourced
% file via @7.}{\cs{osglecture\_project\_intercept\_package\_commands:}}
% {\cs{\_\_osglecture\_project\_config\_tl}}{examples/series-modern}
% {\cs{osglecture\_project\_restore\_package\_commands:}}
% {\cs{osglecture\_project\_intercept\_package\_commands:}}
% {\cs{g\_\_osglecture\_metadata\_pending\_seq}}{\cs{IncludeOsgLecturePreamble}}
\cs_new_eq:NN \__osglecture_project_real_RequirePackage \RequirePackage
\cs_new_protected:Npn \__osglecture_project_package_record:nn #1#2
  { \__osglecture_project_real_RequirePackage [#1] {#2} }
\NewDocumentCommand \__osglecture_project_proxy_RequirePackage { O{} m }
  {
    \seq_gput_right:Nx \g__osglecture_metadata_pending_seq
      {
        \exp_not:N \__osglecture_project_package_record:nn
          {\exp_not:n{#1}} {\exp_not:n{#2}}
      }
  }
\cs_new_protected:Npn \osglecture_project_intercept_package_commands:
  {
    \cs_set_eq:NN \RequirePackage \__osglecture_project_proxy_RequirePackage
    \cs_set_eq:NN \usepackage \__osglecture_project_proxy_RequirePackage
  }
\cs_new_protected:Npn \osglecture_project_restore_package_commands:
  {
    \cs_set_eq:NN \RequirePackage \__osglecture_project_real_RequirePackage
    \cs_set_eq:NN \usepackage \__osglecture_project_real_RequirePackage
  }

\int_gset:Nn \g__osglecture_configuration_priority_int {20}
\bool_gset_false:N \g__osglecture_configuration_enforced_bool

\cs_new_protected:Npn \__osglecture_policy_set:nn #1#2
  {
    \prop_get:NnNTF \g__osglecture_project_policy_priority_prop {#1} \l_tmpa_tl
      {
        \prop_get:NnN \g__osglecture_project_policy_prop {#1} \l_tmpb_tl
        \prop_get:NnN \g__osglecture_project_policy_enforced_prop {#1} \l_tmpc_tl
        \int_compare:nNnTF {\g__osglecture_configuration_priority_int} < {\l_tmpa_tl}
          {
            \str_if_eq:VnT \l_tmpc_tl {true}
              { \str_if_eq:VnF \l_tmpb_tl {#2}
                  { \msg_fatal:nnxxx {osglecture-project}{enforced-conflict}
                      {#1}{\l_tmpb_tl}{#2} } }
          }
          { \__osglecture_policy_store:nn {#1}{#2} }
      }
      { \__osglecture_policy_store:nn {#1}{#2} }
  }
\cs_new_protected:Npn \__osglecture_policy_store:nn #1#2
  {
    \prop_gput:Nnn \g__osglecture_project_policy_prop {#1}{#2}
    \prop_gput:Nnx \g__osglecture_project_policy_priority_prop {#1}
      {\int_use:N \g__osglecture_configuration_priority_int}
    \prop_gput:Nnx \g__osglecture_project_policy_enforced_prop {#1}
      {\bool_if:NTF \g__osglecture_configuration_enforced_bool {true}{false}}
  }

% \ldeen*{@1 sammelt die Optionen für @2 ein, ohne sie sofort auszulösen:
% Die einzelnen Unterschlüssel von @3 (siehe unten) hängen jeweils nur ein
% weiteres \meta{Schlüssel}=\marg{Wert}-Fragment an @1 an. Dadurch ist die
% Reihenfolge, in der die Unterschlüssel innerhalb eines einzelnen
% \cs{LectureProjectSetup}-Aufrufs angegeben werden, unerheblich, und es
% entsteht -- anders als bei sofort auslösenden Einzelschlüsseln -- nur
% ein einziger @2-Aufruf. Ein zweiter, späterer @2-Aufruf für dasselbe
% Paket würde von \LaTeX{} für bereits verarbeitete Optionen stillschweigend
% ignoriert, nicht mit einer Fehlermeldung quittiert; das Sammeln vor dem
% Auslösen ist deshalb kein Komfort, sondern verhindert genau diesen
% stillen Optionsverlust.}{@1 collects the options for @2 without
% triggering them immediately: each sub-key of @3 (see below) only
% appends one more \meta{key}=\marg{value} fragment to @1. This makes the
% order in which the sub-keys are given within a single
% \cs{LectureProjectSetup} call irrelevant, and -- unlike immediately
% triggering individual keys -- results in exactly one @2 call. A second,
% later @2 call for the same package would have its options for
% already-processed keys silently ignored by \LaTeX{}, not flagged with an
% error; collecting before triggering therefore is not a convenience but
% prevents exactly that silent option loss.}{\cs{g\_\_osglecture\_language\_setup\_tl}}
% {\cs{RequirePackage}}{\code{osglecture/project/languages}}
\tl_new:N \g__osglecture_language_setup_tl
\cs_new_protected:Npn \__osglecture_language_setup_append:nn #1#2
  {
    \tl_if_empty:NF \g__osglecture_language_setup_tl
      { \tl_gput_right:Nn \g__osglecture_language_setup_tl {,} }
    \tl_gput_right:Nn \g__osglecture_language_setup_tl {#1={#2}}
  }
% \ldeen*{Die Unterschlüssel spiegeln @1s eigene Paketoptionen, mit einer
% bewussten Ausnahme: @2 heißt hier @3, um beim äußeren Projekt-Schlüssel
% @4 nicht \keyis{languages}{languages=\ldots} schreiben zu müssen -- @3
% trifft die Semantik ohnehin genauer (@1 selbst nennt die möglichen
% Zielsprachen in seiner eigenen Dokumentation "`Auswahlsprachen"').}{The
% sub-keys mirror @1's own package options, with one deliberate exception:
% @2 is named @3 here so that the outer project key @4 does not require
% writing \keyis{languages}{languages=\ldots} -- @3 matches the semantics
% more precisely anyway (@1 itself calls the possible target languages
% "`selectable languages"' in its own documentation).}{\pkg*{langselect}}
% {\option{languages}}{\option{selectable}}{\option{languages}}
\keys_define:nn { osglecture/project/languages }
  {
    selectable .code:n =
      { \__osglecture_language_setup_append:nn {languages}{#1} },
    targetlang .code:n =
      { \__osglecture_language_setup_append:nn {targetlang}{#1} },
    prefix .code:n =
      { \__osglecture_language_setup_append:nn {prefix}{#1} },
    auto .code:n =
      { \__osglecture_language_setup_append:nn {auto}{#1} },
    trim .code:n =
      { \__osglecture_language_setup_append:nn {trim}{#1} },
    map .code:n =
      { \__osglecture_language_setup_append:nn {map}{#1} },
    load~babel .code:n =
      { \__osglecture_language_setup_append:nn {load~babel}{#1} },
    load~polyglossia .code:n =
      { \__osglecture_language_setup_append:nn {load~polyglossia}{#1} },
    unified~shorthands .code:n =
      { \__osglecture_language_setup_append:nn {unified~shorthands}{#1} },
  }

\keys_define:nn { osglecture/project }
  {
    class-options .tl_gset:N = \g__osglecture_project_class_options_tl,
    presentation-profile .code:n =
      { \__osglecture_policy_set:nn {presentation-profile}{#1} },
    longform-profile .code:n =
      { \__osglecture_policy_set:nn {longform-profile}{#1} },
    identity-profile .code:n =
      { \__osglecture_policy_set:nn {identity-profile}{#1} },
    theme .code:n = { \__osglecture_policy_set:nn {theme}{#1} },
    numbering .code:n = { \__osglecture_policy_set:nn {numbering}{#1} },
    references .code:n = { \__osglecture_policy_set:nn {references}{#1} },
% \ldeen*{@1 verbirgt, dass die Mehrsprachumsetzung technisch von @2
% geleistet wird -- Projektdateien nennen nur die Unterschlüssel von
% @6 (@7, \code{map}, \code{load babel}, \ldots), ohne den Paketnamen zu
% kennen. Der gesammelte @3-Aufruf profitiert von derselben
% Abfangen-und-nachtragen-Warteschlange wie ein direkt in
% \file{projectconfig.tex} geschriebenes @4, da @1 ausschließlich
% innerhalb dieses Lesefensters (oder in einem verzögert abgespielten
% Präambelfragment nach @5, wo dieselbe Sicherheit ohnehin gilt)
% aufgerufen wird.}{@1 hides that multilingual support is technically
% implemented by @2 -- project files only name the sub-keys of @6 (@7,
% \code{map}, \code{load babel}, \ldots), without knowing the package
% name. The collected @3 call benefits from the same intercept-and-replay
% queue as a @4 written directly in \file{projectconfig.tex}, since @1 is
% only ever called within that reading window (or in a deferred preamble
% fragment after @5, where the same safety already applies
% otherwise).}{\cs{LectureProjectSetup}}{\pkg*{langselect}}
% {\cs{RequirePackage}}{\cs{usepackage}}{\cs{LoadClass}}
% {\code{osglecture/project/languages}}{\option{selectable}}
    languages .code:n =
      {
        \tl_gclear:N \g__osglecture_language_setup_tl
        \keys_set:nn { osglecture/project/languages } {#1}
        \__osglecture_policy_set:nn {languages}{#1}
        \tl_set:Nx \l_tmpa_tl
          {
            \exp_not:N \RequirePackage
              [ \g__osglecture_language_setup_tl ] {langselect}
          }
        \l_tmpa_tl
      },
  }

\NewDocumentCommand \LectureProjectSetup { m }
  { \keys_set:nn { osglecture/project } {#1} }

% \ldeen*{Der äußere Setup-Namensraum ist absichtlich explizit:
% @1 und @2 bleiben eindeutig unterscheidbar, selbst wenn
% \LaTeX-Ökosysteme denselben Namen für Basisklasse und Modus
% wiederverwenden.}{The outer setup namespace is deliberately explicit:
% @1 and @2 stay unambiguous even when \LaTeX\ ecosystems reuse the same
% name for a base class and a mode.}{\code{class/article}}{\code{mode/article}}
\cs_new_protected:Npn \__osglecture_setup_append:Nnn #1#2#3
  {
    \prop_get:NnNTF #1 {#2} \l_tmpa_tl
      { \prop_gput:Nnx #1 {#2} { \exp_not:V \l_tmpa_tl , \exp_not:n {#3} } }
      { \prop_gput:Nnn #1 {#2} {#3} }
  }

\cs_new_protected:Npn \__osglecture_setup_store:nn #1#2
  {
    \regex_extract_once:nnNTF
      { \A (class|mode) / ([^/]+) \Z } {#1} \l_tmpa_seq
      {
        \tl_set:Ne \l_tmpa_tl { \seq_item:Nn \l_tmpa_seq {2} }
        \tl_set:Ne \l_tmpb_tl { \seq_item:Nn \l_tmpa_seq {3} }
        \str_case:Vn \l_tmpa_tl
          {
            {class}
              {
                \exp_args:NV \OsgLectureProfileClassKnownTF \l_tmpb_tl
                  { \exp_args:NNV \__osglecture_setup_append:Nnn
                      \g__osglecture_setup_class_prop \l_tmpb_tl {#2} }
                  { \msg_error:nnV {osglecture-project}{unknown-setup-class}
                      \l_tmpb_tl }
              }
            {mode}
              { \exp_args:NNV \__osglecture_setup_append:Nnn
                  \g__osglecture_setup_mode_prop \l_tmpb_tl {#2} }
          }
      }
      { \msg_error:nnn {osglecture-project}{invalid-setup-key}{#1} }
  }

\cs_new_protected:Npn \__osglecture_setup_missing_value:n #1
  { \msg_error:nnn {osglecture-project}{invalid-setup-key}{#1} }

\NewDocumentCommand \OsgLectureSetup { m }
  {
    \keyval_parse:NNn
      \__osglecture_setup_missing_value:n
      \__osglecture_setup_store:nn
      {#1}
  }

\NewExpandableDocumentCommand \OsgLectureSetupClassOptions { m }
  { \prop_item:Nn \g__osglecture_setup_class_prop {#1} }

\NewDocumentCommand \DeclareOsgLectureModeSetupArea { m m }
  {
    \keys_define:nn { osglecture/setup/mode }
      { #1 .code:n = { #2 {##1} } }
  }
\cs_new_protected:Npn \osglecture_project_apply_mode_setup:
  {
    \exp_args:Ne \clist_map_inline:nn { \ActiveLectureModes }
      {
        \prop_get:NnNT \g__osglecture_setup_mode_prop {##1} \l_tmpa_tl
          { \exp_args:NnV \keys_set:nn { osglecture/setup/mode } \l_tmpa_tl }
      }
  }

% \ldeen*{Präambelfragmente werden registriert, während @1 gelesen wird,
% bevor die Basisklasse bekannt ist. @2 spielt die Warteschlange erst ab,
% nachdem die Klasse geladen, der Mode-Graph finalisiert und der Adapter
% installiert ist.}{Preamble fragments are registered while @1 is read,
% before the base class is known. @2 replays the queue only after loading
% the class, finalizing the mode graph, and installing the
% adapter.}{\file{projectconfig.tex}}{\file{osglecture.cls}}
\keys_define:nn { osglecture/preamble-include }
  {
    class .tl_set:N = \l__osglecture_preamble_class_tl,
  }

% \ldeen*{Die unsternierte Form bleibt der Normalfall: @1 registriert das
% Fragment nur und spielt es -- ggf. bedingt durch Modus/Klasse -- erst
% nach @2 ab (siehe @3 oben). Die Sternvariante ist der bewusste Ausstieg
% für Nutzer, die wissen, dass ihr Fragment vor @2 laufen muss (z.\,B.
% weil es selbst \cs{LectureTargetSetup} aufruft); sie lädt sofort, mit
% derselben Dateisuche wie die unsternierte Form, aber ohne Modus-/
% Klassenbedingung, da Modus-Graph und Basisklasse an dieser Stelle noch
% nicht feststehen.}{The unstarred form remains the normal case: @1 only
% registers the fragment and replays it -- possibly conditioned on
% mode/class -- only after @2 (see @3 above). The starred variant is the
% deliberate opt-out for users who know their fragment must run before
% @2 (e.g. because it itself calls \cs{LectureTargetSetup}); it loads
% immediately, using the same file lookup as the unstarred form, but
% without a mode/class condition, since the mode graph and base class
% are not yet known at that point.}{\cs{IncludeOsgLecturePreamble}}
% {\cs{LoadClass}}{\cs{osglecture\_project\_intercept\_package\_commands:}}
\NewDocumentCommand \IncludeOsgLecturePreamble { s d<> O{} m }
  {
    \IfBooleanTF{#1}
      {
        \IfNoValueF{#2}
          { \msg_error:nnn {osglecture-project}{preamble-star-condition}{#4} }
        \tl_if_empty:nF {#3}
          { \msg_error:nnn {osglecture-project}{preamble-star-condition}{#4} }
        \__osglecture_preamble_input:n {#4}
      }
      {
        \seq_gput_right:Nx \g__osglecture_project_preamble_seq
          {
            \exp_not:N \__osglecture_preamble_consider:nnn
              { \IfNoValueTF{#2}{}{\exp_not:n{#2}} }
              { \exp_not:n{#3} }
              { \exp_not:n{#4} }
          }
      }
  }

\cs_new_protected:Npn \__osglecture_preamble_consider:nnn #1#2#3
  {
    \bool_set_true:N \l__osglecture_preamble_include_bool
    \tl_clear:N \l__osglecture_preamble_class_tl
    \keys_set:nn { osglecture/preamble-include } {#2}
    \tl_if_blank:nF {#1}
      {
        \IfLectureModeTF {#1} { }
          { \bool_set_false:N \l__osglecture_preamble_include_bool }
      }
    \tl_if_empty:NF \l__osglecture_preamble_class_tl
      {
        \exp_args:NNV \clist_if_in:NnF \l__osglecture_preamble_class_tl
          \l__osglecture_preamble_current_class_tl
          { \bool_set_false:N \l__osglecture_preamble_include_bool }
      }
    \bool_if:NT \l__osglecture_preamble_include_bool
      { \__osglecture_preamble_input:n {#3} }
  }

\cs_new_protected:Npn \__osglecture_preamble_input:n #1
  {
    \tl_clear:N \l_tmpa_tl
    \cs_if_exist:NT \OsgLectureManifestAvailableFlag
      {
        \str_if_eq:eeT { \OsgLectureManifestAvailableFlag } { 1 }
          { \tl_set:Ne \l_tmpa_tl { \OsgLectureManifestSharedTexDirectory } }
      }
    \tl_if_empty:NF \l_tmpa_tl
      {
        \tl_put_right:Nn \l_tmpa_tl { /#1.tex }
        \file_if_exist:VTF \l_tmpa_tl
          {
            \iow_term:x
              { osglecture:~loading~preamble~fragment~'#1'~from~'\l_tmpa_tl' }
            \exp_args:NV \file_input:n \l_tmpa_tl
            \tl_set:Nn \l_tmpa_tl { found }
          }
          { \tl_clear:N \l_tmpa_tl }
      }
    \tl_if_empty:NT \l_tmpa_tl
      {
        \tl_set:Nn \l_tmpb_tl { osglecture-preamble-#1.tex }
        \file_if_exist:VTF \l_tmpb_tl
          {
            \iow_term:x
              { osglecture:~loading~installed~preamble~fragment~'#1' }
            \exp_args:NV \file_input:n \l_tmpb_tl
          }
          { \msg_error:nnn {osglecture-project}{preamble-not-found}{#1} }
      }
  }

\cs_new_protected:Npn \osglecture_project_include_preambles:n #1
  {
    \tl_set:Nn \l__osglecture_preamble_current_class_tl {#1}
    \seq_map_inline:Nn \g__osglecture_project_preamble_seq
      {##1}
  }
% \ldeen*{@1 steuert die Farbe, in der @2 dokumentübergreifende
% @3-Ziele einfärbt (hyperrefs \code{filecolor}, getrennt von
% \code{linkcolor}). Unkonfiguriert bleibt @2s eigener Default aktiv --
% @4 setzt \code{filecolor} nur, wenn diese Eigenschaft für das aktuelle
% Target tatsächlich gesetzt wurde.}{@1 controls the color @2 uses for
% cross-document @3 targets (hyperref's \code{filecolor}, distinct from
% \code{linkcolor}). Left unconfigured, @2's own default stays in
% effect -- @4 only sets \code{filecolor} when this property was
% actually set for the current target.}{\code{filecolor}}{hyperref}{\cs{olref}}{\code{osglecture-references.sty}}
\keys_define:nn { osglecture/project-target }
  {
    profile .code:n =
      { \prop_gput:NVn \g__osglecture_project_target_profile_prop
          \l__osglecture_project_target_tl {#1} },
    filecolor .code:n =
      { \prop_gput:NVn \g__osglecture_project_target_filecolor_prop
          \l__osglecture_project_target_tl {#1} },
  }
\NewDocumentCommand \LectureTargetSetup { m m }
  {
    \tl_set:Nn \l__osglecture_project_target_tl {#1}
    \keys_set:nn {osglecture/project-target} {#2}
  }
\NewExpandableDocumentCommand \OsgLectureTargetProfile { m }
  { \prop_item:Nn \g__osglecture_project_target_profile_prop {#1} }
\NewExpandableDocumentCommand \OsgLectureTargetFileColor { m }
  { \prop_item:Nn \g__osglecture_project_target_filecolor_prop {#1} }
\NewDocumentCommand \LectureProjectEnforce { m }
  {
    \int_gset:Nn \g__osglecture_configuration_priority_int {40}
    \bool_gset_true:N \g__osglecture_configuration_enforced_bool
    \keys_set:nn { osglecture/project } {#1}
    \int_gset:Nn \g__osglecture_configuration_priority_int {20}
    \bool_gset_false:N \g__osglecture_configuration_enforced_bool
  }
\NewDocumentCommand \EnforceLectureProjectConfiguration { m }
  {
    \int_gset:Nn \g__osglecture_configuration_priority_int {40}
    \bool_gset_true:N \g__osglecture_configuration_enforced_bool
    #1
    \int_gset:Nn \g__osglecture_configuration_priority_int {20}
    \bool_gset_false:N \g__osglecture_configuration_enforced_bool
  }
\NewExpandableDocumentCommand \OsgLectureProjectValue { m }
  { \prop_item:Nn \g__osglecture_project_policy_prop {#1} }
\NewDocumentCommand \OsgLectureProjectEnforcedTF { m m m }
  {
    \prop_get:NnNTF \g__osglecture_project_policy_enforced_prop {#1} \l_tmpa_tl
      { \str_if_eq:VnTF \l_tmpa_tl {true}{#2}{#3} }
      {#3}
  }

\NewExpandableDocumentCommand \OsgLectureProjectClassOptions { }
  { \tl_use:N \g__osglecture_project_class_options_tl }

\cs_new_protected:Npn \__osglecture_metadata_dispatch:nnnn #1#2#3#4
  {
    \__osglecture_metadata_receiver:nnnnn {#1} {#2} {#3} {#4}
      { \int_use:N \g__osglecture_configuration_priority_int }
  }

\NewDocumentCommand \DeclareOsgLectureMetadataField { m }
  {
    \clist_if_in:NnF \g__osglecture_metadata_fields_clist {#1}
      {
        \clist_gput_right:Nn \g__osglecture_metadata_fields_clist {#1}
        \clist_if_in:nnT {title,author,date} {#1}
          {
            \cs_if_exist:cT {#1}
              { \cs_new_eq:cc { __osglecture_project_preexisting_#1 } {#1} }
          }
      }
    \cs_if_exist:cF { __osglecture_project_proxy_#1 }
      {
        \exp_args:Nc \NewDocumentCommand
          { __osglecture_project_proxy_#1 } { d<> O{} m }
          {
            \IfNoValueTF{##1}
              { \__osglecture_metadata_dispatch:nnnn {#1} {} {##2} {##3} }
              { \__osglecture_metadata_dispatch:nnnn {#1} {##1} {##2} {##3} }
          }
      }
    \cs_set_eq:cc {#1} { __osglecture_project_proxy_#1 }
  }

\cs_new_protected:Npn \osglecture_project_suspend_metadata_commands:
  {
    \clist_map_inline:Nn \g__osglecture_metadata_fields_clist
      {
        \cs_if_eq:ccT {##1} { __osglecture_project_proxy_##1 }
          {
            \cs_if_exist:cTF { __osglecture_project_preexisting_##1 }
              { \cs_set_eq:cc {##1} { __osglecture_project_preexisting_##1 } }
              { \cs_undefine:c {##1} }
          }
      }
  }

\cs_new_protected:Npn \osglecture_configuration_priority_set:n #1
  {
    \int_gset:Nn \g__osglecture_configuration_priority_int {#1}
    \bool_gset_false:N \g__osglecture_configuration_enforced_bool
  }

\int_gset:Nn \g__osglecture_configuration_priority_int {0}
\keys_set:nn {osglecture/project}
  {
    presentation-profile=beamer,
    longform-profile=book,
    identity-profile=osg,
    theme=osg,
    numbering=chapter,
    references=external-on-demand
  }
\int_gset:Nn \g__osglecture_configuration_priority_int {20}

\clist_map_inline:nn { title, subtitle, author, institute, date, course, event }
  { \DeclareOsgLectureMetadataField{#1} }

\ExplSyntaxOff
%</project>
%<*metadata>

\RequirePackage{osglecture-project}
\RequirePackage{osglecture-modes}

\ExplSyntaxOn

\tl_new:N \g__osglecture_metadata_course_target_tl
\tl_new:N \g__osglecture_metadata_event_target_tl
\prop_new:N \g__osglecture_metadata_value_prop
\prop_new:N \g__osglecture_metadata_short_prop
\prop_new:N \g__osglecture_metadata_priority_prop
\prop_new:N \g__osglecture_metadata_enforced_prop

\keys_define:nn { osglecture/metadata }
  {
    course-target .choices:nn = { title, subtitle, none }
      { \tl_gset_eq:NN \g__osglecture_metadata_course_target_tl \l_keys_choice_tl },
    course-target .initial:n = none,
    event-target .choices:nn = { title, subtitle, none }
      { \tl_gset_eq:NN \g__osglecture_metadata_event_target_tl \l_keys_choice_tl },
    event-target .initial:n = none,
  }

\NewDocumentCommand \OsgLectureConfigureMetadata { m }
  { \keys_set:nn { osglecture/metadata } {#1} }

% \ldeen*{Sichert die von der Basisklasse bereitgestellten
% Metadatenbefehle (z.\,B. \cs{title}, \cs{author}) unter eigenem Namen,
% bevor @1 sie unten durch die einheitlichen Proxys ersetzt: Nur so kann
% @2 später den nativen Befehl der Basisklasse tatsächlich aufrufen,
% statt sich selbst rekursiv zu erreichen.}{Saves the metadata commands
% provided by the base class (e.g. \cs{title}, \cs{author}) under a
% private name before @1 replaces them below with the unified proxies:
% only this lets @2 later actually call the base class's native command,
% instead of recursing into itself.}
% {\cs{DeclareOsgLectureMetadataField}}{\cs{\_\_osglecture\_metadata\_call\_native:nnn}}
\clist_map_inline:Nn \g__osglecture_metadata_fields_clist
  {
    \cs_if_exist:cT {#1}
      {
        \cs_if_eq:ccF {#1} { __osglecture_project_proxy_#1 }
          { \cs_new_eq:cc { __osglecture_metadata_native_#1 } {#1} }
      }
  }
% \ldeen{Manche Zielbasisklassen kennen kein natives \cs{institute}
% (beamerspezifisch), wohl aber \cs{publishers} für dieselbe
% Titelseiten-Zeile; siehe die Rückfalloption in
% \cs{\_\_osglecture\_metadata\_apply\_institute:nn}
% unten.}{Some target base classes have no native \cs{institute}
% (a beamer-specific command), but do have \cs{publishers} for the same
% title-page line; see the fallback in
% \cs{\_\_osglecture\_metadata\_apply\_institute:nn}
% below.}
\cs_if_exist:NT \publishers
  { \cs_new_eq:NN \__osglecture_metadata_native_publishers \publishers }

\cs_new_protected:Npn \__osglecture_metadata_call_native:nnn #1#2#3
  {
    \cs_if_exist:cT { __osglecture_metadata_native_#1 }
      {
        \str_case:enF { \OsgLectureBackend }
          {
            {beamer}   { \use:c { __osglecture_metadata_native_#1 } [#2] {#3} }
            {ltx-talk} { \use:c { __osglecture_metadata_native_#1 } [#2] {#3} }
          }
          { \use:c { __osglecture_metadata_native_#1 } {#3} }
      }
  }

\cs_new_protected:Npn \__osglecture_metadata_apply_institute:nn #1#2
  {
    \cs_if_exist:cTF { __osglecture_metadata_native_institute }
      { \__osglecture_metadata_call_native:nnn {institute} {#1} {#2} }
      {
        \cs_if_exist:NT \__osglecture_metadata_native_publishers
          { \__osglecture_metadata_native_publishers {#2} }
      }
  }

\cs_new_protected:Npn \__osglecture_metadata_apply_target:nnn #1#2#3
  {
    \str_case:nn {#1}
      {
        {title}    { \__osglecture_metadata_call_native:nnn {title} {#2} {#3} }
        {subtitle} { \__osglecture_metadata_call_native:nnn {subtitle} {#2} {#3} }
      }
  }
\cs_generate_variant:Nn \__osglecture_metadata_apply_target:nnn { Vnn }

\cs_new_protected:Npn \__osglecture_metadata_apply:nnn #1#2#3
  {
    \str_case:nnF {#1}
      {
        {course}
          { \__osglecture_metadata_apply_target:Vnn
              \g__osglecture_metadata_course_target_tl {#2} {#3} }
        {event}
          { \__osglecture_metadata_apply_target:Vnn
              \g__osglecture_metadata_event_target_tl {#2} {#3} }
        {institute} { \__osglecture_metadata_apply_institute:nn {#2} {#3} }
      }
      { \__osglecture_metadata_call_native:nnn {#1} {#2} {#3} }
  }

\cs_new_protected:Npn \__osglecture_metadata_set:nnnn #1#2#3#4
  {
    \prop_get:NnNTF \g__osglecture_metadata_priority_prop {#1} \l_tmpa_tl
      {
        \int_compare:nNnF {#4} < {\l_tmpa_tl}
          { \__osglecture_metadata_store:nnnn {#1}{#2}{#3}{#4} }
        \int_compare:nNnT {#4} < {\l_tmpa_tl}
          {
            \prop_get:NnN \g__osglecture_metadata_enforced_prop {#1} \l_tmpb_tl
            \prop_get:NnN \g__osglecture_metadata_value_prop {#1} \l_tmpc_tl
            \str_if_eq:VnT \l_tmpb_tl {true}
              { \tl_if_eq:VnF \l_tmpc_tl {#3}
                  { \msg_fatal:nnxxx {osglecture-project}{enforced-conflict}
                      {#1}{\l_tmpc_tl}{#3} } }
          }
      }
      { \__osglecture_metadata_store:nnnn {#1}{#2}{#3}{#4} }
  }

\cs_new_protected:Npn \__osglecture_metadata_store:nnnn #1#2#3#4
  {
    \tl_if_blank:nTF {#2}
      { \prop_gput:Nnn \g__osglecture_metadata_short_prop {#1} {#3} }
      { \prop_gput:Nnn \g__osglecture_metadata_short_prop {#1} {#2} }
    \prop_gput:Nnn \g__osglecture_metadata_value_prop {#1} {#3}
    \prop_gput:Nnn \g__osglecture_metadata_priority_prop {#1} {#4}
    \prop_gput:Nnx \g__osglecture_metadata_enforced_prop {#1}
      {\bool_if:NTF \g__osglecture_configuration_enforced_bool {true}{false}}
    \tl_if_blank:nTF {#2}
      { \__osglecture_metadata_apply:nnn {#1} {#3} {#3} }
      { \__osglecture_metadata_apply:nnn {#1} {#2} {#3} }
  }

\cs_new_protected:Npn \__osglecture_metadata_receive:nnnnn #1#2#3#4#5
  {
    \tl_if_blank:nTF {#2}
      { \__osglecture_metadata_set:nnnn {#1} {#3} {#4} {#5} }
      {
        \IfLectureModeTF{#2}
          { \__osglecture_metadata_set:nnnn {#1} {#3} {#4} {#5} }
          { }
      }
  }

\cs_set_eq:NN \__osglecture_metadata_receiver:nnnnn
  \__osglecture_metadata_receive:nnnnn

\clist_map_inline:Nn \g__osglecture_metadata_fields_clist
  { \cs_set_eq:cc {#1} { __osglecture_project_proxy_#1 } }

\cs_new_eq:NN \__osglecture_metadata_record:nnnnn
  \__osglecture_metadata_receive:nnnnn
\seq_use:Nn \g__osglecture_metadata_pending_seq { }
\seq_gclear:N \g__osglecture_metadata_pending_seq

\NewExpandableDocumentCommand \OsgLectureMetadataValue { m }
  { \prop_item:Nn \g__osglecture_metadata_value_prop {#1} }
\NewExpandableDocumentCommand \OsgLectureMetadataShortValue { m }
  { \prop_item:Nn \g__osglecture_metadata_short_prop {#1} }

\cs_new_protected:Npn \__osglecture_metadata_compat:nnnn #1#2#3#4
  {
    \IfNoValueTF{#2}
      { \__osglecture_metadata_dispatch:nnnn {#1} {} {#3} {#4} }
      { \__osglecture_metadata_dispatch:nnnn {#1} {#2} {#3} {#4} }
  }
\NewDocumentCommand \OsgLectureSetTitle { d<> O{} m }
  { \__osglecture_metadata_compat:nnnn {title}{#1}{#2}{#3} }
\NewDocumentCommand \OsgLectureSetSubtitle { d<> O{} m }
  { \__osglecture_metadata_compat:nnnn {subtitle}{#1}{#2}{#3} }
\NewDocumentCommand \OsgLectureSetAuthor { d<> O{} m }
  { \__osglecture_metadata_compat:nnnn {author}{#1}{#2}{#3} }
\NewDocumentCommand \OsgLectureSetInstitute { d<> O{} m }
  { \__osglecture_metadata_compat:nnnn {institute}{#1}{#2}{#3} }
\NewDocumentCommand \OsgLectureSetDate { d<> O{} m }
  { \__osglecture_metadata_compat:nnnn {date}{#1}{#2}{#3} }
\NewDocumentCommand \OsgLectureSetCourse { d<> O{} m }
  { \__osglecture_metadata_compat:nnnn {course}{#1}{#2}{#3} }
\NewDocumentCommand \OsgLectureSetEvent { d<> O{} m }
  { \__osglecture_metadata_compat:nnnn {event}{#1}{#2}{#3} }

\NewExpandableDocumentCommand \OsgLectureTitle {} { \OsgLectureMetadataValue{title} }
\NewExpandableDocumentCommand \OsgLectureShortTitle {} { \OsgLectureMetadataShortValue{title} }
\NewExpandableDocumentCommand \OsgLectureSubtitle {} { \OsgLectureMetadataValue{subtitle} }
\NewExpandableDocumentCommand \OsgLectureShortSubtitle {} { \OsgLectureMetadataShortValue{subtitle} }
\NewExpandableDocumentCommand \OsgLectureAuthor {} { \OsgLectureMetadataValue{author} }
\NewExpandableDocumentCommand \OsgLectureShortAuthor {} { \OsgLectureMetadataShortValue{author} }
\NewExpandableDocumentCommand \OsgLectureInstitute {} { \OsgLectureMetadataValue{institute} }
\NewExpandableDocumentCommand \OsgLectureShortInstitute {} { \OsgLectureMetadataShortValue{institute} }
\NewExpandableDocumentCommand \OsgLectureDate {} { \OsgLectureMetadataValue{date} }
\NewExpandableDocumentCommand \OsgLectureShortDate {} { \OsgLectureMetadataShortValue{date} }
\NewExpandableDocumentCommand \OsgLectureCourse {} { \OsgLectureMetadataValue{course} }
\NewExpandableDocumentCommand \OsgLectureShortCourse {} { \OsgLectureMetadataShortValue{course} }
\NewExpandableDocumentCommand \OsgLectureEvent {} { \OsgLectureMetadataValue{event} }
\NewExpandableDocumentCommand \OsgLectureShortEvent {} { \OsgLectureMetadataShortValue{event} }

% \ldeen{Legacy-Aliasnamen aus der Vor-osglecture-Ära für
% \cs{course}/\cs{event}, weiterhin unterstützt für bestehende
% Dokumentquellen.}{Legacy alias names from the pre-osglecture era for
% \cs{course}/\cs{event}, still supported for existing document
% sources.}
\cs_set_eq:NN \lehrveranstaltung \course
\cs_set_eq:NN \conference \event

\ExplSyntaxOff
%</metadata>
%<*references>

\RequirePackage{expl3}
\ExplSyntaxOn
\tl_new:N \l__olref_initial_language_tl
\cs_if_exist:NTF \OsgLectureLanguage
  { \tl_set:Ne \l__olref_initial_language_tl { \OsgLectureLanguage } }
  { \tl_set:Nn \l__olref_initial_language_tl { en } }
\cs_if_exist:NF \extrasenglish { \cs_new:Npn \extrasenglish { } }
\cs_if_exist:NF \extrasngerman { \cs_new:Npn \extrasngerman { } }
\PassOptionsToPackage{nospace,ngerman}{varioref}
\ExplSyntaxOff
\RequirePackage{hyperref}
% \ldeen*{Unter \cs{DocumentMetadata}s Bootstrap bleibt
% \cs{@pdfborder} leer statt, wie unter \code{colorlinks} sonst üblich,
% auf "0 0 0" gesetzt zu werden -- die Zielfarben (\cs{@linkcolor},
% \cs{@filecolor}) werden dabei korrekt gesetzt, nur die Rahmenbreite
% nicht. Ein leeres \code{/Border[]} ist kein gültiges PDF-Array;
% Betrachter fallen dann auf einen sichtbaren Standardrahmen zurück --
% sichtbar nur bei über \cs{xr} importierten (dokumentübergreifenden)
% Verweisen wie \cs{olref}, da nur diese den Datei-Verweis-Codepfad
% (\code{GoToR}) durchlaufen. @1 erzwingt explizit denselben
% randlosen Zustand, den \code{colorlinks} für alle anderen Verweisarten
% bereits liefert.}{Under \cs{DocumentMetadata}'s bootstrap,
% \cs{@pdfborder} stays empty instead of being set to "0 0 0" as usual
% under \code{colorlinks} -- the target colors (\cs{@linkcolor},
% \cs{@filecolor}) get set correctly, just not the border width. An empty
% \code{/Border[]} is not a valid PDF array; viewers then fall back to a
% visible default border -- visible only for cross-document references
% imported via \cs{xr}, such as \cs{olref}, since only those go through
% the file-link (\code{GoToR}) code path. @1 explicitly forces the same
% borderless state that \code{colorlinks} already provides for every
% other link type.}{\cs{hypersetup}\code{\{pdfborder=0~0~0\}}}
\hypersetup{pdfborder=0 0 0}
\RequirePackage{varioref}
\RequirePackage{xr}

\ExplSyntaxOn

\str_if_eq:VnTF \l__olref_initial_language_tl {de}
  { \cs_if_exist:NT \extrasngerman { \extrasngerman } }
  { \cs_if_exist:NT \extrasenglish { \extrasenglish } }

% \ldeen*{Die Farbe, in der \cs{olref} dokumentübergreifende Ziele
% einfärbt, ist projektseitig je Target konfigurierbar
% (\cs{LectureTargetSetup}\code{\{<target>\}\{filecolor=<Farbe>\}} in
% \file{projectconfig.tex}), damit unterschiedliche Doctypes
% unterschiedliche Farbschemata verwenden können. Ohne Konfiguration
% bleibt hyperrefs eigener \code{filecolor}-Default (cyan) unangetastet.
% @1 setzt \cs{@filecolor} bewusst direkt statt über
% \cs{hypersetup}\code{\{filecolor=...\}}: Unter \cs{DocumentMetadata}s
% Bootstrap hat sich empirisch gezeigt, dass \code{hypersetup}s
% \code{filecolor}-Schlüssel wirkungslos bleibt (unabhängig vom
% Aufrufzeitpunkt -- Präambel, \code{begindocument}-Hooks oder
% Dokumentanfang wurden alle getestet), während eine direkte Zuweisung
% zuverlässig funktioniert. Das ist keine Abkürzung auf Kosten der
% Farbmodell-Auflösung: Für einen einfachen Farbnamen ohne führende
% eckige Klammer geht \cs{HyColor@UseColor} (in \code{hycolor.sty}, der
% Konsument von \cs{@filecolor}) ohnehin über denselben
% \cs{color}\code{\{...\}}-Pfad, den \code{hypersetup} für diesen Fall
% erzeugt hätte.}{The color \cs{olref} uses for cross-document targets
% is configurable per target from the project side
% (\cs{LectureTargetSetup}\code{\{<target>\}\{filecolor=<color>\}} in
% \file{projectconfig.tex}), so different doctypes can use different
% color schemes. Left unconfigured, hyperref's own \code{filecolor}
% default (cyan) stays untouched. @1 deliberately sets \cs{@filecolor}
% directly instead of via \cs{hypersetup}\code{\{filecolor=...\}}: under
% \cs{DocumentMetadata}'s bootstrap, \code{hypersetup}'s
% \code{filecolor} key was found empirically to have no effect
% (regardless of call timing -- preamble, \code{begindocument} hooks,
% and document start were all tested), while a direct assignment works
% reliably. This is not a shortcut at the expense of color-model
% resolution: for a plain color name without a leading square bracket,
% \cs{HyColor@UseColor} (in \code{hycolor.sty}, the consumer of
% \cs{@filecolor}) goes through the same \cs{color}\code{\{...\}} path that
% \code{hypersetup} would have produced for this case anyway.}
% {\cs{\_\_olref\_apply\_filecolor:n}}
\cs_new_protected:Npn \__olref_apply_filecolor:n #1
  { \cs_set:Npn \@filecolor {#1} }
\OsgLectureBuildLoadedTF
  {
    \tl_set:Ne \l_tmpa_tl
      { \exp_args:Ne \OsgLectureTargetFileColor { \OsgLectureBuildValue{target} } }
    \tl_if_empty:NF \l_tmpa_tl
      { \exp_args:NV \__olref_apply_filecolor:n \l_tmpa_tl }
  }
  { }

\prop_new:N \g__olref_documents_prop
\prop_new:N \g__olref_imported_prop
\prop_new:N \g__olref_used_prop
\clist_new:N \g__olref_replace_clist
\bool_new:N \g__olref_legacy_bool
\bool_new:N \g__olref_replacements_installed_bool
\iow_new:N \g__olref_export_iow
\iow_new:N \g__olref_used_iow
\cs_generate_variant:Nn \prop_if_in:NnTF { NV }
\cs_generate_variant:Nn \prop_get:NnN { NV }

\tl_new:N \l__olref_unit_tl
\tl_new:N \l__olref_type_tl
\tl_new:N \l__olref_lang_tl
\tl_new:N \l__olref_aux_tl
\tl_new:N \l__olref_pdf_tl
\tl_new:N \l__olref_generation_tl
\tl_new:N \l__olref_identity_tl
\tl_new:N \l__olref_prefix_tl
\tl_new:N \l__olref_label_tl

\msg_new:nnn { osglecture-references } { invalid-unit }
  {
    Invalid~logical~unit~identifier~'#1'.~
    Unit~identifiers~must~not~contain~whitespace,~comma,~equals,~or~'...'.
  }
\msg_new:nnn { osglecture-references } { unknown-option }
  { Unknown~reference-address~key~'#1'. }
\msg_new:nnn { osglecture-references } { unknown-unit }
  {
    Logical~unit~'#1'~is~not~known~in~this~series.~Build~that~unit~once~
    before~resolving~references~to~it;~using~an~undefined~reference.
  }
\msg_new:nnn { osglecture-references } { missing-projection }
  {
    Logical~unit~'#1'~is~known~as~physical~unit~'#2',~but~no~reference~
    export~is~available~for~the~requested~document~type~and~language.~
    Build~that~projection;~using~an~undefined~reference.
  }
\msg_new:nnn { osglecture-references } { missing-registration-key }
  { Reference-document~registration~does~not~define~'#1'. }
\msg_new:nnn { osglecture-references } { replace-standard }
  {
    Replacing~standard~reference~commands~with~the~osglecture~
    optional-address~syntax:~#1.
  }

\cs_new_eq:NN \__olref_original_ref \ref
\cs_new_eq:NN \__olref_original_pageref \pageref
\cs_new_eq:NN \__olref_original_nameref \nameref
\cs_new_eq:NN \__olref_original_autoref \autoref

\cs_new_protected:Npn \__olref_replace_set:n #1
  {
    \clist_gclear:N \g__olref_replace_clist
    \clist_map_inline:nn {#1}
      {
        \str_case:nnF {##1}
          {
            {none} { \clist_gclear:N \g__olref_replace_clist }
            {all}
              {
                \clist_gset:Nn \g__olref_replace_clist
                  {ref,pageref,nameref,autoref}
              }
            {ref}      { \clist_gput_right:Nn \g__olref_replace_clist {ref} }
            {pageref}  { \clist_gput_right:Nn \g__olref_replace_clist {pageref} }
            {nameref}  { \clist_gput_right:Nn \g__olref_replace_clist {nameref} }
            {autoref}  { \clist_gput_right:Nn \g__olref_replace_clist {autoref} }
          }
          { \msg_error:nnn {osglecture-references}{unknown-option}{##1} }
      }
    \clist_gremove_duplicates:N \g__olref_replace_clist
  }

\keys_define:nn { osglecture/references }
  {
    legacy .bool_gset:N = \g__olref_legacy_bool,
    legacy .default:n = true,
    replace .code:n = \__olref_replace_set:n {#1},
    replace .default:n = all,
  }

\ProcessKeyOptions [ osglecture/references ]

\NewDocumentCommand \OsgLectureReferencesSetup { m }
  { \keys_set:nn {osglecture/references} {#1} }

\cs_new:Npn \__olref_current_type:
  {
    \cs_if_exist:NTF \OsgLectureDocumentType
      { \OsgLectureDocumentType }
      { }
  }
\cs_new:Npn \__olref_current_lang:
  {
    \cs_if_exist:NTF \OsgLectureLanguage
      { \OsgLectureLanguage }
      { }
  }

\keys_define:nn { osglecture/reference-address }
  {
    unit .tl_set:N = \l__olref_unit_tl,
    type .tl_set:N = \l__olref_type_tl,
    lang .tl_set:N = \l__olref_lang_tl,
    unknown .code:n =
      { \msg_error:nnn {osglecture-references}{unknown-option}{\l_keys_key_str} },
  }

\cs_new_protected:Npn \__olref_address_parse:n #1
  {
    \tl_clear:N \l__olref_unit_tl
    \tl_set:Ne \l__olref_type_tl { \__olref_current_type: }
    \tl_set:Ne \l__olref_lang_tl { \__olref_current_lang: }
    \clist_map_inline:nn {#1}
      {
        \tl_if_in:nnTF {##1} {=}
          { \keys_set:nn {osglecture/reference-address} {##1} }
          { \tl_set:Nn \l__olref_unit_tl {##1} }
      }
    \tl_trim_spaces:N \l__olref_unit_tl
    \tl_trim_spaces:N \l__olref_type_tl
    \tl_trim_spaces:N \l__olref_lang_tl
    \tl_if_empty:NF \l__olref_unit_tl
      {
        \regex_match:nVTF
          { \A [^\s,=]+ \Z }
          \l__olref_unit_tl
          {
            \tl_if_in:VnT \l__olref_unit_tl {...}
              {
                \msg_error:nnV {osglecture-references}{invalid-unit}
                  \l__olref_unit_tl
              }
          }
          {
            \msg_error:nnV {osglecture-references}{invalid-unit}
              \l__olref_unit_tl
          }
      }
  }

\keys_define:nn { osglecture/reference-document }
  {
    unit .tl_set:N = \l__olref_unit_tl,
    type .tl_set:N = \l__olref_type_tl,
    lang .tl_set:N = \l__olref_lang_tl,
    aux .tl_set:N = \l__olref_aux_tl,
    pdf .tl_set:N = \l__olref_pdf_tl,
    generation .tl_set:N = \l__olref_generation_tl,
  }

\cs_new:Npn \__olref_identity:nnn #1#2#3
  { unit=#1,type=#2,lang=#3 }

\NewDocumentCommand \OsgLectureReferenceDocument { m }
  {
    \tl_clear:N \l__olref_unit_tl
    \tl_clear:N \l__olref_type_tl
    \tl_clear:N \l__olref_lang_tl
    \tl_clear:N \l__olref_aux_tl
    \tl_clear:N \l__olref_pdf_tl
    \tl_clear:N \l__olref_generation_tl
    \keys_set:nn {osglecture/reference-document} {#1}
    \clist_map_inline:nn {unit,type,lang,aux,pdf,generation}
      {
        \tl_if_empty:cT {l__olref_##1_tl}
          {
            \msg_error:nnn {osglecture-references}{missing-registration-key}
              {##1}
          }
      }
    \tl_set:Ne \l__olref_identity_tl
      {
        \__olref_identity:nnn
          {\l__olref_unit_tl}{\l__olref_type_tl}{\l__olref_lang_tl}
      }
    \prop_gput:NVx \g__olref_documents_prop \l__olref_identity_tl
      {
        {\exp_not:V \l__olref_aux_tl}
        {\exp_not:V \l__olref_pdf_tl}
        {\exp_not:V \l__olref_generation_tl}
      }
    \tl_set:Ne \l__olref_prefix_tl
      { olr-\str_mdfive_hash:e {\l__olref_identity_tl}: }
    \__olref_import_now:
  }

\NewDocumentCommand \OsgLectureReferenceDocumentTF { m m m m m }
  {
    \tl_set:Ne \l__olref_identity_tl
      { \__olref_identity:nnn {#1}{#2}{#3} }
    \prop_if_in:NVTF \g__olref_documents_prop \l__olref_identity_tl
      {#4} {#5}
  }
\cs_new:Npn \__olref_document_aux:nnn #1#2#3 {#1}
\cs_new:Npn \__olref_document_pdf:nnn #1#2#3 {#2}
\cs_new:Npn \__olref_document_generation:nnn #1#2#3 {#3}
\NewDocumentCommand \OsgLectureReferenceDocumentAuxTo { m m m m }
  {
    \tl_set:Ne \l__olref_identity_tl
      { \__olref_identity:nnn {#1}{#2}{#3} }
    \prop_get:NVN \g__olref_documents_prop \l__olref_identity_tl #4
    \tl_set:Ne #4
      { \exp_after:wN \__olref_document_aux:nnn #4 }
  }
\NewDocumentCommand \OsgLectureReferenceDocumentPdfTo { m m m m }
  {
    \tl_set:Ne \l__olref_identity_tl
      { \__olref_identity:nnn {#1}{#2}{#3} }
    \prop_get:NVN \g__olref_documents_prop \l__olref_identity_tl #4
    \tl_set:Ne #4
      { \exp_after:wN \__olref_document_pdf:nnn #4 }
  }
\NewDocumentCommand \OsgLectureReferenceDocumentGenerationTo { m m m m }
  {
    \tl_set:Ne \l__olref_identity_tl
      { \__olref_identity:nnn {#1}{#2}{#3} }
    \prop_get:NVN \g__olref_documents_prop \l__olref_identity_tl #4
    \tl_set:Ne #4
      { \exp_after:wN \__olref_document_generation:nnn #4 }
  }

\NewDocumentCommand \OsgLectureWriteReferenceExport { m }
  { \iow_now:Ne \g__olref_export_iow { \exp_not:n {#1} } }
\NewDocumentCommand \OsgLectureWriteReferenceUse { m }
  { \iow_now:Ne \g__olref_used_iow { \exp_not:n {#1} } }

\ProvideDocumentCommand \OsgLectureReferenceExport
  { m m m m m m m } { }
\ProvideDocumentCommand \OsgLectureReferenceUse
  { m m m m m m m } { }
\ProvideDocumentCommand \OsgLectureContinuationCounter { m m } { }
\ProvideDocumentCommand \OsgLectureIntegrationUse { m m m m m } { }

\cs_new_protected:Npn \__olref_use_write:nn #1#2
  {
    \tl_set:Ne \l_tmpa_tl
      {
        \l__olref_identity_tl,label=#1,property=#2
      }
    \prop_if_in:NVF \g__olref_used_prop \l_tmpa_tl
      {
        \OsgLectureBuildLoadedTF
          {
            \tl_if_empty:eF { \OsgLectureBuildValue{generation-id} }
              {
                \iow_now:Ne \g__olref_used_iow
                  {
                    \token_to_str:N \OsgLectureReferenceUse
                    {\OsgLectureBuildValue{generation-id}}
                    {\exp_not:V \l__olref_generation_tl}
                    {\exp_not:V \l__olref_unit_tl}
                    {\exp_not:V \l__olref_type_tl}
                    {\exp_not:V \l__olref_lang_tl}
                    {\exp_not:n {#1}}
                    {\exp_not:n {#2}}
                  }
              }
          }
          { }
        \prop_gput:NVn \g__olref_used_prop \l_tmpa_tl {true}
      }
  }

% \ldeen*{@1 (\cs{externaldocument}) ist nach Konvention ein
% Präambelbefehl. Ein mitten im Absatz ausgeführtes @1 hinterlässt dort
% unsichtbares, aber dehnbares Material, das der Zeilenumbruch-Algorithmus
% beim Justieren ungleichmäßig auf die Zeile verteilt -- sichtbar als ein
% überproportional großer Zwischenraum vor der eigentlichen Referenz. Ein
% erneuter @3-Aufruf auf eine bereits importierte Einheit zeigt dieses
% Symptom nicht, da @1 dann nicht nochmal aufgerufen wird -- beim Testen
% also leicht zu übersehen. @4 kapselt den Importcode deshalb in einer
% eigenen, wiederverwendbaren Funktion: @2 ruft sie lazy als Fallback auf
% (für noch nicht importierte Einheiten), \cs{OsgLectureReferenceDocument}
% ruft sie zusätzlich eager auf -- Letzteres läuft im
% \code{begindocument/before}-Hook und damit sicher im vertikalen Modus,
% bevor der erste Absatz beginnt.}{@1 (\cs{externaldocument}) is a
% preamble-only command by convention. An @1 executed mid-paragraph
% leaves invisible but stretchable material behind that the
% line-breaking algorithm distributes unevenly across the line when
% justifying -- visible as a disproportionately large gap right before
% the actual reference. A repeated @3 call for an already-imported unit
% does not show this symptom, since @1 is not called again then -- easy
% to miss while testing. @4 therefore factors the import code into its
% own, reusable function: @2 calls it lazily as a fallback (for units
% not yet imported), and \cs{OsgLectureReferenceDocument} additionally
% calls it eagerly -- the latter runs in the \code{begindocument/before}
% hook and is therefore safely in vertical mode, before the first
% paragraph begins.}{\cs{externaldocument}}{\cs{\_\_olref\_import:nn}}{\cs{olref}}{\cs{\_\_olref\_import\_now:}}
\cs_new_protected:Npn \__olref_import_now:
  {
    \tl_clear:N \l__olref_generation_tl
    \prop_if_in:NVF \g__olref_imported_prop \l__olref_identity_tl
      {
        \prop_get:NVNTF
          \g__olref_documents_prop \l__olref_identity_tl \l_tmpa_tl
          {
            \exp_after:wN \__olref_import_aux:nnn \l_tmpa_tl
          }
          {
            \cs_if_exist:NTF \OsgLectureSeriesUnitKnownTF
              {
                \exp_args:NV \OsgLectureSeriesUnitKnownTF \l__olref_unit_tl
                  {
                    \msg_warning:nnxx
                      {osglecture-references}{missing-projection}
                      {\l__olref_unit_tl}
                      {\OsgLectureSeriesPhysicalUnit{\l__olref_unit_tl}}
                  }
                  {
                    \msg_warning:nnV {osglecture-references}{unknown-unit}
                      \l__olref_unit_tl
                  }
              }
              {
                \msg_warning:nnV {osglecture-references}{unknown-unit}
                  \l__olref_unit_tl
              }
          }
        \prop_gput:NVn
          \g__olref_imported_prop \l__olref_identity_tl {true}
      }
  }
\cs_new_protected:Npn \__olref_import:nn #1#2
  {
    \tl_set:Ne \l__olref_identity_tl
      {
        \__olref_identity:nnn
          {\l__olref_unit_tl}{\l__olref_type_tl}{\l__olref_lang_tl}
      }
    \tl_set:Ne \l__olref_prefix_tl
      { olr-\str_mdfive_hash:e {\l__olref_identity_tl}: }
    \__olref_import_now:
    \prop_get:NVNT \g__olref_documents_prop \l__olref_identity_tl \l_tmpa_tl
      { \exp_after:wN \__olref_generation_set:nnn \l_tmpa_tl }
    \__olref_use_write:nn {#1}{#2}
  }

\cs_new_protected:Npn \__olref_generation_set:nnn #1#2#3
  { \tl_set:Nn \l__olref_generation_tl {#3} }

\cs_new_protected:Npn \__olref_import_aux:nnn #1#2#3
  {
    \tl_set:Nn \l__olref_generation_tl {#3}
    \tl_set:Nn \l_tmpa_tl {#1}
    \regex_replace_once:nnN { [.]aux \Z } { } \l_tmpa_tl
    \externaldocument[\l__olref_prefix_tl]{\l_tmpa_tl}[#2]
  }

\cs_new_protected:Npn \__olref_target_set:nnn #1#2#3
  {
    \__olref_address_parse:n {#1}
    \tl_set:Nn \l__olref_label_tl {#2}
    \tl_if_empty:NF \l__olref_unit_tl
      {
        \__olref_import:nn {#2}{#3}
        \tl_put_left:NV \l__olref_label_tl \l__olref_prefix_tl
      }
  }

% \ldeen*{Ein dokumentübergreifender @1-Verweis erzeugt normalerweise
% über hyperrefs eigenes \cs{hyper@linkfile} einen reinen
% \code{GoToR}-Link (Sprung in eine externe Datei) -- unabhängig davon,
% ob das Ziel inzwischen (via \cs{includeunit}) real Teil desselben
% integrierten Dokuments geworden ist. tagpax übernimmt beim
% PDF-Zusammenführen jede Annotation unverändert und hat keine
% Möglichkeit zu erkennen, dass ein \code{GoToR} nun eigentlich ein
% interner Sprung sein müsste. @2 baut den Link deshalb nicht mehr über
% hyperrefs Standardpfad (\cs{hyper@linkfile}), sondern bettet
% zusätzlich die logische Zielidentität (Unit/Doctyp/Sprache) sowie die
% Quellseite als eigene Schlüssel in die Annotation ein -- dieselbe
% Identität, die \cs{includeunit} bereits kennt. So kann tagpax (bzw.
% dessen Erweiterung) am Link selbst erkennen, ob das Ziel intern
% verfügbar ist, ganz ohne fragile Dateipfad-Vergleiche. Zwei
% Implementierungen sind nötig, je nachdem, ob \cs{DocumentMetadata}
% für das aktuelle Target aktiv ist (in OSGLecture targetabhängig, siehe
% \code{document\_metadata} in \file{ollmconfig.toml}): @3 dupliziert
% hyperrefs klassische \cs{hyper@linkfile}-Struktur direkt
% (\cs{pdfstartlink}/\cs{close@pdflink}, LuaTeX-spezifisch --
% OSGLecture ist ohnehin LuaTeX-only), da diese Primitiv-Wrapper unter
% \cs{DocumentMetadata} \emph{nicht} definiert sind (die
% PDF-Management-Schicht übernimmt dort die Link-Erzeugung vollständig).
% @4 nutzt stattdessen l3pdfannot, das unter \cs{DocumentMetadata}
% zur Verfügung steht (aber ohne dieses selbst nicht -- ein
% nachträgliches \cs{RequirePackage}\code{\{pdfmanagement\}} nach hyperref
% scheitert an einem Hook-Sortierungskonflikt, die Schicht muss vor
% hyperref/xcolor aktiviert werden). @5 wählt zur Aufrufzeit anhand von
% \cs{IfDocumentMetadataTF} zwischen beiden. Der sichtbare Text wird in
% beiden Fällen über das reguläre (aber unverlinkte,
% \cs{ref}*-artige) \cs{\_\_olref\_original\_ref} aus der externen
% \file{.aux} aufgelöst -- nur die Verlinkung selbst wird
% ersetzt.}{A cross-unit @1 reference normally creates a plain
% \code{GoToR} link (jump into an external file) via hyperref's own
% \cs{hyper@linkfile} -- regardless of whether the target has since
% (via \cs{includeunit}) become a real part of the same combined
% document. tagpax carries every annotation over unchanged when merging
% PDFs and has no way to tell that a \code{GoToR} should now actually
% be an internal jump. @2 therefore no longer builds the link via
% hyperref's standard path (\cs{hyper@linkfile}), but additionally
% embeds the logical target identity (unit/doctype/language) plus the
% source page as custom keys on the annotation -- the same identity
% \cs{includeunit} already knows. This lets tagpax (or an extension of
% it) recognize from the link itself whether the target is available
% internally, without fragile file-path comparisons. Two
% implementations are needed, depending on whether
% \cs{DocumentMetadata} is active for the current target (in
% OSGLecture this is target-dependent, see \code{document\_metadata} in
% \file{ollmconfig.toml}): @3 duplicates hyperref's classic
% \cs{hyper@linkfile} structure directly
% (\cs{pdfstartlink}/\cs{close@pdflink}, LuaTeX-specific --
% OSGLecture is LuaTeX-only anyway), since these primitive wrappers are
% \emph{not} defined under \cs{DocumentMetadata} (the PDF management
% layer takes over link creation entirely there). @4 uses l3pdfannot
% instead, which is available under \cs{DocumentMetadata} (but not
% without it -- loading it after the fact via \cs{RequirePackage}\code{\{pdfmanagement\}} after hyperref fails with a hook-ordering conflict;
% the layer must be activated before hyperref/xcolor). @5 picks between
% the two at call time via \cs{IfDocumentMetadataTF}. The visible text
% is resolved in both cases from the external \file{.aux} via the
% regular (but unlinked, \cs{ref}*-like)
% \cs{\_\_olref\_original\_ref} -- only the link creation itself is
% replaced.}{\cs{olref}}{\cs{\_\_olref\_link\_goToR:nnnn}}
% {\cs{\_\_olref\_link\_goToR\_classic:nnnn}}
% {\cs{\_\_olref\_link\_goToR\_pdfannot:nnnn}}
% {\cs{\_\_olref\_link\_goToR:nnnn}}
% \ldeen{\#1 = PDF-Pfad (xr-URL), \#2 = Zielname, \#3 = Quellseite,
% \#4 = Anzeigetext}{\#1 = PDF path (xr URL), \#2 = destination name,
% \#3 = source page, \#4 = display text}
\cs_new_protected:Npn \__olref_link_goToR_classic:nnnn #1#2#3#4
  {
    \begingroup
      \def\Hy@pstringF{#1}
      \Hy@CleanupFile\Hy@pstringF
      \Hy@pstringdef\Hy@pstringF\Hy@pstringF
      \Hy@pstringdef\Hy@pstringD{#2}
      \leavevmode
      \pdfstartlink
        attr{%
          /Border[\@pdfborder]%
          /OSGLectureUnit(\l__olref_unit_tl)%
          /OSGLectureType(\l__olref_type_tl)%
          /OSGLectureLang(\l__olref_lang_tl)%
          /OSGLectureSourcePage(#3)%
        }%
        user{%
          /Subtype/Link%
          /A<<%
            /F(\Hy@pstringF)%
            /S/GoToR%
            /D(\Hy@pstringD)%
          >>%
        }%
        \relax
      \Hy@colorlink\@filecolor #4\Hy@xspace@end
      \close@pdflink
    \endgroup
  }
% \ldeen{\#1 = PDF-Pfad (xr-URL), \#2 = Zielname, \#3 = Quellseite,
% \#4 = Anzeigetext}{\#1 = PDF path (xr URL), \#2 = destination name,
% \#3 = source page, \#4 = display text}
\cs_new_protected:Npn \__olref_link_goToR_pdfannot:nnnn #1#2#3#4
  {
    \group_begin:
      \def\Hy@pstringF{#1}
      \Hy@CleanupFile\Hy@pstringF
      \Hy@pstringdef\Hy@pstringF\Hy@pstringF
      \Hy@pstringdef\Hy@pstringD{#2}
      \pdfannot_dict_put:nne {link/GoToR} {Border} {[\@pdfborder]}
      \pdfannot_dict_put:nne {link/GoToR} {OSGLectureUnit} {(\l__olref_unit_tl)}
      \pdfannot_dict_put:nne {link/GoToR} {OSGLectureType} {(\l__olref_type_tl)}
      \pdfannot_dict_put:nne {link/GoToR} {OSGLectureLang} {(\l__olref_lang_tl)}
      \pdfannot_dict_put:nnn {link/GoToR} {OSGLectureSourcePage} {(#3)}
      \pdfannot_link:nnn { GoToR }
        {
          /A<<
            /F(\Hy@pstringF)
            /S/GoToR
            /D(\Hy@pstringD)
          >>
        }
        { \Hy@colorlink \@filecolor #4 \Hy@endcolorlink \Hy@xspace@end }
    \group_end:
  }
\cs_new_protected:Npn \__olref_link_goToR:nnnn #1#2#3#4
  {
    \IfDocumentMetadataTF
      { \__olref_link_goToR_pdfannot:nnnn {#1}{#2}{#3}{#4} }
      { \__olref_link_goToR_classic:nnnn {#1}{#2}{#3}{#4} }
  }
% \ldeen{\#1 = präfigiertes Label; unterscheidet danach, ob xr es
% kennt}{\#1 = prefixed label; dispatches on whether xr knows it}
\cs_new_protected:Npn \__olref_ref_external:n #1
  {
    \cs_if_exist:cTF { r@#1 }
      {
        \tl_set_eq:Nc \l_tmpa_tl { r@#1 }
        \exp_after:wN \__olref_ref_external_fields:nnnnnn
          \l_tmpa_tl {#1}
      }
      { \__olref_original_ref {#1} }
  }
% \ldeen{\#1 Nummer \#2 Seite \#3 Titel \#4 Ankername \#5 xr-URL
% \#6 Original-Label}{\#1 number \#2 page \#3 title \#4 anchor-name
% \#5 xr-url \#6 orig label}
\cs_new_protected:Npn \__olref_ref_external_fields:nnnnnn #1#2#3#4#5#6
  {
    \tl_if_empty:nTF {#5}
      { \__olref_original_ref {#6} }
      { \__olref_link_goToR:nnnn {#5}{#4}{#2}{#1} }
  }
\cs_new_protected:Npn \__olref_ref:n #1
  {
    \tl_if_empty:NTF \l__olref_unit_tl
      { \__olref_original_ref {#1} }
      { \__olref_ref_external:n {#1} }
  }
\cs_new_protected:Npn \__olref_ref_star:n #1
  { \__olref_original_ref * {#1} }
\cs_new_protected:Npn \__olref_pageref:n #1
  { \__olref_original_pageref {#1} }
\cs_new_protected:Npn \__olref_pageref_star:n #1
  { \__olref_original_pageref * {#1} }
\cs_new_protected:Npn \__olref_nameref:n #1
  { \__olref_original_nameref {#1} }
\cs_new_protected:Npn \__olref_nameref_star:n #1
  { \__olref_original_nameref * {#1} }
\cs_new_protected:Npn \__olref_autoref:n #1
  { \__olref_original_autoref {#1} }
\cs_new_protected:Npn \__olref_autoref_star:n #1
  { \__olref_original_autoref * {#1} }

\NewDocumentCommand \olref { s O{} m }
  {
    \__olref_target_set:nnn {#2}{#3}{ref}
    \IfBooleanTF{#1}
      { \exp_args:NV \__olref_ref_star:n \l__olref_label_tl }
      { \exp_args:NV \__olref_ref:n \l__olref_label_tl }
  }
\NewDocumentCommand \olpageref { s O{} m }
  {
    \__olref_target_set:nnn {#2}{#3}{page}
    \IfBooleanTF{#1}
      { \exp_args:NV \__olref_pageref_star:n \l__olref_label_tl }
      { \exp_args:NV \__olref_pageref:n \l__olref_label_tl }
  }
\NewDocumentCommand \olnameref { s O{} m }
  {
    \__olref_target_set:nnn {#2}{#3}{name}
    \IfBooleanTF{#1}
      { \exp_args:NV \__olref_nameref_star:n \l__olref_label_tl }
      { \exp_args:NV \__olref_nameref:n \l__olref_label_tl }
  }
\NewDocumentCommand \olautoref { s O{} m }
  {
    \__olref_target_set:nnn {#2}{#3}{autoref}
    \IfBooleanTF{#1}
      { \exp_args:NV \__olref_autoref_star:n \l__olref_label_tl }
      { \exp_args:NV \__olref_autoref:n \l__olref_label_tl }
  }

\cs_new_protected:Npn \__olref_install_replacements:
  {
    \bool_if:NF \g__olref_replacements_installed_bool
      {
        \clist_if_empty:NF \g__olref_replace_clist
          {
            \msg_warning:nnx {osglecture-references}{replace-standard}
              {\clist_use:Nn \g__olref_replace_clist {,~}}
          }
        \clist_if_in:NnT \g__olref_replace_clist {ref}
          {
            \RenewDocumentCommand \ref { s O{} m }
              {
                \IfBooleanTF{##1}
                  {\olref*[##2]{##3}}
                  {\olref[##2]{##3}}
              }
          }
        \clist_if_in:NnT \g__olref_replace_clist {pageref}
          {
            \RenewDocumentCommand \pageref { s O{} m }
              {
                \IfBooleanTF{##1}
                  {\olpageref*[##2]{##3}}
                  {\olpageref[##2]{##3}}
              }
          }
        \clist_if_in:NnT \g__olref_replace_clist {nameref}
          {
            \RenewDocumentCommand \nameref { s O{} m }
              {
                \IfBooleanTF{##1}
                  {\olnameref*[##2]{##3}}
                  {\olnameref[##2]{##3}}
              }
          }
        \clist_if_in:NnT \g__olref_replace_clist {autoref}
          {
            \RenewDocumentCommand \autoref { s O{} m }
              {
                \IfBooleanTF{##1}
                  {\olautoref*[##2]{##3}}
                  {\olautoref[##2]{##3}}
              }
          }
        \bool_gset_true:N \g__olref_replacements_installed_bool
      }
  }

\cs_new_protected:Npn \__olref_install_legacy:
  {
    \bool_if:NT \g__olref_legacy_bool
      {
        \ProvideDocumentCommand \xref { O{} m } { \olref[##1]{##2} }
        \ProvideDocumentCommand \xrefchap { m }
          { \olref[##1]{osglecture} }
        \ProvideDocumentCommand \xrefsmart { s O{} m }
          {
            \IfBooleanTF{##1}
              { \olpageref*[##2]{##3} }
              { \olpageref[##2]{##3} }
          }
        \ProvideDocumentCommand \xrefdist { s O{} m m }
          { \__olref_legacy_dist:nnnn {##1}{##2}{##3}{##4} }
        \ProvideDocumentCommand \xarticleref { O{} m }
          { \olref[##1,type=script]{##2} }
        \ProvideDocumentCommand \xpresentationref { O{} m }
          { \olref[##1,type=slides]{##2} }
      }
  }

\cs_new_protected:Npn \__olref_legacy_dist_text:nnn #1#2#3
  {
    \IfBooleanT{#1}{(}
    #2~#3
    \IfBooleanT{#1}{)}
  }

\cs_new_protected:Npn \__olref_legacy_dist:nnnn #1#2#3#4
  {
    \group_begin:
      \cs_set:Npn \reftextfaceafter
        { \__olref_legacy_dist_text:nnn {#1}{#4}{\thevpagerefnum} }
      \cs_set_eq:NN \reftextafter \reftextfaceafter
      \cs_set_eq:NN \reftextfacebefore \reftextfaceafter
      \cs_set_eq:NN \reftextbefore \reftextfaceafter
      \cs_set:Npn \reftextfaraway ##1
        {
          \__olref_legacy_dist_text:nnn
            {#1}{#4}{\__olref_original_pageref{##1}}
        }
      \IfBooleanTF{#1}
        { \vpageref*[#2]{#3} }
        { \vpageref[#2]{#3} }
    \group_end:
  }

\cs_new_protected:Npn \__olref_export_label:n #1
  {
    \iow_now:Ne \g__olref_export_iow
      {
        \token_to_str:N \newlabel
        { \exp_not:n {#1} }
          {
            { \exp_not:o {\@currentlabel} }
            { \thepage }
            { \exp_not:o {\@currentlabelname} }
            { \exp_not:o {\@currentHref} }
            {}
          }
      }
  }

\cs_new_protected:Npn \__olref_export_finish:
  {
    \OsgLectureBuildLoadedTF
      {
        \tl_if_empty:eF { \OsgLectureBuildValue{generation-id} }
          {
            \iow_now:Ne \g__olref_export_iow
              {
                \token_to_str:N \OsgLectureReferenceExport
                {\OsgLectureBuildValue{schema}}
                {\OsgLectureBuildValue{generation-id}}
                {\OsgLectureBuildValue{job-id}}
                {\OsgLectureBuildValue{series-id}}
                {\OsgLectureUnitId}
                {\OsgLectureBuildValue{doctype}}
                {\OsgLectureBuildValue{language}}
              }
          }
      }
      { }
    \iow_close:N \g__olref_export_iow
    \OsgLectureBuildLoadedTF
      {
        \tl_if_empty:eF { \OsgLectureBuildValue{generation-id} }
          { \iow_close:N \g__olref_used_iow }
      }
      { }
  }

\AddToHook{begindocument/before}[osglecture-references/setup]
  {
    \file_if_exist:nT {\jobname.osgrefs.tex}
      { \file_input:n {\jobname.osgrefs.tex} }
    \iow_open:Nn \g__olref_export_iow {\jobname.osgref.aux}
    \OsgLectureBuildLoadedTF
      {
        \tl_if_empty:eF { \OsgLectureBuildValue{generation-id} }
          { \iow_open:Nn \g__olref_used_iow {\jobname.osgref-used.aux} }
      }
      { }
    \__olref_install_legacy:
    \__olref_install_replacements:
  }
\AddToHookWithArguments{label}[osglecture-references/export]
  { \__olref_export_label:n {#1} }
\AddToHook{enddocument/afterlastpage}[osglecture-references/export]
  { \__olref_export_finish: }

\ExplSyntaxOff
%</references>
%<*continuation>

\RequirePackage{osglecture-project}
\RequirePackage{osglecture-series}
\RequirePackage{osglecture-references}

\ExplSyntaxOn

\clist_new:N \g__osglecture_continuation_counters_clist
\tl_new:N \l__osglecture_continuation_physical_tl
\tl_new:N \l__osglecture_continuation_unit_tl
\tl_new:N \l__osglecture_continuation_aux_tl
\prop_new:N \g__osglecture_continuation_values_prop

\msg_new:nnn { osglecture-continuation } { unknown-previous-unit }
  {
    Previous~physical~unit~'#1'~has~no~known~logical~identity.~Build~that~unit~
    once~before~continuing~its~counters;~using~the~normal~initial~values.
  }
\msg_new:nnn { osglecture-continuation } { missing-projection }
  {
    Previous~logical~unit~'#1'~has~no~promoted~projection~for~the~current~
    document~type~and~language;~using~the~normal~initial~counter~values.
  }
\msg_new:nnn { osglecture-continuation } { unknown-counter }
  { Cannot~continue~unknown~LaTeX~counter~'#1'. }

\keys_define:nn { osglecture/setup/continuation }
  {
    counters .clist_gset:N = \g__osglecture_continuation_counters_clist,
    none .code:n = { \clist_gclear:N \g__osglecture_continuation_counters_clist },
    none .value_forbidden:n = true,
  }
\cs_new_protected:Npn \__osglecture_continuation_setup:n #1
  { \keys_set:nn { osglecture/setup/continuation } {#1} }
\DeclareOsgLectureModeSetupArea{continuation}
  {\__osglecture_continuation_setup:n}

\NewExpandableDocumentCommand \OsgLectureContinuationCounters { }
  { \clist_use:Nn \g__osglecture_continuation_counters_clist {,} }

\cs_new_protected:Npn \__osglecture_continuation_restore:nn #1#2
  {
    \clist_if_in:NnT \g__osglecture_continuation_counters_clist {#1}
      {
        \cs_if_exist:cTF { c@#1 }
          { \prop_gput:Nnn \g__osglecture_continuation_values_prop {#1}{#2} }
          { \msg_warning:nnn {osglecture-continuation}{unknown-counter}{#1} }
      }
  }

\cs_new_protected:Npn \__osglecture_continuation_apply:
  {
    \prop_map_inline:Nn \g__osglecture_continuation_values_prop
      { \int_gset:cn { c@##1 } {##2} }
  }

\cs_new_protected:Npn \__osglecture_continuation_input:
  {
    \OsgLectureReferenceDocumentAuxTo
      {\l__osglecture_continuation_unit_tl}
      {\OsgLectureDocumentType}
      {\OsgLectureLanguage}
      \l__osglecture_continuation_aux_tl
    \group_begin:
      \cs_set_protected:Npn \newlabel ##1##2 { }
      \cs_set_protected:Npn \OsgLectureReferenceExport
        ##1##2##3##4##5##6##7 { }
      \cs_set_eq:NN \OsgLectureContinuationCounter
        \__osglecture_continuation_restore:nn
      \exp_args:NV \file_input:n \l__osglecture_continuation_aux_tl
    \group_end:
  }

\cs_new_protected:Npn \__osglecture_continuation_import:
  {
    \clist_if_empty:NF \g__osglecture_continuation_counters_clist
      {
        \tl_set:Ne \l__osglecture_continuation_physical_tl
          { \OsgLectureSeriesPreviousContinuationPhysicalUnit }
        \tl_if_empty:NF \l__osglecture_continuation_physical_tl
          {
            \tl_set:Ne \l__osglecture_continuation_unit_tl
              {
                \exp_args:NV \OsgLectureSeriesLogicalUnit
                  \l__osglecture_continuation_physical_tl
              }
            \tl_if_empty:NTF \l__osglecture_continuation_unit_tl
              {
                \msg_warning:nnV {osglecture-continuation}
                  {unknown-previous-unit}
                  \l__osglecture_continuation_physical_tl
              }
              {
                \OsgLectureReferenceDocumentTF
                  {\l__osglecture_continuation_unit_tl}
                  {\OsgLectureDocumentType}
                  {\OsgLectureLanguage}
                  { \__osglecture_continuation_input: }
                  {
                    \msg_warning:nnV {osglecture-continuation}
                      {missing-projection}
                      \l__osglecture_continuation_unit_tl
                  }
              }
          }
      }
  }

\cs_new_protected:Npn \__osglecture_continuation_export:
  {
    \clist_map_inline:Nn \g__osglecture_continuation_counters_clist
      {
        \cs_if_exist:cTF { c@##1 }
          {
            % \ldeen{Der Export erfolgt nach dem letzten Shipout, daher
            % enthält der \code{page}-Zähler dort die zuletzt
            % ausgelieferte Seitenzahl; für \code{page} wird deshalb
            % ausdrücklich der um eins erhöhte Wert exportiert, damit
            % eine folgende einseitige Unit auf der nächsten Seite
            % beginnt.}{The export happens after the last shipout, so
            % the \code{page} counter holds the number of the last
            % emitted page at that point; for \code{page} the exported
            % value is therefore explicitly incremented by one, so a
            % following single-sided unit starts on the next
            % page.}
            \exp_args:Ne \OsgLectureWriteReferenceExport
              {
                \token_to_str:N \OsgLectureContinuationCounter
                {\exp_not:n {##1}}
                {
                  \str_if_eq:nnTF {##1}{page}
                    { \int_eval:n { \int_use:c {c@page} + 1 } }
                    { \int_use:c {c@##1} }
                }
              }
          }
          { \msg_warning:nnn {osglecture-continuation}{unknown-counter}{##1} }
      }
  }

% \ldeen{Die Hook-Reihenfolge ist bewusst festgelegt: Der Import muss nach
% dem Setup von \pkg*{osglecture-references} laufen (dessen
% \cs{jobname}\code{.osgrefs.tex} eingelesen sein muss), und der Export muss
% vor dem Referenz-Export laufen, damit dessen Datei bereits offen bzw.
% konsistent ist.}{The hook ordering is deliberate: the import must run
% after \pkg*{osglecture-references}'s setup (whose
% \cs{jobname}\code{.osgrefs.tex} must already be read), and the export must
% run before the reference export, so its file is already open and
% consistent.}
\AddToHook{begindocument/before}[osglecture-continuation/import]
  { \__osglecture_continuation_import: }
\DeclareHookRule{begindocument/before}{osglecture-continuation/import}
  {after}{osglecture-references/setup}
\AddToHook{osglecture/lecture/after}[osglecture-continuation/apply]
  { \__osglecture_continuation_apply: }
\AddToHook{enddocument/afterlastpage}[osglecture-continuation/export]
  { \__osglecture_continuation_export: }
\DeclareHookRule{enddocument/afterlastpage}{osglecture-continuation/export}
  {before}{osglecture-references/export}

\ExplSyntaxOff
%</continuation>
%<*integration>

\RequirePackage{osglecture-config}
\RequirePackage{osglecture-references}
\RequirePackage{tagpax}
% \ldeen*{Werden unbedingt geladen, nicht erst wenn ein ungetaggter Doctype
% tatsächlich auftritt: beide Pakete müssen vor \cs{begin}\Marg{document}
% stehen, der tatsächliche Tagging-Status einer Unit ist aber erst zur
% Laufzeit von \cs{includeunit} bekannt. @1 nutzt \cs{tagpax} nicht mit
% dessen eigener \code{pdfpages}-Option, daher entsteht kein Konflikt mit
% dem echten @2.}{Loaded unconditionally, not only once an untagged doctype
% actually occurs: both packages must be loaded before \cs{begin}\Marg{document},
% but a unit's actual tagging status is known only when \cs{includeunit}
% runs. @1 does not load \cs{tagpax} with its own \code{pdfpages} option, so
% there is no conflict with the real @2.}{\code{osglecture-integration}}{\pkg{pdfpages}}
\RequirePackage{pdfpages}
\RequirePackage{newpax}

\ExplSyntaxOn

\tl_new:N \l__osglecture_integration_pdf_tl
\tl_new:N \l__osglecture_integration_ir_tl
\tl_new:N \l__osglecture_integration_generation_tl
\tl_new:N \l__osglecture_integration_prefix_tl
\tl_new:N \l__osglecture_integration_mode_tl
\tl_new:N \l__osglecture_integration_tagged_word_tl
\tl_new:N \l__osglecture_integration_mixed_units_tl
\tl_new:N \l__osglecture_integration_mixed_doctypes_tl
\prop_new:N \g__osglecture_integration_mode_prop
\seq_new:N \g__osglecture_integration_mixed_seq

% \ldeen*{@1 kennt für jeden \cs{olref}-Link auf eine andere Einheit
% deren logische Identität (Unit/Doctyp/Sprache), aber nicht, ob diese
% Einheit im aktuellen Dokument tatsächlich per \cs{includeunit}
% eingebunden wird oder wurde -- das weiß nur @2. @3 macht diese
% Information für tagpax verfügbar: einmal pro Lauf aus der im
% vorigen Lauf geschriebenen Registrierungsdatei (löst
% Rückwärts- \emph{und}, nach einem weiteren Lauf,
% Vorwärtsreferenzen auf), und live nach jedem
% \cs{includeunit}-Aufruf in diesem Lauf.}{@1 knows, for every
% \cs{olref} link to another unit, that unit's logical identity
% (unit/doctype/language), but not whether that unit is actually (or
% will be) included in the current document via \cs{includeunit} --
% only @2 knows that. @3 makes this information available to tagpax:
% once per run from the registration file written on the previous run
% (resolves backward \emph{and}, after one more run, forward
% references), and live after every \cs{includeunit} call in this
% run.}{\cs{olref}}{\cs{includeunit}}{\code{osglecture-integration.lua}}
\directlua{require("osglecture-integration").load_registry("\luaescapestring{\jobname}")}

\msg_new:nnn {osglecture-integration}{missing-projection}
  {
    Logical~unit~'#1'~has~no~promoted~projection~for~the~current~document~type~
    and~language.~Build~that~unit~before~integrating~it.
  }
\msg_new:nnn {osglecture-integration}{mixed-tagging}
  {
    Doctype~'#1'~mixed~tagged~and~untagged~source~PDFs~in~the~previous~run~
    (previously~included~units:~#2).~Falling~back~to~the~untagged~inclusion~
    path~for~every~unit~of~this~doctype~in~this~run:~cross-unit~navigation~
    is~best-effort~and~the~accessibility~structure~is~degraded~for~the~whole~
    doctype.~Rebuild~every~unit~of~'#1'~as~a~tagged~PDF~to~restore~the~
    native~import~path.
  }
\msg_new:nnn {osglecture-integration}{mixed-tagging-summary}
  {
    Doctype(s)~with~mixed~tagged/untagged~source~PDFs~in~the~previous~run:~
    #1.~These~were~built~through~the~untagged~fallback~path~in~this~run;~
    see~earlier~warnings~for~details.~Rebuild~every~unit~of~the~affected~
    doctype(s)~as~a~tagged~PDF~and~rerun~to~restore~the~native~import~path.
  }

% \ldeen*{Sorgt bei Bedarf dafür, dass @2 einen Eintrag für @1 enthält --
% ausgehend vom im vorigen Lauf aggregierten Status, inklusive Warnung und
% erzwungenem \code{untagged} im Mix-Fall. Ist bereits ein Eintrag
% vorhanden (auch ein diesen Lauf live gesetzter), geschieht nichts. Bleibt
% @1 danach ohne Eintrag, kannte der vorige Lauf keine Unit dieses
% Doctypes; der Aufrufer muss diesen Bootstrap-Fall live anhand der ersten
% tatsächlichen Unit auflösen. Als eigenständiger protokollierter Befehl
% statt als Teil einer scheinbar rein expandierbaren Abfrage, weil die
% Entscheidung Seiteneffekte (Warnung, Zustandsschreiben) hat, die in
% einem Expansionskontext nicht ausgeführt würden.}{Ensures, when needed,
% that @2 has an entry for @1 -- derived from the previous run's
% aggregated status, including the warning and the forced \code{untagged}
% mode in the mixed case. Does nothing if an entry already exists
% (including one set live earlier in this run). If @1 still has no entry
% afterwards, the previous run knew no unit of this doctype; the caller
% must resolve that bootstrap case live from the first actual unit. A
% standalone protected command rather than part of a seemingly pure
% lookup, because the decision has side effects (warning, state write)
% that would not execute in an expansion context.}{\meta{doctype}}
% {\cs{g\_\_osglecture\_integration\_mode\_prop}}
\cs_new_protected:Npn \__osglecture_integration_decide_mode:n #1
  {
    \prop_if_in:NnF \g__osglecture_integration_mode_prop {#1}
      {
        \tl_set:Ne \l__osglecture_integration_mode_tl
          {
            \directlua{
              tex.sprint(require("osglecture-integration").previous_mode(
                "\luaescapestring{#1}"))
            }
          }
        \tl_if_eq:VnT \l__osglecture_integration_mode_tl {mixed}
          {
            \tl_set:Ne \l__osglecture_integration_mixed_units_tl
              {
                \directlua{
                  tex.sprint(require("osglecture-integration").previous_units_of(
                    "\luaescapestring{#1}"))
                }
              }
            \msg_warning:nnnn {osglecture-integration}{mixed-tagging}
              {#1} {\l__osglecture_integration_mixed_units_tl}
            \seq_if_in:NnF \g__osglecture_integration_mixed_seq {#1}
              { \seq_gput_right:Nn \g__osglecture_integration_mixed_seq {#1} }
            \tl_set:Nn \l__osglecture_integration_mode_tl {untagged}
          }
        \tl_if_eq:VnT \l__osglecture_integration_mode_tl {tagged}
          { \prop_gput:NnV \g__osglecture_integration_mode_prop {#1} \l__osglecture_integration_mode_tl }
        \tl_if_eq:VnT \l__osglecture_integration_mode_tl {untagged}
          { \prop_gput:NnV \g__osglecture_integration_mode_prop {#1} \l__osglecture_integration_mode_tl }
      }
  }

% \ldeen{Getaggter Pfad: nativer tagpax-Import inklusive des für
% spätere Cross-Referenzen benötigten Präfixes.}{Tagged path: native
% tagpax import, including the prefix needed for later
% cross-references.}
\cs_new_protected:Npn \__osglecture_integration_include_tagged:nn #1#2
  {
    \tl_set:Ne \l__osglecture_integration_ir_tl
      { \jobname-osglecture-#1.tagpax }
    \tagpaxextract[\l__osglecture_integration_ir_tl] {#2}
    \tagpaxinclude[ir=\l__osglecture_integration_ir_tl] {#2}
    \tl_set:Ne \l__osglecture_integration_prefix_tl {\tagpaxlastprefix}
  }

% \ldeen*{Ungetaggter Rückfallpfad: die Seiten werden mit der
% \code{artifact}-Option von \pkg{pdfpages} als Artefakte importiert, damit
% der \code{StructTreeRoot} des Masterdokuments gültig
% bleibt -- der importierte Block ist bewusst nicht barrierefrei
% erschlossen. \pkg{newpax} rekonstruiert die in der Quelle bereits
% vorhandenen Links und das Lesezeichen-/Gliederungsverzeichnis auf die
% neuen Seitenpositionen; das ergibt ein klickbares Inhaltsverzeichnis auch
% ohne Struct-Tree, aber keine Vorlesereihenfolge oder
% Überschriftenebenen.}{Untagged fallback path: pages are imported as
% artifacts using \pkg{pdfpages}' \code{artifact} option so the master
% document's \code{StructTreeRoot} stays
% valid -- the imported block is deliberately not made accessible.
% \pkg{newpax} reconstructs the links and the bookmark/outline entries
% already present in the source onto the new page positions, giving a
% clickable table of contents even without a struct tree, but no reading
% order or heading levels.}{}
\cs_new_protected:Npn \__osglecture_integration_include_untagged:nn #1#2
  {
    \group_begin:
      % pdfpages uses non-final \includegraphics calls to inspect the PDF.
      % Older latex-lab versions do not yet recognize pdfpages' status flag
      % and otherwise warn about missing alternative text for such a probe.
      % The final call is independently marked by pdfpages' artifact option.
      \cs_if_exist:NT \l_tag_graphic_mode_tl
        { \tl_set:Nn \l_tag_graphic_mode_tl {artifact} }
      \includepdf[pages=-,artifact] {#2}
    \group_end:
    \tl_clear:N \l__osglecture_integration_prefix_tl
  }

\cs_new_protected:Npn \__osglecture_integration_include:n #1
  {
    \OsgLectureReferenceDocumentTF
      {#1}{\OsgLectureDocumentType}{\OsgLectureLanguage}
      {
        \OsgLectureReferenceDocumentPdfTo
          {#1}{\OsgLectureDocumentType}{\OsgLectureLanguage}
          \l__osglecture_integration_pdf_tl
        \OsgLectureReferenceDocumentGenerationTo
          {#1}{\OsgLectureDocumentType}{\OsgLectureLanguage}
          \l__osglecture_integration_generation_tl
        \__osglecture_integration_decide_mode:n {\OsgLectureDocumentType}
        \prop_if_in:NnTF \g__osglecture_integration_mode_prop {\OsgLectureDocumentType}
          {
            \tl_set:Ne \l__osglecture_integration_mode_tl
              { \prop_item:Nn \g__osglecture_integration_mode_prop {\OsgLectureDocumentType} }
          }
          {
            % \ldeen{Bootstrap: Der vorige Lauf kannte keine Unit dieses
            % Doctypes, und in diesem Lauf wurde bislang auch keine
            % eingebunden. Das aktuelle PDF dieser Unit wird live geprüft
            % und legt den Modus für den Rest des Laufs fest, analog zum
            % Referenzauflösungs-Bootstrap in Abschnitt~14.3 des
            % Architekturdokuments.}{Bootstrap: the previous run knew no
            % unit of this doctype, and no unit of it has been included
            % earlier in this run either. This unit's own current PDF is
            % probed live and establishes the mode for the rest of the
            % run, exactly like the reference-resolution bootstrap case
            % in section 14.3 of the architecture document.}
            \tl_set:Ne \l__osglecture_integration_mode_tl
              {
                \directlua{
                  tex.sprint(require("osglecture-integration").probe_tagged(
                    "\luaescapestring{\l__osglecture_integration_pdf_tl}")
                    and "tagged" or "untagged")
                }
              }
            \prop_gput:NnV \g__osglecture_integration_mode_prop
              {\OsgLectureDocumentType} \l__osglecture_integration_mode_tl
          }
        \tl_if_eq:VnTF \l__osglecture_integration_mode_tl {tagged}
          {
            \__osglecture_integration_include_tagged:nn {#1}{\l__osglecture_integration_pdf_tl}
            \tl_set:Nn \l__osglecture_integration_tagged_word_tl {true}
          }
          {
            \__osglecture_integration_include_untagged:nn {#1}{\l__osglecture_integration_pdf_tl}
            \tl_set:Nn \l__osglecture_integration_tagged_word_tl {false}
          }
        \directlua{
          require("osglecture-integration").register_unit(
            "\luaescapestring{#1}",
            "\luaescapestring{\OsgLectureDocumentType}",
            "\luaescapestring{\OsgLectureLanguage}",
            "\luaescapestring{\l__osglecture_integration_prefix_tl}",
            \l__osglecture_integration_tagged_word_tl)
        }
        \exp_args:Ne \OsgLectureWriteReferenceUse
          {
            \token_to_str:N \OsgLectureIntegrationUse
            {\OsgLectureBuildValue{generation-id}}
            {\exp_not:V \l__osglecture_integration_generation_tl}
            {\exp_not:n {#1}}
            {\OsgLectureDocumentType}
            {\OsgLectureLanguage}
          }
      }
      { \msg_error:nnn {osglecture-integration}{missing-projection}{#1} }
  }

\NewDocumentCommand \includeunit { m }
  { \__osglecture_integration_include:n {#1} }

\AddToHook{enddocument}
  {
    \seq_if_empty:NF \g__osglecture_integration_mixed_seq
      {
        \tl_set:Ne \l__osglecture_integration_mixed_doctypes_tl
          { \seq_use:Nn \g__osglecture_integration_mixed_seq {,~} }
        \msg_warning:nnn {osglecture-integration}{mixed-tagging-summary}
          {\l__osglecture_integration_mixed_doctypes_tl}
      }
  }

\ExplSyntaxOff
%</integration>
%<*structure>

\ExplSyntaxOn

\tl_new:N \g__osglecture_unit_id_tl
\tl_new:N \g__osglecture_unit_title_tl
\tl_new:N \g__osglecture_unit_short_title_tl
\tl_new:N \g__osglecture_deployment_chapter_tl
\bool_new:N \g__osglecture_unit_declared_bool
\bool_new:N \g__osglecture_unit_valid_bool
\iow_new:N \g__osglecture_result_iow

\msg_new:nnnn { osglecture-structure } { invalid-unit }
  { Invalid~logical~unit~identifier~'#1'. }
  { Unit~identifiers~must~not~contain~whitespace,~comma,~equals,~or~'...'. }
\msg_new:nnnn { osglecture-structure } { duplicate-unit }
  { The~document~declares~more~than~one~lecture/unit. }
  { A~series~unit~must~contain~exactly~one~canonical~lecture~declaration. }
\msg_new:nnnn { osglecture-structure } { missing-unit }
  { This~series~build~does~not~declare~a~logical~unit. }
  { Add~\token_to_str:N\lecture[short]\{title\}\{unit-id\}. }
\msg_new:nnnn { osglecture-structure } { unit-mismatch }
  { Declared~unit~'#1'~does~not~match~the~known~build~unit~'#2'. }
  { Keep~the~explicit~unit-id~stable~or~remove~the~stale~OLLM~state. }

\bool_gset_true:N \g__osglecture_unit_valid_bool
\NewHook{osglecture/lecture/after}

\cs_new_protected:Npn \__osglecture_unit_validate:n #1
  {
    \regex_match:nnTF { \A [^\s,=]+ \Z } {#1}
      {
        \tl_if_in:nnT {#1} {...}
          {
            \bool_gset_false:N \g__osglecture_unit_valid_bool
            \msg_error:nnn {osglecture-structure}{invalid-unit}{#1}
          }
      }
      {
        \bool_gset_false:N \g__osglecture_unit_valid_bool
        \msg_error:nnn {osglecture-structure}{invalid-unit}{#1}
      }
  }

\cs_new_protected:Npn \__osglecture_unit_declare:nnn #1#2#3
  {
    \bool_if:NTF \g__osglecture_unit_declared_bool
      {
        \bool_gset_false:N \g__osglecture_unit_valid_bool
        \msg_error:nn {osglecture-structure}{duplicate-unit}
      }
      {
        \__osglecture_unit_validate:n {#3}
        \tl_gset:Nn \g__osglecture_unit_short_title_tl {#1}
        \tl_gset:Nn \g__osglecture_unit_title_tl {#2}
        \tl_gset:Nn \g__osglecture_unit_id_tl {#3}
        \bool_gset_true:N \g__osglecture_unit_declared_bool
        \OsgLectureBuildLoadedTF
          {
            \tl_set:Ne \l_tmpa_tl { \OsgLectureBuildValue{unit-id} }
            \tl_if_empty:NF \l_tmpa_tl
              {
                \str_if_eq:VnF \l_tmpa_tl {#3}
                  {
                    \bool_gset_false:N \g__osglecture_unit_valid_bool
                    \msg_error:nnnn {osglecture-structure}{unit-mismatch}
                      {#3}{\l_tmpa_tl}
                  }
              }
          }
          { }
      }
  }

\cs_new_protected:Npn \__osglecture_lecture:nnn #1#2#3
  {
    \__osglecture_unit_declare:nnn {#1}{#2}{#3}
    \osglecture_adapter_lecture:nnn {#1}{#2}{#3}
    \UseHook{osglecture/lecture/after}
    % \ldeen*{Fängt die von der Basisklasse beim @1-Aufruf tatsächlich
    % vergebene Kapitelnummer ab (falls vorhanden), damit sie unten in
    % \code{osgresult.aux} exportiert werden kann. \cs{OsgLectureDeploymentChapter}
    % erlaubt es, diesen Wert manuell zu überschreiben, falls die Basisklasse
    % keinen \cs{thechapter}-Zähler kennt oder ein anderer Wert
    % maßgeblich ist.}{Captures the chapter number actually assigned by
    % the base class's @1 call (if any), so it can be exported below in
    % \code{osgresult.aux}. \cs{OsgLectureDeploymentChapter} allows this
    % value to be overridden manually when the base class has no
    % \cs{thechapter} counter, or a different value is
    % authoritative.}{\cs{lecture}}
    \tl_if_empty:NT \g__osglecture_deployment_chapter_tl
      {
        \cs_if_exist:NT \thechapter
          { \tl_gset:Ne \g__osglecture_deployment_chapter_tl {\thechapter} }
      }
  }

\NewDocumentCommand \OsgLectureDeploymentChapter { m }
  { \tl_gset:Nn \g__osglecture_deployment_chapter_tl {#1} }

\cs_if_exist:NTF \lecture
  { \RenewDocumentCommand \lecture { O{} m m }
      { \__osglecture_lecture:nnn {#1}{#2}{#3} } }
  { \NewDocumentCommand \lecture { O{} m m }
      { \__osglecture_lecture:nnn {#1}{#2}{#3} } }

\cs_new_protected:Npn \__osglecture_result_write:
  {
    \OsgLectureBuildLoadedTF
      {
        \tl_set:Ne \l_tmpa_tl { \OsgLectureBuildValue{generation-id} }
        \tl_if_empty:NF \l_tmpa_tl
          {
            \str_if_eq:eeF
              { \OsgLectureUnitRole } {i}
              {
                \bool_if:NF \g__osglecture_unit_declared_bool
                  {
                    \bool_gset_false:N \g__osglecture_unit_valid_bool
                    \msg_error:nn {osglecture-structure}{missing-unit}
                  }
              }
            \bool_if:NT \g__osglecture_unit_valid_bool
              {
                \tl_set:Ne \l_tmpb_tl { \OsgLectureSeriesPosition }
                % \ldeen*{Integrationsunits (Rolle @1) haben keine eigene
                % logische Unit-Id, kein eigenes Kapitel und keine eigene
                % Serienposition -- sie bündeln nur andere Units. Diese
                % Felder werden deshalb vor dem Export geleert statt die
                % zuletzt deklarierte reguläre Unit fälschlich
                % mitzuschleppen.}{Integration units (role @1) have no
                % logical unit id, chapter, or series position of their
                % own -- they only bundle other units. These fields are
                % therefore cleared before export instead of incorrectly
                % carrying over the last regular unit
                % declared.}{\code{i}}
                \str_if_eq:eeT
                  { \OsgLectureUnitRole } {i}
                  {
                    \tl_gclear:N \g__osglecture_unit_id_tl
                    \tl_gclear:N \g__osglecture_deployment_chapter_tl
                    \tl_clear:N \l_tmpb_tl
                  }
                \iow_open:Nn \g__osglecture_result_iow
                  {\jobname.osgresult.aux}
                \iow_now:Ne \g__osglecture_result_iow
                  {
                        \token_to_str:N \OsgLectureResult
                        {\OsgLectureBuildValue{schema}}
                        {\OsgLectureBuildValue{generation-id}}
                        {\OsgLectureBuildValue{job-id}}
                        {\OsgLectureBuildValue{series-id}}
                        {\exp_not:V \g__osglecture_unit_id_tl}
                        {\OsgLectureUnitPhysicalUnit}
                        {\OsgLectureUnitRole}
                        {\OsgLectureBuildValue{doctype}}
                        {\OsgLectureBuildValue{language}}
                  }
                \iow_now:Ne \g__osglecture_result_iow
                  {
                    \token_to_str:N \OsgLectureDeploymentResult
                    {\exp_not:V \g__osglecture_deployment_chapter_tl}
                    {\exp_not:V \l_tmpb_tl}
                  }
                \iow_close:N \g__osglecture_result_iow
              }
            }
      }
      { }
  }

\NewExpandableDocumentCommand \OsgLectureUnitId { }
  { \tl_use:N \g__osglecture_unit_id_tl }

\AddToHook{enddocument/afterlastpage}[osglecture-structure/result]
  { \__osglecture_result_write: }

\ExplSyntaxOff
%</structure>
%<*adapters>

\ExplSyntaxOn

\msg_new:nnnn { osglecture-adapters } { unknown-adapter }
  { Unknown~osglecture~adapter~'#1'. }
  { Install~osglecture-adapter-#1.def~or~select~another~document~profile. }
\msg_new:nnnn { osglecture-adapters } { incomplete-adapter }
  { Adapter~'#1'~does~not~provide~'#2'. }
  { Every~structure~adapter~must~implement~the~common~adapter~contract. }

\tl_new:N \g__osglecture_adapter_name_tl

% \ldeen*{Vertrag für jede \file{osglecture-adapter-*.def}: @1 ist
% Pflicht (überträgt \cs{lecture} auf die Zielbasisklasse); @2 ist
% optional und nur nötig, wenn der Adapter zusätzlich vorhandene Befehle
% der Zielklasse patchen oder deaktivieren muss (etwa
% \env{frame} bei den Langform-Adaptern).}{Contract for every
% \file{osglecture-adapter-*.def}: @1 is mandatory (translates
% \cs{lecture} to the target base class); @2 is optional and only needed
% when the adapter must additionally patch or disable commands already
% provided by the target class (e.g.\ \env{frame} for the long-form
% adapters).}{\cs{osglecture\_adapter\_lecture:nnn}}
% {\cs{osglecture\_adapter\_install:}}
\NewDocumentCommand \OsgLectureLoadAdapter { m }
  {
    \file_if_exist:nTF {osglecture-adapter-#1.def}
      {
        \tl_gset:Nn \g__osglecture_adapter_name_tl {#1}
        \file_input:n {osglecture-adapter-#1.def}
        \cs_if_exist:NF \osglecture_adapter_lecture:nnn
          { \msg_fatal:nnnn {osglecture-adapters}{incomplete-adapter}
              {#1}{lecture} }
        \cs_if_exist:NT \osglecture_adapter_install:
          { \osglecture_adapter_install: }
      }
      { \msg_fatal:nnn {osglecture-adapters}{unknown-adapter}{#1} }
  }

\ExplSyntaxOff
%</adapters>
%<*profiles>

\ExplSyntaxOn

\prop_new:N \g__osglecture_profiles_prop
\seq_new:N \g__osglecture_profile_classes_seq
\tl_new:N \l__osglecture_profile_backend_tl
\tl_new:N \l__osglecture_profile_adapter_tl
\tl_new:N \l__osglecture_profile_engine_tl
\tl_new:N \l__osglecture_profile_capabilities_tl
\tl_new:N \l__osglecture_profile_class_tl
\tl_new:N \l__osglecture_profile_doctypes_tl
\tl_new:N \l__osglecture_profile_options_tl
\tl_new:N \l__osglecture_profile_setup_tl
\tl_new:N \l__osglecture_profile_metadata_tl
\tl_new:N \l__osglecture_profile_modes_tl
\tl_new:N \l__osglecture_profile_mode_setup_tl
\tl_new:N \l__osglecture_profile_course_target_tl
\tl_new:N \l__osglecture_profile_event_target_tl
\cs_generate_variant:Nn \clist_if_in:nnF { Vn }

\msg_new:nnnn { osglecture-profiles } { unknown-profile }
  { Unknown~document~profile~'#1'. }
  { Install~or~declare~the~profile~before~selecting~it. }
\msg_new:nnnn { osglecture-profiles } { unsupported-doctype }
  { Document~profile~'#1'~does~not~support~document~type~'#2'. }
  { Select~a~profile~which~declares~that~document~type. }
\msg_new:nnnn { osglecture-profiles } { incomplete-profile }
  { Document~profile~'#1'~does~not~define~'#2'. }
  { Every~profile~must~define~backend,~class,~and~doctypes. }
\msg_new:nnnn { osglecture-profiles } { unsupported-engine }
  { Document~profile~'#1'~requires~unsupported~engine~'#2'. }
  { This~version~of~osglecture~supports~only~engine~'lualatex'. }

\keys_define:nn { osglecture/profile }
  {
    engine .tl_set:N = \l__osglecture_profile_engine_tl,
    engine .initial:n = lualatex,
    capabilities .clist_set:N = \l__osglecture_profile_capabilities_tl,
    backend .tl_set:N = \l__osglecture_profile_backend_tl,
    adapter .tl_set:N = \l__osglecture_profile_adapter_tl,
    class .tl_set:N = \l__osglecture_profile_class_tl,
    doctypes .clist_set:N = \l__osglecture_profile_doctypes_tl,
    class-options .tl_set:N = \l__osglecture_profile_options_tl,
    setup-file .tl_set:N = \l__osglecture_profile_setup_tl,
    modes .clist_set:N = \l__osglecture_profile_modes_tl,
    mode-setup-file .tl_set:N = \l__osglecture_profile_mode_setup_tl,
    course-target .choices:nn = { title, subtitle, none }
      { \tl_set_eq:NN \l__osglecture_profile_course_target_tl \l_keys_choice_tl },
    course-target .initial:n = none,
    event-target .choices:nn = { title, subtitle, none }
      { \tl_set_eq:NN \l__osglecture_profile_event_target_tl \l_keys_choice_tl },
    event-target .initial:n = none,
    document-metadata .choices:nn = { supported, required, forbidden }
      { \tl_set_eq:NN \l__osglecture_profile_metadata_tl \l_keys_choice_tl },
    document-metadata .initial:n = supported,
  }

\NewDocumentCommand \DeclareOsgLectureProfile { m m }
  {
    \tl_clear:N \l__osglecture_profile_backend_tl
    \tl_clear:N \l__osglecture_profile_adapter_tl
    \tl_set:Nn \l__osglecture_profile_engine_tl {lualatex}
    \tl_clear:N \l__osglecture_profile_capabilities_tl
    \tl_clear:N \l__osglecture_profile_class_tl
    \tl_clear:N \l__osglecture_profile_doctypes_tl
    \tl_clear:N \l__osglecture_profile_options_tl
    \tl_clear:N \l__osglecture_profile_setup_tl
    \tl_clear:N \l__osglecture_profile_modes_tl
    \tl_clear:N \l__osglecture_profile_mode_setup_tl
    \tl_set:Nn \l__osglecture_profile_course_target_tl {none}
    \tl_set:Nn \l__osglecture_profile_event_target_tl {none}
    \tl_set:Nn \l__osglecture_profile_metadata_tl {supported}
    \keys_set:nn { osglecture/profile } {#2}
    \tl_if_empty:NT \l__osglecture_profile_backend_tl
      { \msg_fatal:nnnn {osglecture-profiles}{incomplete-profile}{#1}{backend} }
    \tl_if_empty:NT \l__osglecture_profile_adapter_tl
      { \tl_set_eq:NN \l__osglecture_profile_adapter_tl
          \l__osglecture_profile_backend_tl }
    \tl_if_empty:NT \l__osglecture_profile_class_tl
      { \msg_fatal:nnnn {osglecture-profiles}{incomplete-profile}{#1}{class} }
    \tl_if_empty:NT \l__osglecture_profile_doctypes_tl
      { \msg_fatal:nnnn {osglecture-profiles}{incomplete-profile}{#1}{doctypes} }
    \prop_gput:NnV \g__osglecture_profiles_prop {#1/backend}
      \l__osglecture_profile_backend_tl
    \prop_gput:NnV \g__osglecture_profiles_prop {#1/adapter}
      \l__osglecture_profile_adapter_tl
    \prop_gput:NnV \g__osglecture_profiles_prop {#1/engine}
      \l__osglecture_profile_engine_tl
    \prop_gput:NnV \g__osglecture_profiles_prop {#1/capabilities}
      \l__osglecture_profile_capabilities_tl
    \prop_gput:NnV \g__osglecture_profiles_prop {#1/class}
      \l__osglecture_profile_class_tl
    \seq_if_in:NVF \g__osglecture_profile_classes_seq
      \l__osglecture_profile_class_tl
      { \seq_gput_right:NV \g__osglecture_profile_classes_seq
          \l__osglecture_profile_class_tl }
    \prop_gput:NnV \g__osglecture_profiles_prop {#1/doctypes}
      \l__osglecture_profile_doctypes_tl
    \prop_gput:NnV \g__osglecture_profiles_prop {#1/class-options}
      \l__osglecture_profile_options_tl
    \prop_gput:NnV \g__osglecture_profiles_prop {#1/setup-file}
      \l__osglecture_profile_setup_tl
    \prop_gput:NnV \g__osglecture_profiles_prop {#1/modes}
      \l__osglecture_profile_modes_tl
    \prop_gput:NnV \g__osglecture_profiles_prop {#1/mode-setup-file}
      \l__osglecture_profile_mode_setup_tl
    \prop_gput:NnV \g__osglecture_profiles_prop {#1/course-target}
      \l__osglecture_profile_course_target_tl
    \prop_gput:NnV \g__osglecture_profiles_prop {#1/event-target}
      \l__osglecture_profile_event_target_tl
    \prop_gput:NnV \g__osglecture_profiles_prop {#1/document-metadata}
      \l__osglecture_profile_metadata_tl
  }

\NewDocumentCommand \OsgLectureProfileClassKnownTF { m m m }
  { \seq_if_in:NnTF \g__osglecture_profile_classes_seq {#1} {#2} {#3} }

\NewDocumentCommand \DeclareOsgLectureProfileClassOptions { m m m }
  { \prop_gput:Nnn \g__osglecture_profiles_prop {#1/class-options/#2} {#3} }

\cs_new_protected:Npn \osglecture_profile_select:nn #1#2
  {
    \OsgLectureLoadProfile{#1}
    \prop_if_in:NnF \g__osglecture_profiles_prop {#1/backend}
      { \msg_fatal:nnn {osglecture-profiles}{unknown-profile}{#1} }
    \prop_get:NnN \g__osglecture_profiles_prop {#1/engine} \l_tmpb_tl
    \str_if_eq:VnF \l_tmpb_tl {lualatex}
      { \msg_fatal:nnnn {osglecture-profiles}{unsupported-engine}{#1}{\l_tmpb_tl} }
    \prop_get:NnN \g__osglecture_profiles_prop {#1/doctypes} \l_tmpa_tl
    \clist_if_in:VnF \l_tmpa_tl {#2}
      { \msg_fatal:nnnn {osglecture-profiles}{unsupported-doctype}{#1}{#2} }
  }

\NewDocumentCommand \OsgLectureLoadProfile { m }
  {
    \prop_if_in:NnF \g__osglecture_profiles_prop {#1/backend}
      {
        \file_if_exist:nT {osglecture-profile-#1.def}
          { \file_input:n {osglecture-profile-#1.def} }
      }
  }

\cs_new:Npn \osglecture_profile_value:nn #1#2
  { \prop_item:Nn \g__osglecture_profiles_prop {#1/#2} }

% \ldeen*{Liefert zwei durch \cs{prop\_item:Nn} nacheinander expandierte
% Werte, die der Aufrufer als eine kommaseparierte Optionsliste
% zusammensetzt: die allgemeinen Klassenoptionen des Profils sowie die
% (ggf. leeren) doctypspezifischen Zusatzoptionen aus @1. Zwei
% \cs{prop\_item:Nn}-Aufrufe statt eines zusammengesetzten Werts, weil
% Letzterer bei leerem doctypspezifischen Anteil ein unerwünschtes
% führendes Komma hinterließe.}{Yields two values expanded in sequence
% by \cs{prop\_item:Nn}, which the caller concatenates into one
% comma-separated option list: the profile's general class options plus
% the (possibly empty) doctype-specific extra options from @1. Two
% \cs{prop\_item:Nn} calls instead of one combined value, because the
% latter would leave an unwanted leading comma when the doctype-specific
% part is empty.}{\cs{DeclareOsgLectureProfileClassOptions}}
\cs_new:Npn \osglecture_profile_class_options:nn #1#2
  {
    \prop_item:Nn \g__osglecture_profiles_prop {#1/class-options}
    \prop_item:Nn \g__osglecture_profiles_prop {#1/class-options/#2}
  }

\cs_generate_variant:Nn \osglecture_profile_select:nn { VV }
\cs_generate_variant:Nn \osglecture_profile_value:nn { Vn }
\cs_generate_variant:Nn \osglecture_profile_class_options:nn { VV }

\file_input:n {osglecture-profile-beamer.def}
\file_input:n {osglecture-profile-ltx-talk.def}
\file_input:n {osglecture-profile-book.def}
\file_input:n {osglecture-profile-scrbook.def}

\ExplSyntaxOff
%</profiles>
%<*presitemize>

\RequirePackage{osglecture-modes}

\ExplSyntaxOn

\bool_new:N \l__osglecture_pres_longform_bool
\bool_new:N \l__osglecture_pres_subordinate_bool
\bool_new:N \l__osglecture_pres_has_item_bool
\bool_new:N \l__osglecture_pres_auto_punctuation_bool
\tl_new:N \l__osglecture_pres_language_tl
\tl_new:N \l__osglecture_pres_language_override_tl
\tl_new:N \l__osglecture_pres_join_tl
\tl_new:N \l__osglecture_pres_finite_tl
\tl_new:N \l__osglecture_pres_relation_tl
\tl_new:N \l__osglecture_pres_initial_tl
\prop_new:N \g__osglecture_pres_language_alias_prop
\prop_new:N \g__osglecture_pres_relation_prop
\prop_new:N \g__osglecture_pres_initial_command_prop
\cs_generate_variant:Nn \prop_get:NnNTF {NVNTF}

\msg_new:nnnn { osglecture-presitemize } { unsupported-language }
  { Language~'#1'~has~no~presitemize~profile. }
  { Falling~back~to~English.~Use~language=de,~en,~or~fr. }
\msg_new:nnnn { osglecture-presitemize } { initial-not-letter }
  { Cannot~change~the~case~of~the~first~token~'#1'. }
  {
    Use~initial=keep,~\token_to_str:N \preskeepcase,~or~
    \token_to_str:N \prescase.
  }
\msg_new:nnnn { osglecture-presitemize } { duplicate-finite }
  { More~than~one~finite~verb~was~marked~in~one~subordinate~item. }
  { Combine~the~verb~cluster~in~one~\token_to_str:N \finite\ argument. }

\cs_new_protected:Npn \__osglecture_pres_alias:nn #1#2
  { \prop_gput:Nnn \g__osglecture_pres_language_alias_prop {#1} {#2} }

\clist_map_inline:nn {de,de-de,de-at,de-ch,german,ngerman}
  { \__osglecture_pres_alias:nn {#1} {de} }
\clist_map_inline:nn
  {en,en-gb,en-us,en-au,english,british,american,australian}
  { \__osglecture_pres_alias:nn {#1} {en} }
\clist_map_inline:nn {fr,fr-fr,fr-ca,french,francais}
  { \__osglecture_pres_alias:nn {#1} {fr} }

\cs_new_protected:Npn \__osglecture_pres_declare_relation:nnnn #1#2#3#4
  {
    \prop_gput:Nnn \g__osglecture_pres_relation_prop {#1/#2/join} {#3}
    \prop_gput:Nnn \g__osglecture_pres_relation_prop {#1/#2/finite} {#4}
  }

\__osglecture_pres_declare_relation:nnnn {de}{sentence}
  {\unskip.\c_space_tl}{inplace}
\__osglecture_pres_declare_relation:nnnn {de}{and}
  {\unskip\c_space_tl und\c_space_tl}{inplace}
\__osglecture_pres_declare_relation:nnnn {de}{but}
  {\unskip,\c_space_tl aber\c_space_tl}{inplace}
\__osglecture_pres_declare_relation:nnnn {de}{therefore}
  {\unskip,\c_space_tl weshalb\c_space_tl}{final}

\__osglecture_pres_declare_relation:nnnn {en}{sentence}
  {\unskip.\c_space_tl}{inplace}
\__osglecture_pres_declare_relation:nnnn {en}{and}
  {\unskip\c_space_tl and\c_space_tl}{inplace}
\__osglecture_pres_declare_relation:nnnn {en}{but}
  {\unskip,\c_space_tl but\c_space_tl}{inplace}
\__osglecture_pres_declare_relation:nnnn {en}{therefore}
  {\unskip,\c_space_tl so\c_space_tl}{inplace}

\__osglecture_pres_declare_relation:nnnn {fr}{sentence}
  {\unskip.\c_space_tl}{inplace}
\__osglecture_pres_declare_relation:nnnn {fr}{and}
  {\unskip\c_space_tl et\c_space_tl}{inplace}
\__osglecture_pres_declare_relation:nnnn {fr}{but}
  {\unskip,\c_space_tl mais\c_space_tl}{inplace}
\__osglecture_pres_declare_relation:nnnn {fr}{therefore}
  {\unskip,\c_space_tl donc\c_space_tl}{inplace}

\keys_define:nn { osglecture/presitemize }
  {
    language .tl_set:N = \l__osglecture_pres_language_override_tl,
    punctuation .choice:,
    punctuation / auto .code:n =
      { \bool_set_true:N \l__osglecture_pres_auto_punctuation_bool },
    punctuation / manual .code:n =
      { \bool_set_false:N \l__osglecture_pres_auto_punctuation_bool },
    punctuation .initial:n = auto,
  }

\keys_define:nn { osglecture/presitem }
  {
    relation .choices:nn = {sentence,and,but,therefore}
      { \tl_set:NV \l__osglecture_pres_relation_tl \l_keys_choice_tl },
    relation .initial:n = sentence,
    initial .choices:nn = {auto,keep}
      { \tl_set:NV \l__osglecture_pres_initial_tl \l_keys_choice_tl },
    initial .initial:n = auto,
  }

\cs_new_protected:Npn \__osglecture_pres_normalize_language:
  {
    \tl_set:Ne \l__osglecture_pres_language_tl
      { \text_lowercase:n {\tl_use:N \l__osglecture_pres_language_tl} }
    \tl_trim_spaces:N \l__osglecture_pres_language_tl
    \regex_replace_all:nnN { _ } {-} \l__osglecture_pres_language_tl
    \prop_get:NVNTF \g__osglecture_pres_language_alias_prop
      \l__osglecture_pres_language_tl \l_tmpa_tl
      { \tl_set_eq:NN \l__osglecture_pres_language_tl \l_tmpa_tl }
      {
        \msg_warning:nnV {osglecture-presitemize}{unsupported-language}
          \l__osglecture_pres_language_tl
        \tl_set:Nn \l__osglecture_pres_language_tl {en}
      }
  }

\cs_new_protected:Npn \__osglecture_pres_resolve_language:
  {
    \tl_if_empty:NTF \l__osglecture_pres_language_override_tl
      {
        \cs_if_exist:cTF {ver@babel.sty}
          {
            \cs_if_exist:NTF \localename
              { \tl_set:Ne \l__osglecture_pres_language_tl {\localename} }
              { \tl_set:Ne \l__osglecture_pres_language_tl {\languagename} }
          }
          {
            \cs_if_exist:cTF {ver@polyglossia.sty}
              {
                \tl_set:Ne \l__osglecture_pres_language_tl {\languagename}
              }
              {
                \cs_if_exist:NTF \olsTargetLanguage
                  { \tl_set:Ne \l__osglecture_pres_language_tl
                      {\olsTargetLanguage} }
                  {
                    \cs_if_exist:NTF \OsgLectureLanguage
                      { \tl_set:Ne \l__osglecture_pres_language_tl
                          {\OsgLectureLanguage} }
                      { \tl_set:Nn \l__osglecture_pres_language_tl {en} }
                  }
              }
          }
      }
      {
        \tl_set_eq:NN \l__osglecture_pres_language_tl
          \l__osglecture_pres_language_override_tl
      }
    \tl_if_empty:NT \l__osglecture_pres_language_tl
      { \tl_set:Nn \l__osglecture_pres_language_tl {en} }
    \__osglecture_pres_normalize_language:
  }

\NewDocumentCommand \longform {m}
  { \bool_if:NT \l__osglecture_pres_longform_bool {#1} }
\NewDocumentCommand \slides {m}
  { \bool_if:NF \l__osglecture_pres_longform_bool {#1} }
\NewDocumentCommand \longslides {mm}
  { \bool_if:NTF \l__osglecture_pres_longform_bool {#1} {#2} }
\NewDocumentCommand \prescase {mm}
  { \bool_if:NTF \l__osglecture_pres_longform_bool {#2} {#1} }
\NewDocumentCommand \preskeepcase {m} {#1}

% \ldeen{Diese beiden Deklarationen machen Befehle für die automatische
% Anpassung der Schreibung am Itemanfang transparent. Die erste Form
% beschreibt einen Befehl mit einem Textargument, die zweite ein
% Sprachwahlmakro mit zwei oder drei Textzweigen. Registrierte Hüllen werden
% rekursiv durchlaufen; in der Präsentation bleiben sie unverändert.}{These
% two declarations make commands transparent to automatic case adjustment at the
% beginning of joined items. The first form describes a command with one text
% argument, the second a language-selection command with two or three text
% branches. Registered wrappers are traversed recursively and remain unchanged
% in presentations.}
\NewDocumentCommand \DeclarePresInitialTransparentCommand { m }
  {
    \prop_gput:Nee \g__osglecture_pres_initial_command_prop
      { \cs_to_str:N #1 } { transparent }
  }
\NewDocumentCommand \DeclarePresInitialBranchingCommand { m m }
  {
    \int_case:nnF {#2}
      {
        {2}
          {
            \prop_gput:Nee \g__osglecture_pres_initial_command_prop
              { \cs_to_str:N #1 } { branches/2 }
          }
        {3}
          {
            \prop_gput:Nee \g__osglecture_pres_initial_command_prop
              { \cs_to_str:N #1 } { branches/3 }
          }
      }
      {
        \msg_error:nnn { osglecture-presitemize }
          { invalid-initial-branch-count } {#2}
      }
  }

\msg_new:nnnn { osglecture-presitemize } { invalid-initial-branch-count }
  { An~initial-transparent~branching~command~needs~two~or~three~branches,~not~'#1'. }
  { Use~2~for~a~bilingual~and~3~for~a~trilingual~selection~command. }

\clist_map_inline:nn
  { \emph,\textbf,\textit,\textsl,\textrm,\textsf,\texttt,\textnormal }
  { \DeclarePresInitialTransparentCommand #1 }

\msg_new:nnn {osglecture-presitemize} {outside-environment}
  {
    Command~'#1'~may~only~be~used~inside~a~presitemize~environment.
  }

\NewDocumentCommand \presitem {O{} d<> o}
  {
    \msg_error:nnn {osglecture-presitemize} {outside-environment}
      {\token_to_str:N \presitem}
  }
\NewDocumentCommand \anditem {d<> o}
  {
    \msg_error:nnn {osglecture-presitemize} {outside-environment}
      {\token_to_str:N \anditem}
  }
\NewDocumentCommand \butitem {d<> o}
  {
    \msg_error:nnn {osglecture-presitemize} {outside-environment}
      {\token_to_str:N \butitem}
  }
\NewDocumentCommand \thereforeitem {d<> o}
  {
    \msg_error:nnn {osglecture-presitemize} {outside-environment}
      {\token_to_str:N \thereforeitem}
  }

\NewDocumentCommand \finite {m}
  {
    \bool_if:nTF
      {
        \l__osglecture_pres_longform_bool
        && \l__osglecture_pres_subordinate_bool
      }
      {
        \tl_if_empty:NF \l__osglecture_pres_finite_tl
          { \msg_warning:nn {osglecture-presitemize}{duplicate-finite} }
        \tl_set:Nn \l__osglecture_pres_finite_tl {#1}
      }
      {#1}
  }

\cs_new_protected:Npn \__osglecture_pres_flush_finite:
  {
    \tl_if_empty:NF \l__osglecture_pres_finite_tl
      {
        \unskip
        \c_space_tl
        \tl_use:N \l__osglecture_pres_finite_tl
        \tl_clear:N \l__osglecture_pres_finite_tl
      }
  }

\cs_new_protected:Npn \__osglecture_pres_case_transparent:Nn #1#2
  { #1 { \__osglecture_pres_change_initial:n #2 } }

\NewDocumentCommand \__osglecture_pres_lower_two_branches: { m s +m +m }
  {
    \IfBooleanTF{#2}
      {
        #1*
          { \__osglecture_pres_change_initial:n #3 }
          { \__osglecture_pres_change_initial:n #4 }
      }
      {
        #1
          { \__osglecture_pres_change_initial:n #3 }
          { \__osglecture_pres_change_initial:n #4 }
      }
  }

\NewDocumentCommand \__osglecture_pres_lower_three_branches: { m s +m +m +m }
  {
    \IfBooleanTF{#2}
      {
        #1*
          { \__osglecture_pres_change_initial:n #3 }
          { \__osglecture_pres_change_initial:n #4 }
          { \__osglecture_pres_change_initial:n #5 }
      }
      {
        #1
          { \__osglecture_pres_change_initial:n #3 }
          { \__osglecture_pres_change_initial:n #4 }
          { \__osglecture_pres_change_initial:n #5 }
      }
  }

\cs_new_protected:Npn \__osglecture_pres_change_initial:n #1
  {
    \token_if_eq_meaning:NNTF #1 \preskeepcase
      { \preskeepcase }
      {
        \token_if_eq_meaning:NNTF #1 \prescase
          { \prescase }
          {
            \prop_get:NeNTF \g__osglecture_pres_initial_command_prop
              { \cs_to_str:N #1 } \l_tmpa_tl
              {
                \str_if_eq:VnTF \l_tmpa_tl { transparent }
                  { \__osglecture_pres_case_transparent:Nn #1 }
                  {
                    \str_if_eq:VnTF \l_tmpa_tl { branches/2 }
                      { \__osglecture_pres_lower_two_branches: #1 }
                      { \__osglecture_pres_lower_three_branches: #1 }
                  }
              }
              {
                \token_if_letter:NTF #1
                  { \text_uppercase:n {#1} }
                  {
                    \msg_warning:nnn {osglecture-presitemize}
                      {initial-not-letter}{#1}
                    #1
                  }
              }
          }
      }
  }

\cs_new_protected:Npn \__osglecture_pres_upper_initial:n #1
  { \__osglecture_pres_change_initial:n #1 }

\cs_new_protected:Npn \__osglecture_pres_set_relation:n #1
  {
    \prop_get:NeN \g__osglecture_pres_relation_prop
      {\l__osglecture_pres_language_tl/#1/join}
      \l__osglecture_pres_join_tl
    \prop_get:NeN \g__osglecture_pres_relation_prop
      {\l__osglecture_pres_language_tl/#1/finite}
      \l_tmpa_tl
    \str_if_eq:VnTF \l_tmpa_tl {final}
      { \bool_set_true:N \l__osglecture_pres_subordinate_bool }
      { \bool_set_false:N \l__osglecture_pres_subordinate_bool }
  }

\cs_new_protected:Npn \__osglecture_pres_original_item:nn #1#2
  {
    \IfNoValueTF{#2}
      { \__osglecture_pres_original_item: }
      { \__osglecture_pres_original_item: [#2] }
  }

\cs_new_protected:Npn \__osglecture_pres_item:nnnn #1#2#3#4
  {
    \__osglecture_pres_flush_finite:
    \keys_set:nn {osglecture/presitem}
      {relation=sentence,initial=auto,#1}
    \bool_if:NTF \l__osglecture_pres_auto_punctuation_bool
      { \exp_args:NV \__osglecture_pres_set_relation:n
          \l__osglecture_pres_relation_tl }
      { \tl_set:Nn \l__osglecture_pres_join_tl {\c_space_tl} }
    \__osglecture_pres_original_item:nn {#2}{#3}
    \bool_set_true:N \l__osglecture_pres_has_item_bool
    \str_if_eq:VnTF \l__osglecture_pres_initial_tl {keep}
      {#4}
      {
        \str_if_eq:VnTF \l__osglecture_pres_relation_tl {sentence}
          { \__osglecture_pres_upper_initial:n #4 }
          {#4}
      }
  }

\cs_new_protected:Npn \__osglecture_pres_compact_item:nnnn #1#2#3#4
  { \__osglecture_pres_item:nnnn {relation=#1}{#2}{#3}{#4} }

\IfLectureModeTF{longform}
  {
    \bool_set_true:N \l__osglecture_pres_longform_bool

    \NewDocumentEnvironment{presitemize}{O{}}
      {
        \group_begin:
        \tl_clear:N \l__osglecture_pres_language_override_tl
        \bool_set_true:N \l__osglecture_pres_auto_punctuation_bool
        \keys_set:nn {osglecture/presitemize}{#1}
        \__osglecture_pres_resolve_language:
        \tl_clear:N \l__osglecture_pres_finite_tl
        \tl_clear:N \l__osglecture_pres_join_tl
        \bool_set_false:N \l__osglecture_pres_subordinate_bool
        \bool_set_false:N \l__osglecture_pres_has_item_bool
        % \ldeen{Die Langform ist semantischer Fließtext, keine Inline-Liste.
        % Der Verzicht auf eine Listenumgebung erlaubt es der \LaTeX-Tagging-
        % Schicht außerdem, das Ergebnis als gewöhnlichen Absatzinhalt zu
        % taggen.}{Longform is semantic prose, not an inline list. Keeping
        % it out of a list environment also lets the \LaTeX\ tagging layer
        % tag the result as ordinary paragraph content.}
        \cs_set_protected:Npn \__osglecture_pres_original_item:
          {
            \bool_if:NT \l__osglecture_pres_has_item_bool
              { \tl_use:N \l__osglecture_pres_join_tl }
          }
        \RenewDocumentCommand \item {d<> o}
          {
            \__osglecture_pres_item:nnnn
              {relation=sentence}{##1}{##2}{}
          }
        \RenewDocumentCommand \presitem {O{} d<> o}
          { \__osglecture_pres_item:nnnn {##1}{##2}{##3} }
        \RenewDocumentCommand \anditem {d<> o}
          { \__osglecture_pres_compact_item:nnnn {and}{##1}{##2} }
        \RenewDocumentCommand \butitem {d<> o}
          { \__osglecture_pres_compact_item:nnnn {but}{##1}{##2} }
        \RenewDocumentCommand \thereforeitem {d<> o}
          { \__osglecture_pres_compact_item:nnnn {therefore}{##1}{##2} }
      }
      {
        \__osglecture_pres_flush_finite:
        \bool_if:nT
          {
            \l__osglecture_pres_has_item_bool
            && \l__osglecture_pres_auto_punctuation_bool
          }
          {\unskip.}
        \group_end:
      }
  }
  {
    \bool_set_false:N \l__osglecture_pres_longform_bool
    \cs_new_protected:Npn \__osglecture_pres_presentation_item:nnn #1#2#3
      {
        \keys_set:nn {osglecture/presitem}
          {relation=sentence,initial=auto,#1}
        \IfNoValueTF{#2}
          {
            \IfNoValueTF{#3}
              {\__osglecture_pres_original_item:}
              {\__osglecture_pres_original_item:[#3]}
          }
          {
            \IfNoValueTF{#3}
              {\__osglecture_pres_original_item:<#2>}
              {\__osglecture_pres_original_item:<#2>[#3]}
          }
        \str_if_eq:VnF \l__osglecture_pres_initial_tl {keep}
          { \__osglecture_pres_upper_initial:n }
      }
    \NewDocumentEnvironment{presitemize}{O{}}
      {
        \group_begin:
        \tl_clear:N \l__osglecture_pres_language_override_tl
        \bool_set_true:N \l__osglecture_pres_auto_punctuation_bool
        \keys_set:nn {osglecture/presitemize}{#1}
        \__osglecture_pres_resolve_language:
        \begin{itemize}
        \cs_set_eq:NN \__osglecture_pres_original_item: \item
        \RenewDocumentCommand \item {d<> o}
          { \__osglecture_pres_presentation_item:nnn {} {##1}{##2} }
        \RenewDocumentCommand \presitem {O{} d<> o}
          { \__osglecture_pres_presentation_item:nnn {##1}{##2}{##3} }
        \RenewDocumentCommand \anditem {d<> o}
          { \__osglecture_pres_presentation_item:nnn {} {##1}{##2} }
        \RenewDocumentCommand \butitem {d<> o}
          { \__osglecture_pres_presentation_item:nnn {} {##1}{##2} }
        \RenewDocumentCommand \thereforeitem {d<> o}
          { \__osglecture_pres_presentation_item:nnn {} {##1}{##2} }
      }
      {
        \end{itemize}
        \group_end:
      }
  }

\ExplSyntaxOff
%</presitemize>
%<*twocolumns>

\RequirePackage{osglecture-modes}

\ExplSyntaxOn

\cs_generate_variant:Nn \tl_if_blank:nF { V }
\cs_generate_variant:Nn \tl_if_eq:nnTF { Vn }
\cs_generate_variant:Nn \seq_put_right:Nn { NV }

\dim_new:N \l_twocol_sep_dim
\dim_new:N \l_twocol_avail_dim
\dim_new:N \l_twocol_a_dim
\dim_new:N \l_twocol_b_dim
\dim_new:N \l_twocol_h_dim

\fp_new:N  \l_twocol_frac_fp
\tl_new:N  \l_twocol_a_align_tl
\tl_new:N  \l_twocol_b_align_tl
\seq_new:N \l_twocol_align_seq
\seq_new:N \l_twocol_body_seq

\box_new:N \l_twocol_a_box
\box_new:N \l_twocol_b_box

\dim_set:Nn \l_twocol_sep_dim { \columnsep }

\cs_new_protected:Npn \twocol_add_option:n #1
  {
    \tl_set:Nx \l_tmpa_tl { \tl_trim_spaces:n {#1} }

    \tl_if_blank:VF \l_tmpa_tl
      {
        \tl_if_eq:VnTF \l_tmpa_tl { t }
          { \seq_put_right:NV \l_twocol_align_seq \l_tmpa_tl }
          {
            \tl_if_eq:VnTF \l_tmpa_tl { c }
              { \seq_put_right:NV \l_twocol_align_seq \l_tmpa_tl }
              {
                \tl_if_eq:VnTF \l_tmpa_tl { b }
                  { \seq_put_right:NV \l_twocol_align_seq \l_tmpa_tl }
                  { \fp_set:Nn \l_twocol_frac_fp { \l_tmpa_tl } }
              }
          }
      }
  }

\cs_new_protected:Npn \twocol_parse:n #1
  {
    \fp_set:Nn \l_twocol_frac_fp { .5 }
    \tl_set:Nn \l_twocol_a_align_tl { c }
    \tl_set:Nn \l_twocol_b_align_tl { c }
    \seq_clear:N \l_twocol_align_seq

    \clist_map_inline:nn {#1}
      { \twocol_add_option:n {##1} }

    \fp_compare:nF
      { 0 < \l_twocol_frac_fp && \l_twocol_frac_fp < 1 }
      { \fp_set:Nn \l_twocol_frac_fp { .5 } }

    \int_case:nn { \seq_count:N \l_twocol_align_seq }
      {
        {1}
          {
            \tl_set:Nx \l_twocol_a_align_tl
              { \seq_item:Nn \l_twocol_align_seq {1} }
            \tl_set_eq:NN \l_twocol_b_align_tl \l_twocol_a_align_tl
          }
        {2}
          {
            \tl_set:Nx \l_twocol_a_align_tl
              { \seq_item:Nn \l_twocol_align_seq {1} }
            \tl_set:Nx \l_twocol_b_align_tl
              { \seq_item:Nn \l_twocol_align_seq {2} }
          }
      }

    \dim_set:Nn \l_twocol_avail_dim
      { \linewidth - \l_twocol_sep_dim }

    \dim_set:Nn \l_twocol_a_dim
      {
        \fp_to_dim:n
          { \l_twocol_frac_fp * \dim_to_fp:n { \l_twocol_avail_dim } }
      }

    \dim_set:Nn \l_twocol_b_dim
      { \l_twocol_avail_dim - \l_twocol_a_dim }
  }

\cs_new_protected:Npn \twocol_set_box:NNn #1#2#3
  {
    \vbox_set:Nn #1
      {
        \hsize=#2
        \dim_set:Nn \linewidth { \hsize }
        \dim_set:Nn \columnwidth { \hsize }
        #3
        \par
      }
  }

\cs_new_protected:Npn \twocol_aligned_box:NN #1#2
  {
    \vbox to \l_twocol_h_dim
      {
        \tl_if_eq:VnTF #2 { b }
          { \vfil \box_use:N #1 }
          {
            \tl_if_eq:VnTF #2 { c }
              { \vfil \box_use:N #1 \vfil }
              { \box_use:N #1 \vfil }
          }
      }
  }

% \ldeen*{Die beiden Spaltenboxen werden nebeneinander als einzelne
% \code{hbox}es gesetzt, nicht mit \LaTeX's echten \env{multicols}. Damit
% der Struct-Tree dennoch einen zusammenhängenden Block statt zweier
% unabhängiger Fragmente sieht, wird die gesamte Konstruktion vorab als
% \code{Div}-Struktur getaggt (nur wenn @1 verfügbar ist, also unter
% \cs{DocumentMetadata}).}{The two column boxes are typeset side by
% side as plain \code{hbox}es, not with \LaTeX's real \env{multicols}. So
% the struct tree still sees one coherent block instead of two
% independent fragments, the whole construction is tagged beforehand as
% a \code{Div} structure (only when @1 is available, i.e.\ under
% \cs{DocumentMetadata}).}{\cs{tagstructbegin}}
\cs_new_protected:Npn \twocol_tag_begin:
  {
    \cs_if_exist:NT \tagstructbegin
      { \tagstructbegin { tag=Div } }
  }

\cs_new_protected:Npn \twocol_tag_end:
  {
    \cs_if_exist:NT \tagstructend
      { \tagstructend }
  }

\cs_new_protected:Npn \twocol_typeset_columns:nn #1#2
  {
    \twocol_parse:n {#1}
    \seq_set_split:Nnn \l_twocol_body_seq { \nextcolumn } {#2}

    \par
    \twocol_tag_begin:

    \twocol_set_box:NNn \l_twocol_a_box \l_twocol_a_dim
      { \seq_item:Nn \l_twocol_body_seq {1} }

    \twocol_set_box:NNn \l_twocol_b_box \l_twocol_b_dim
      { \seq_item:Nn \l_twocol_body_seq {2} }

    \dim_set:Nn \l_twocol_h_dim
      {
        \dim_max:nn
          { \box_ht:N \l_twocol_a_box + \box_dp:N \l_twocol_a_box }
          { \box_ht:N \l_twocol_b_box + \box_dp:N \l_twocol_b_box }
      }

    \hbox to \linewidth
      {
        \hbox to \l_twocol_a_dim
          { \twocol_aligned_box:NN \l_twocol_a_box \l_twocol_a_align_tl \hss }
        \hskip \l_twocol_sep_dim
        \hbox to \l_twocol_b_dim
          { \twocol_aligned_box:NN \l_twocol_b_box \l_twocol_b_align_tl \hss }
      }
    \twocol_tag_end:
    \par
  }

\cs_new_protected:Npn \twocol_typeset_plain:n #1
  {
    \seq_set_split:Nnn \l_twocol_body_seq { \nextcolumn } {#1}
    \seq_use:Nn \l_twocol_body_seq { \par\smallskip }
  }

\NewDocumentEnvironment{twocolumns}{ D<>{presentation} O{} +b }
  {
    \group_begin:
    \tl_if_blank:nTF {#1}
      { \tl_set:Nn \l_tmpa_tl {presentation} }
      { \tl_set:Nn \l_tmpa_tl {#1} }
    \exp_args:NV \IfLectureModeTF \l_tmpa_tl
      { \twocol_typeset_columns:nn {#2} {#3} }
      { \twocol_typeset_plain:n {#3} }
    \group_end:
  }
  {}

\ExplSyntaxOff
%</twocolumns>
%<*series>

\ExplSyntaxOn

\cs_new:Npn \OsgLectureSeriesAvailableFlag { 0 }
\cs_new:Npn \OsgLectureSeriesPosition { }
\cs_new:Npn \OsgLectureSeriesCurrentPhysicalUnit { }
\cs_new:Npn \OsgLectureSeriesCurrentRole { }
\cs_new:Npn \OsgLectureSeriesPreviousPhysicalUnit { }
\cs_new:Npn \OsgLectureSeriesPreviousContinuationPhysicalUnit { }
\cs_new:Npn \OsgLectureSeriesNextPhysicalUnit { }
\cs_new:Npn \OsgLectureUnitPhysicalUnit { }
\cs_new:Npn \OsgLectureUnitRole { }

\prop_new:N \g__osglecture_series_physical_to_logical_prop
\prop_new:N \g__osglecture_series_logical_to_physical_prop
\tl_new:N \l__osglecture_series_physical_tl
\tl_new:N \l__osglecture_series_logical_tl

\msg_new:nnnn { osglecture-series } { index-failed }
  { Cannot~initialize~the~series~index:~#1 }
  { Check~the~source~directory~and~the~target's~unit_scopes. }

\NewDocumentCommand \OsgLectureSeriesAvailableTF { m m }
  {
    \str_if_eq:eeTF { \OsgLectureSeriesAvailableFlag } { 1 } {#1} {#2}
  }

\msg_new:nnn { osglecture-series } { inconsistent-unit }
  { Inconsistent~physical/logical~series-unit~mapping:~#1. }

\keys_define:nn { osglecture/series-unit }
  {
    physical .tl_set:N = \l__osglecture_series_physical_tl,
    unit .tl_set:N = \l__osglecture_series_logical_tl,
  }

\NewDocumentCommand \OsgLectureSeriesUnit { m }
  {
    \tl_clear:N \l__osglecture_series_physical_tl
    \tl_clear:N \l__osglecture_series_logical_tl
    \keys_set:nn { osglecture/series-unit } {#1}
    \prop_get:NVNT \g__osglecture_series_physical_to_logical_prop
      \l__osglecture_series_physical_tl \l_tmpa_tl
      {
        \tl_if_eq:VVF \l_tmpa_tl \l__osglecture_series_logical_tl
          {
            \msg_error:nnx { osglecture-series } { inconsistent-unit }
              {
                physical~unit~'\l__osglecture_series_physical_tl'~maps~to~both~
                '\l_tmpa_tl'~and~'\l__osglecture_series_logical_tl'
              }
          }
      }
    \prop_get:NVNT \g__osglecture_series_logical_to_physical_prop
      \l__osglecture_series_logical_tl \l_tmpa_tl
      {
        \tl_if_eq:VVF \l_tmpa_tl \l__osglecture_series_physical_tl
          {
            \msg_error:nnx { osglecture-series } { inconsistent-unit }
              {
                logical~unit~'\l__osglecture_series_logical_tl'~maps~to~both~
                '\l_tmpa_tl'~and~'\l__osglecture_series_physical_tl'
              }
          }
      }
    \prop_gput:NVV \g__osglecture_series_physical_to_logical_prop
      \l__osglecture_series_physical_tl \l__osglecture_series_logical_tl
    \prop_gput:NVV \g__osglecture_series_logical_to_physical_prop
      \l__osglecture_series_logical_tl \l__osglecture_series_physical_tl
  }

\NewExpandableDocumentCommand \OsgLectureSeriesLogicalUnit { m }
  { \prop_item:Nn \g__osglecture_series_physical_to_logical_prop {#1} }
\NewExpandableDocumentCommand \OsgLectureSeriesPhysicalUnit { m }
  { \prop_item:Nn \g__osglecture_series_logical_to_physical_prop {#1} }
\NewDocumentCommand \OsgLectureSeriesUnitKnownTF { m m m }
  {
    \prop_if_in:NnTF \g__osglecture_series_logical_to_physical_prop
      {#1} {#2} {#3}
  }

% \ldeen*{Der Startpfad stammt implizit vom Arbeitsverzeichnis des
% TeX-Laufs (@1 wäre redundant zu \code{lfs.currentdir()}, siehe
% \code{ARCHITECTURE.md} Abschnitt 12). Der Vorabtest mit @2 sorgt
% dafür, dass die volle Initialisierung (samt Scope-Analyse und
% Fehlermeldung bei echten Problemen) nur versucht wird, wenn überhaupt
% ein Serieneinheit-artiges Vorfahrenverzeichnis existiert. Ohne ein
% solches Vorfahrenverzeichnis bleibt die Serie still nicht verfügbar,
% statt einen Fehler zu erzeugen.}{The start path comes implicitly from
% the TeX run's working directory (@1 would be redundant with
% \code{lfs.currentdir()}, see \code{ARCHITECTURE.md} section 12). The
% pre-check with @2 ensures the full initialization (including scope
% analysis and error reporting for real problems) is only attempted
% when a series-unit-shaped ancestor directory exists at all. Without
% such an ancestor, the series stays silently unavailable, instead of
% raising an error.}{\code{source-directory}}{\code{series\_index.locate\_unit\_ancestor}}
\cs_new_protected:Npn \__osglecture_series_initialize:nn #1#2
  {
    \directlua{
      require("osglecture-series-index").initialize_tex_csv_if_available(
        nil, "\luaescapestring{#1}", "\luaescapestring{#2}")
    }
  }

\OsgLectureBuildLoadedTF
  {
    \exp_args:Nee \__osglecture_series_initialize:nn
      { \OsgLectureBuildValue{doctype} }
      { \OsgLectureBuildValue{applicable-unit-scopes} }
  }
  { }

\ExplSyntaxOff
%</series>
%<*osgbeamer>
\ProvidesFile{osglecture-osgbeamer.code.tex}
  [2026/07/29 v0.6.0 legacy osgbeamer implementation]
% \ldeen{Legacy-Implementierung, unverändert übernommen; nicht Teil der
%   einheitlichen Kommentarkonvention, da zur baldigen Ablösung
%   vorgesehen.}{Legacy implementation, carried over unchanged; not
%   part of the unified comment convention, since it is slated for
%   removal soon.}
\def\classname{osglecture}
%%%%%%%%%%%%%%%%%%%%%%%%%%%%%%%%%%%%%%%%%%%%%%%%%%%%%%%%%%%%%%%%%%%%%%%%%%%%%%
%
% Beamer- und Script-Template für die Professur Betriebssysteme an
% der TU Chemnitz
%
% (c) Matthias Werner (matthias.werner@informatik.tu-chemnitz.de), 2022-2024
%
% Dieses Material steht unter der Creative-Commons-Lizenz CC BY-SA 4.0
% (Namensnennung - Weitergabe unter gleichen Bedingungen).
% Um eine Kopie dieser Lizenz zu sehen, besuchen Sie                                        
% http://creativecommons.org/licenses/by-sa/4.0/deed.de.                      
%
% Hinweis:
% Obwohl diese Klasse eigenständig genutzt werden kann, ist sie primär
% für den Einsatz in Vorlesungsreihen im Kombination mit dem
% OSG LaTeX Lecture Maker (ollm, einem rc-Script für latexmk) gedacht.
%
\IfFormatAtLeastTF{2022/06/01}{}{
\ClassError{\classname}{\MessageBreak
************************************************\MessageBreak
* This class requires at least the TeX format \MessageBreak
* from 2022/06. \MessageBreak
************************************************
 }{Update your LaTeX.}
}
\RequirePackage{iftex}               % für Test auf Lua(La)TeX
\ifLuaTeX\else
\ClassError{\classname}{\MessageBreak
************************************************\MessageBreak
* LuaLaTeX is required to use this class. \MessageBreak
* Sorry! \MessageBreak
************************************************
 }{Use this class with LuaLaTex.}
\fi
\RequirePackage{etoolbox,auxhook}
% Die nächsten beiden Zeilen dienen nur dafür, dass mein Editor die
% etoolbox-if nicht dauernd mit TeX-if verwechselt.
\let\IfBool=\ifbool 
\let\IfToggle=\iftoggle
% \if's verwechselt...

\RequirePackage{varsfromjobname}
\RequirePackage{luacode}
%%%%%%%%%%%%%%%%%%%%%%%%%%%%%%%%%%%%%%%%%%%%%%%%%%%%%%%%%%%%%%%%%%%%%%%%%%%%%%
% OSG LaTeX Lecture Maker
% =======================
% 
% Das OSG LaTeX lecture maker system (OLLM system) steht teilweise über den
% Klassenoptionen. Im OLLM System werden an die Klasse durch den Jobnamen und 
% durch vordefinierte Macros Informationen übergeben. Macros und Jobnamen werden
% durch das `latexmkrc` aufgearbeitet.
% Der Jobname ist wie folgt aufgebaut und wird ausgewertet, wenn das Macro
% `ollm` definiert ist.
%
% <lectureprefix>-<number>-<doctype>-<language>-<topic>
%
% * <lectureprefix>: wahlfreies Kürzel für die Vorlesungsreihe
% * <number>       : zweistellig Nummer des Kapitels
% * <doctype>      : Ausgabeformat (siehe unten)
% * <language>     : de oder en (für deutsch oder englisch)
% * <topic>        : wahlfreie Bezeichnung für das Kapitelthema
%
% Ob das OLLM System aktiv ist, wird anhand des Vorhandenseins des der Datei
% 'ollmconfig.pl' im Elternverzeichnis festgestellt und im Flag 'osgllm'
% gespeichert. Dann wird davon ausgegangen, dass ein korrektes
% Verzeichnis-Layout existiert.
\newbool{osgllm}\boolfalse{osgllm}
\newbool{osglangset}\boolfalse{osglangset}
\newbool{osgdoctypeset}\boolfalse{osgdoctypeset}
\RequirePackage{osglecture-config}
\IfFileExists{\jobname.osgbuild.tex}{
  \OsgLectureLoadBuildFileForJob
  % The legacy code below uses the presence of \ollm as its series marker.
  % Build-file values themselves are read through osglecture-config.
  \providecommand\ollm{}
}{}
\begin{luacode*}
  -- lfs wird für currentdir() benötigt
  require("lfs")
  -- Teste auf das Vorhandensein einer Konfigurationsdatei.
  -- Nur wenn die Datei existiert, lässt sie sich öffnen
  f = io.open("../ollmconfig.pl","r")
  if f ~= nil then
      io.close(f)
      -- Wenn Datei gefunden, wird osgllm-Flag gesetzt 
      tex.print("\\booltrue{osgllm}")
      -- Defaultwerte für OLLM
      sh_path = "../Ref"
      cfg_file = "lectdates"
      fst_chap_a = "1"
      fst_chap_b = "1"
      -- curr_dir hat keinen Defaultwert, wird direkt gesetzt
      cur_dir = lfs.currentdir() 
      -- Die anderen Werte werden aus der Konfigurationsdatei geholt.
      for line in (io.lines("../ollmconfig.pl")) do
        if string.sub(line,1,1) ~= "#" then
           -- Nichtkommentarzeile, versuche parsen
           v = string.match(line,"$first_chapter_number%s-=%s-(%S-);")
           if v~= nil then
              fst_chap_a = v
              fst_chap_b = v
           end
           v1,v2 = string.match(line, "@first_chapter_number%s-=%s-%((%S-)%s-,%s-(%S-)%)%s-;")
           if v1~= nil then
              fst_chap_b = v1
              fst_chap_a = v2
           end
           v = string.match(line,"$lectconfig%s-=%s-(%S-);")
           if v ~= nil then
              cfg_file = v
           end   
           v = string.match(line,"$shared_data_dir%s-=%s-(%S-);")
           if v ~= nil then
              sh_path = v
           end   
        end
        -- \ollm dient als Hook für verschiedene OLLM-spezifische Macros
        tex.print("\\gdef\\ollm{\\gdef\\OsgShareDataPath{".. sh_path.."}"..
          "\\gdef\\OsgLectConfig{"..cfg_file.."}"..
          "\\gdef\\OsgCurrentDir{"..cur_dir.."}"..
          "\\IfToggle{osgscript}{\\gdef\\OsgFirstChapter{"..fst_chap_a.."}}{\\gdef\\OsgFirstChapter{"..fst_chap_b.."}}"..
          "}")
      end
  else 
      tex.print("\\typeout{No proper directory configuration, assume standalone}")
  end    
\end{luacode*}
\ifdef{\ollm}{
   \ClassInfo{\classname}{osgllm set}
   \booltrue{osgllm}
 }{
   \ClassInfo{\classname}{osgllm unset}
  \boolfalse{osgllm}
}
\newtoggle{osgscript}         % Wird bis zum Laden von Beamer als Unterscheidung
                              % von presentation- und article-Mode genutzt

\newbool{osgoptionstandalone} % Zeigt die Nutzung der Klasse außerhalb einer
                              % Lecture-Serie an

\newbool{osglegacy}           % Zeigt die Nutzung des "Kompatibilitätsmodus" an
\boolfalse{osglegacy}         
\newbool{osgforcetoc}         % Legt eine Inhaltsverzeichnisdatei an, ohne ein 
\booltrue{osgforcetoc}        % entsprechendes Verzeichnis zu generieren
\newbool{osgbib}              % true lädt osgbeamerbib.sty
\booltrue{osgbib}

%%%%%%%%%%%%%%%%%%%%%%%%%%%%%%%%%%%%%%%%%%%%%%%%%%%%%%%%%%%%%%%%%%%%%%%%%%%%%%
% Klassenoptionen
%
% Durch das neue Key-Value-interface von LaTeX ab Juni 2022 können (wie in
% pgfkeys oder Koma-Script Klassenoptionen ggf. auch noch später gesetzt
% werden (so sie nicht vor der ersten Seitenserstellung gebraucht werden).
%%%%%%%%%%%%%%%%%%%%%%%%%%%%%%%%%%%%%%%%%%%%%%%%%%%%%%%%%%%%%%%%%%%%%%%%%%%%%%
\DeclareKeys[osgbeamer]{
  % The dispatcher in osglecture.cls has already consumed this option.
  osgbeamer.code = {},
  % Soll die Kompatibilität zur alten osgbeamer-Klasse erhöhen (z.Z. nicht genutzt)
  legacy.code={
    \ClassWarning{\classname}{Legacy mode activated}
    \booltrue{osglegacy}
  },
  % Optionen, die an die beamer-Klasse weitergeleitet werden.
  beamer.code = {
    \PassOptionsToClass{#1}{beamer}
  },
  % Da 'aspectratio' vermutlich häufiger als andere Optionen abweicht, erhält es
  % eine eigene Option
  aspectratio.code={
    \PassOptionsToClass{aspectratio=#1}{beamer}
  },
  % Optionen, die an die scrbook-Klasse weitergeleitet werden.
  book.code = {
    \PassOptionsToClass{#1}{scrbook}
  },
  % Optionen, die an das tucbeamer-Package weitergeleitet werden.
  tuc.code={
    \PassOptionsToPackage{#1}{beamerthemetuc2019}
    %\PassOptionsToPackage{#1}{beamerarticletuc2019}
  },
  % Bei nobib wird osgbeamerbib.sty nicht geladen. Die bib-Option wird
  % ignoriert. Die Integration von \cite etc. funktioniert nicht.
  % Dies ist sinnvoll, wenn
  % - keine Literaturreferenzen genutzt werden, oder
  % - eine eigene Verwaltung für Literaturreferenzen genutzt werden soll,
  %   z.B. biblatex
  nobib.code = {
    \boolfalse{osgbib}
  },
  % Optionen für Umgang mit Literaturreferenzen (siehe osgbeamerbib.sty)
  bib.code = {
    \PassOptionsToPackage{#1}{osgbeamerbib}
  },
  % Werden Seiten kapitelübergreifend numeriert?
  continuation.code={
    \PassOptionsToClass{continuation=#1}{osgbeamerref}
  },
  docid.code = {
    \PassOptionsToPackage{docid=#1}{osgbeamerref}
  },
  % Eine Standalone-Präsentation ist nicht Teil einer Serie
  standalone.code = {
    \booltrue{osgoptionstandalone}
    \PassOptionsToPackage{standalone}{osgbeamerref}
    \PassOptionsToPackage{standalone}{osgbeamerbib}
  },
  % Die ggf. durch OLLM gegeben Konfigurationen werden
  % ignoriert. Erzwingt den Standalone-Modus
  noollm.code = {
    \booltrue{osgoptionstandalone}
    \undef{\ollm}
    \boolfalse{osgllm}
    \ClassInfo{\classname}{OLLM disabled or not detected, deactivate OLLM features}
  },
  % Optionen, die sinnvolle Einstellungen für die OSG zusammenfassen
  osgdefaults.code = {
    \PassOptionsToPackage{fakcolor=if,latexfonts}{beamerthemetuc2019}
    \PassOptionsToPackage{fakcolor=if,latexfonts}{beamerarticletuc2019}
    \SetLogo{OSG-Logo}
    \PassOptionsToClass{professionalfonts}{beamer}
  },
  final.code{
    \PassOptionsToPackage{final}{osgbeamerref}
    \PassOptionsToClass{final}{beamer}
    \PassOptionsToPackage{final}{scrbook}
  }
  noforcetoc.code = {
    \boolfalse{osgforcetoc}
  }
}
% Da Choice-Keys benötigt werden, die noch nicht nach LaTeX2e
% zurückportiert wurden, wird LaTeX3 benötigt. Daher wird auf die
% alternative Syntax umgestellt. Ein explizites Laden eines
% LaTeX3-Pakets ist nicht nötig, da die aktuelle LaTeX-Version auf
% LaTeX3 beruht.
\ExplSyntaxOn
\cs_generate_variant:Nn \keys_set:nn { nx } % SetKey mit Expansion der Werteargumente
\let\SetKeyEx=\keys_set:nx                  % Alias mit LaTeX2e-Syntax
\keys_define:nn {osgbeamer}{
  % Doctype
  % =======
  %
  % Es können mehrere Arten von Ausgabeformaten (Doctypes) erzeugt werden:
  % * slides  => Foliensatz,
  % * handout => mehrere (Standard:4) Slides auf einem Blatt,
  % * screen  => Slides mit Kommentaren und Vorschau für die
  %              Vorführung (z.Z. nicht richtig unterstützt)
  % * script  => Kapitel eines Scripts
  % * web     => Scriptkapitel in HTML (z.Z. nicht unterstützt)
  %
  % In der Regel sollte diese Option unter OLLM nicht per Hand gesetzt werden.
  doctype .choice: ,
  doctype/slides .code:n={
    \booltrue{osgdoctypeset}
    \gdef\osgparentclass{beamer}
    \PassOptionsToPackage{doctype=slides}{osgbeamerref}
    \togglefalse{osgscript}
  },
  doctype/handout .code:n={
    \booltrue{osgdoctypeset}
    \PassOptionsToClass{handout}{beamer}
    \PassOptionsToPackage{doctype=handout}{osgbeamerref}
    \togglefalse{osgscript}
  },
  doctype/script .code:n={
    \booltrue{osgdoctypeset}
    \PassOptionsToPackage{doctype=script}{osgbeamerref}
    \toggletrue{osgscript}
  },
  doctype/article .code:n={
    \booltrue{osgdoctypeset}
    \PassOptionsToPackage{doctype=script}{osgbeamerref}
    \toggletrue{osgscript}
  },
  doctype/screen .code:n={
    \booltrue{osgdoctypeset}
    \PassOptionsToClass{second}{beamer}
    \PassOptionsToPackage{doctype=slides}{osgbeamerref}
    \togglefalse{osgscript}
  },
  doctype/web .code:n={
    % \PassOptionsToClass{second}{beamer}
    \booltrue{osgdoctypeset}
    \ClassError{\classname}{Document type 'web' is not supported.}
    %\toggletrue{osgscript}
    %\toggletrue{osgweb}
  },
  % Die Sprache _kann_ explizit mit der Option lang=<lang> gesetzt werden, jedoch
  % sollte dies primär im 'standalone'-Mode genutzt werden
  % Im OLLM wird die entsprechende Sprache i.d.R. über den Jobnamen
  % gesetzt (siehe osgbeamerlanguage.sty)
  lang .choice:,
  lang/de .code:n = {
    \booltrue{osglangset}
    \providecommand\selectedbeamerlanguage{de}
   \PassOptionsToPackage{de}{osgbeamerlanguage}
   \PassOptionsToPackage{ngerman}{translator}
  },
  lang/en .code:n = {
    \booltrue{osglangset}
    \providecommand\selectedbeamerlanguage{en}
    \PassOptionsToPackage{en}{osgbeamerlanguage}
    \PassOptionsToPackage{english}{translator}
  },
  
  % Mit der Option 'handout format' kann angegeben werden, wieviel Folien
  % auf einer A4 Seite dargestellt werden.
  % Es gibt die Auswahl zwischen '2 on 1', '4 on 1' und '6 on 1'
  handout~format.choice:,
  handout~format/2~on~1 .code ={
    \gdef\osg@pgfpageoptions{{2~on~1}[a4paper, border~shrink=2mm]}
  },
  handout~format/4~on~1 .code ={
    \gdef\osg@pgfpageoptions{{4~on~1}[a4paper, landscape, border~shrink=2mm]}
  },
  handout~format/6~on~1 .code ={
    \gdef\osg@pgfpageoptions{{6~on~1}[a4paper, border~shrink=2mm]}
  },
}
\ExplSyntaxOff

% In \ollm werden die Pfade zum Verzeichnis für gemeinsame Daten (für
% vorlesungsübergreifende Referenzierung), der eigene Verzeichnisname
% sowie die erste logische Kapitelnummer (i.d.R. 0 oder 1) übergeben.
\IfBool{osgllm}{
  \ClassInfo{\classname}{Acivated OLLM}  
}{
  \boolfalse{osgllm}
  % Ohne OLLM ist das Dokument alleinstehend (standalone), auch
  % wenn dies nicht explizit als Option angegeben wurde.
  \ClassInfo{\classname}{Osgbeamer without OLLM, activate standalone}  
  \SetKeys[osgbeamer]{standalone=true}
}
%%%%%%%%%%%%%%%%%%%%%%%%%%%%%%%%%%%%%%%%%%%%%%%%%%%%%%%%%%%%%%%%%%%%%%%%%%%%%%
% Default-Einstellungen:
%
% Hier werden (aus Sicht des Klassenautoren) sinnvolle
% Voreinstellungen getroffen.
%
\SetKeys[osgbeamer]{
  aspectratio=1610,
  tuc={
    fakcolor=if,
  },
  beamer={
    ignorenonframetext,
    10pt,
    unknownkeysallowed %,notheorems
  },
  book={
    chapterprefix=true,
    twoside,
    open=right,
    cleardoublepage=current,
    parskip=half,
    DIV=9,
    fontsize=10pt
  },
  handout format=4 on 1,
}
\SetKeyEx{osgbeamer}{
  bib={file = \getonefromjobname,
    off}
}

%%%%%%%%%%%%%%%%%%%%%%%%%%%%%%%%%%%%%%%%%%%%%%%%%%%%%%%%%%%%%%%%%%%%%%%%%%%%%%
% Es gibt drei Möglichkeiten, wie spezifische Klassenoptionen gesetzt
% werden können:
%   1. Mit \SetGlobalClassOptions in lectdate.tex. Die hier gesetzten
%      Optionen gelten für *alle* Kapitel.
%   2. Als 'normale' Klassenoption. Die so gesetzten Optionen gelten
%      nur in dem Kapitel und überschreiben die Optionen von 1.
%   3. Mit \EnforceGlobalClassOptions in lectdate.tex. Die hier gesetzten
%      Optionen gelten für *alle* Kapitel und überschreiben Optionen
%      von 1. und 2.
%
% Frühe Optionen
% ==============
% 
% Um Klassen-Optionen für *alle* Kapitel/Lectures gemeinsam setzen zu können,
% soll die 'lectdate.tex' bereits jetzt eingelesen werden.
% Dort werden die Macros \SetGlobalClassOptions und
% \EnforceGlobalClassOptions zur Verfügung gestellt.
\NewDocumentCommand{\SetGlobalClassOptions}{m}
  {\SetKeys[osgbeamer]{#1}}
\NewDocumentCommand{\EnforceGlobalClassOptions}{m}{
  \gappto{\osg@forceoptionhook}{\SetKeys[osgbeamer]{#1}}}

% Allerdings sind zu diesem Zeitpunkt die meisten Titelbefehle noch
% nicht bekannt oder noch nicht von Beamer umdefiniert.
% Daher werden sie hier temporär definiert und ihre eigentliche
% Ausführung auf später verschoben.
\gdef\osg@titlematterhook{}            % Hook für Titelbefehle
\gdef\osg@forceoptionhook{}            % Hook für Options

% Bereite die Befehle auf ihre spätere Ausführung vor
% 
% Anmerkung: Praktisch alle Titelmacros werden durch Beamer, Koma oder OSG-Beamer
% überschrieben. Da jedoch (vermutlich) häufig auf die Kernel-Definition zurückgegriffen
% wird, ist es sinnvoll, diese zu erhalten.
% 
% In der Stern-Variante definiert \osg@makedelayed das Ersatzmacro mit einem optionalen
% Argument, so wie es später auch durch Beamer und co. definiert wird.
% 
% Alle Titlemacros sind gewahr (wie übersetzt man 'aware'?) der Beamer-Modi, es lassen
% sich also mode-spezifische Titel (z.B. \title<article>{Foo}) spezifizieren. Dies wird
% *nicht* im Macro selbst verankert, sondern über eine \mode-Umgebung zur verzögerten
% Ausführung realisiert; entsprechend stehen die Modes außerhalb von 'lectdates.tex' nicht
% zur Verfügung (auch nicht in von 'lectdates.tex' includierten Dateien).
\NewDocumentCommand{\osg@makedelayed}{s m}{
  % Wenn das Macro schon (durch den Kernel) definiert ist, speichere es unter einem anderen
  % Namen, sonst setzte es auf \relax.
  \ifcsdef{#2}{\ExpandArgs{cc}\NewCommandCopy{original@#2}{#2}}{\cslet{original@#2}{\relax}}
  \IfBooleanTF{#1}{ % Stern => mit optionalem Argument
     \ExpandArgs{c}\DeclareDocumentCommand{#2}{d<> O{} m}{ % Definiere Ersatzmacro
      \IfNoValueTF{##1}{
        \gappto{\osg@titlematterhook}{\csuse{#2}[##2]{##3}}
      }{
        \gappto{\osg@titlematterhook}{\mode<##1>{\csuse{#2}[##2]{##3}}}
      }
    }
  }{ % kein Stern => ohne optionales Argument
    \ExpandArgs{c}\DeclareDocumentCommand{#2}{d<> m}{      
      \IfNoValueTF{##1}{
        \gappto{\osg@titlematterhook}{\csuse{#2}{##2}}
      }{
        \gappto{\osg@titlematterhook}{\mode<##1>{\csuse{#2}{##2}}}
      }
    }
  }
}
% Setze Macro auf die alte Definition zurück oder (wenn vorher nicht definiert) lösche es.
\NewDocumentCommand{\osg@restoredelayed}{m}{
  \ifcsvoid{original@#1}{\csundef{#1}}{\ExpandArgs{cc}\DeclareCommandCopy{#1}{original@#1}}
}
% Da \logo schon anderweitig von Beamer belegt, von der Logik aber in eine Reihe
% wie \author,\date etc. passt, wird für lectdates.tex \logo auf \SetLogo umgelenkt
\NewDocumentCommand{\SetLogo}{O{#2} m}{
  \gdef\osg@logo{#2}
  \gdef\osg@tinylogo{#1}
}
\let\logo=\SetLogo

\IfBool{osgllm}{ % Zu diesem Zeitpunkt sind die Klassenoptionen noch nicht
  % ausgewertet. lectdates.tex wird also auch geladen, wenn OLLM korrekte
  % Verzeichnisstrukturen vorfindet und der Nutzer nicht 'standalone' 
  % als OLLM-Option übergibt, sondern nur als Klassenoption
  \OsgLectureBuildLoadedTF{}{
    \ollm % Definiert \OsgShareDataPath, \OsgLectConfig und \OsgFirstChapter
    \expandafter\InputIfFileExists{\OsgLectConfig}{
      \osg@makedelayed*{author}
      \osg@makedelayed*{date}
      \osg@makedelayed*{course}
      \osg@makedelayed*{event}
      \osg@makedelayed*{lehrveranstaltung}
      \osg@makedelayed*{institute}
      \osg@makedelayed*{tucurl}
      \osg@makedelayed{input}
      \RenewDocumentCommand{\SetGlobalClassOptions}{m}
        {\SetKeys[osgbeamer]{##1}}
    }{
      \ClassWarning{\classname}{'\OsgLectConfig' not found.}
    }
    \expandafter\IfFileExists{\OsgLectConfig}{
      \osg@restoredelayed{author}
      \osg@restoredelayed{date}
      \osg@restoredelayed{course}
      \osg@restoredelayed{event}
      \osg@restoredelayed{lehrveranstaltung}
      \osg@restoredelayed{institute}
      \osg@restoredelayed{tucurl}
      \osg@restoredelayed{input}
    }{}
  }
}{}
\undef\logo          % \logo wird wieder freigegeben, so dass Beamer es neu
                     % definieren kann,
%%%%%%%%%%%%%%%%%%%%%%%%%%%%%%%%%%%%%%%%%%%%%%%%%%%%%%%%%%%%%%%%%%%%%%%%%%%%%%
% Verarbeite Optionen
\ProcessKeyOptions[osgbeamer] % Klassenoptionen können die Keys von \SetGlobalClassOptions
                     % überschreiben... 
\osg@forceoptionhook % ... die wiederum durch \EnforceGlobalClassOptions überschrieben
                     % werden können, da im Hook \osg@forceoptionhook Key-Settings
                     % gesammelt worden, und hier *nach* \Processoptions ausgeführt werden.
\IfBool{osgllm}{
  \OsgLectureBuildLoadedTF{
    % New builds use the explicit, job-bound request file.
    \IfBool{osglangset}{}{
      \SetKeyEx{osgbeamer}{lang=\OsgLectureBuildValue{language}}
    }
    \IfBool{osgdoctypeset}{}{
      \SetKeyEx{osgbeamer}{doctype=\OsgLectureBuildValue{doctype}}
    }
  }{
    % Legacy builds continue to derive both values from the jobname.
    \IfBool{osglangset}{}{
      \SetKeyEx{osgbeamer}{lang=\getfourfromjobname}
    }
    \IfBool{osgdoctypeset}{}{
      \SetKeyEx{osgbeamer}{doctype=\getthreefromjobname}
    }
  }
}{}

%%%%%%%%%%%%%%%%%%%%%%%%%%%%%%%%%%%%%%%%%%%%%%%%%%%%%%%%%%%%%%%%%%%%%%%%%%%%%%
% Basisklasse
%
% Je nach Dokumentmodus wird entweder beamer.cls oder scrbook.cls (zusammen mit
% beamerarticle.sty) geladen.
\IfToggle{osgscript}{
  % Script-Mode 
  \LoadClass[twoside]{scrbook}
  %\PreventPackageFromLoading{bookmark}
  %\RequirePackage{hyperref}
  \RequirePackage{beamerarticle}
  \RequirePackage{beamerarticleosg}
  %\RequirePackage{scrhack}
  % 
}{
  \LoadClass{beamer}
  %%% Handout
  \mode<handout>{
    \RequirePackage{pgfpages}
    \expandafter\pgfpagesuselayout\osg@pgfpageoptions
  }
}

% Die Basisklasse (+ggf. beamerarticle) ist geladen.
% Ab jetzt enspricht osgscript = true dem article-Mode,
% false dem presentation-Mode.
% Im Prinzip würde jetzt in das Beamer-Theme einerseits und in
% beamerarticle.sty andererseits verzweigt werden. Jedoch sind einige
% Einstellungen in allen Modi gemeinsam, daher erfolgen 
% sie hier in der Klasse und nicht im Theme/beamerarticle

% Jetzt werden erst mal ein paar Standard-Packages geladen
\RequirePackage{fontspec}
\RequirePackage{xcolor}
\RequirePackage{luacolor}
% Zum Laden von package-lokalen Lua-Modulen wird der Pfad benötigt
\luadirect{
  function osggetpath(type)
    local pkgpath = file.pathpart(file.collapsepath(status.filename, true))
    local allpaths ={
      ['lua']  = file.collapsepath(file.join(pkgpath,'../lua')),
      ['clua'] = file.collapsepath(file.join(pkgpath,'../lua/lib')),
      ['pkg']  = pkgpath
    }
    return allpaths[type]
  end    
  token.set_macro('osgpkgpath', osggetpath('pkg'))
  token.set_macro('osgluapath', osggetpath('lua'))
}
%%%%%%%%%%%%%%%%%%%%%%%%%%%%%%%%%%%%%%%%%%%%%%%%%%%%%%%%%%%%%%%%%%%%%%%%%%%%%%
% Legacy Code
% ===========
%
% Mit der 'legacy'-Option wird die Klasse kompatibler zu den Dokumenten der
% alten OSGBeamer-Klasse gemacht, jedoch fehleranfälliger, da z.T. das erwartete
% Standardverhalten unter LaTeX ausgehebelt wird. Alte Dokumente sollten daher
% so schnell wie möglich angepasst und auf das Setzen der 'legacy'-Option
% verzichtet werden.
%
% In der alten OSGBeamer-Klasse gab es einige Macros und Konstrukte,
% die mittelfristig wegfallen sollen. Mit Hilfe der 'legacy'-Option
% werden einige davon aktiviert, um die Kompatibilität zu erhöhen.
% Beim Einsatz dieser Macros wird eine Warnung gegeben.
\IfBool{osglegacy}{
  \NewDocumentCommand{\DeprecationWarning}{m o}{%
    \IfNoValueTF{#2}{\def\@suf{}}{\def\@suf{, #2}}%
    \ClassWarning{osgbeamer}{Use of '\expandafter\@firstoftwo\string\@#1' is deprecated\@suf}%
  }
  % 
  \NewDocumentCommand\NewCSDeprecatedCommand{m m m O{}}{%
    \ExpandArgs{c}\NewDocumentCommand{#1}{#2}{#3\DeprecationWarning{#1}[#4]}
  }
  \NewDocumentCommand\RenewCSDeprecatedCommand{m m m O{}}{%
    \ExpandArgs{c}\RenewDocumentCommand{#1}{#2}{#3\DeprecationWarning{#1}[#4]}
  }
  % Mit scrlfile (aus dem KOMA-Packet) können unter LuaLaTex und OSGBeamer
  % überflüssige oder störende Pakete gezielt unterdrückt oder angepasst werden.
  \RequirePackage{scrlfile}
  \PreventPackageFromLoading*[
  \ClassWarning{\classname}{Package 'fontenc' is not needed for LuaLaTeX.\MessageBreak
   Please remove \string\usepackage[...]{fontenc}\MessageBreak}
  ]{fontenc}
  \PreventPackageFromLoading*[
  \ClassWarning{\classname}{Package 'inputenc' is not needed for LuaLaTeX.\MessageBreak
   Please remove \string\usepackage[...]{inputenc}\MessageBreak}
  ]{inputenc}
}{}

%%%%%%%%%%%%%%%%%%%%%%%%%%%%%%%%%%%%%%%%%%%%%%%%%%%%%%%%%%%% 
% Umgang mit Beamer-Modi
%%%%%%%%%%%%%%%%%%%%%%%%%%%%%%%%%%%%%%%%%%%%%%%%%%%%%%%%%%%% 
% Hilfsbefehl \osgmakeselectable: Redefiniert ein Macro für die
% Verwendung mit beamer-Aktionen. Optionaler Parameter gibt den
% Default-Mode (*mit* spitzen Klammern) 
% % TODO: Behandlung von optionalen Parametern
\NewDocumentCommand\osgmakeselectable{m o}{
  \expandafter\let\csname\detokenize{#1}@osg@orig\endcsname=#1%
  \renewcommand<>{#1}{%
    \only##1{\csname\detokenize{#1}@osg@orig\endcsname}}%
  \IfNoValueTF{#2}{%
    \renewcommand<>{#1}{\only##1{\csname\detokenize{#1}@osg@orig\endcsname}}%
  }{%
    \renewcommand<>{#1}{\IfBlankTF{##1}{%
        \only#2{\csname\detokenize{#1}@osg@orig\endcsname}%
      }{%
        \only##1{\csname\detokenize{#1}@osg@orig\endcsname}%
      }%
    }%
  } %
}  

\osgmakeselectable{\smallskip}[<presentation>]
\osgmakeselectable{\medskip}[<presentation>]
\osgmakeselectable{\bigskip}[<presentation>]
\osgmakeselectable{\TINY}
\osgmakeselectable{\Tiny}
\osgmakeselectable{\tiny}
\osgmakeselectable{\scriptsize}
\osgmakeselectable{\footnotesize}
\osgmakeselectable{\small}
\osgmakeselectable{\normalsize}
\osgmakeselectable{\large}
\osgmakeselectable{\Large} 
\osgmakeselectable{\LARGE}
\osgmakeselectable{\huge}
\osgmakeselectable{\Huge}

% \vspace kann wegen der Stern-Variante nicht direkt selektierbar
% gemacht werden. Dafür werden die unterliegenden Macros \@vspacer und
% \@vspace gepatcht.
\let\osg@original@@vspacer=\@vspacer
\let\osg@original@@vspace=\@vspace
\renewcommand<>{\@vspacer}{\alt#1{\osg@original@@vspacer}{\@gobble}}
\renewcommand<>{\@vspace}{\alt#1{\osg@original@@vspace}{\@gobble}}

% Newline (\\) ist schon durch Beamer selectierbar gemacht. Allerdings
% definieren einige Umgebungen das Macro neu. Wenn innerhalb einer
% solchen Umgebung die Beamer-Variante gebraucht wird, kann sie mit
% \activateSelDS und \deactivateSelDS aktiviert bzw. wieder
% deaktiviert werden.
\let\beameroriginal@dsl=\\
\let\temp@dsl=\relax
\NewDocumentCommand\activateSelDS{}{
  \temp@dsl=\\%
  \let\\=\beameroriginal@dsl%
}
\NewDocumentCommand\deactivateSelDS{}{
  \let\\=\temp@dsl
}

% \osgpresart erhält zwei Argumente: das erste wird im
% presentation-Mode verwendet, das zweite im article-Mode. Das jeweils
% andere wird verschluckt.
\NewDocumentCommand\osgpresart{+m +m}{%
  \only<presentation>{#1}\only<article>{#2}%
}

%%%%%%%%%%%%%%%%%%%%%%%%%%%%%%%%%%%%%%%%%%%%%%%%%%%%%%%%%%%% 
% Gliederung
%%%%%%%%%%%%%%%%%%%%%%%%%%%%%%%%%%%%%%%%%%%%%%%%%%%%%%%%%%%% 
% Im Original geht Beamer nur von Parts und Sections aus.
% Bei osgbeamer entspricht jede Lecture einem Kapitel.
% Entsprechend werden die Zähler zusammengelegt.
\RequirePackage{aliascnt}
\mode<article>{\newaliascnt{lecture}{chapter}}
\mode<presentation>{\newaliascnt{chapter}{lecture}}
% Die Zähleranzeige der Sections etc. enthält das Kapitel, außer bei 'standalone'
\IfBool{osgoptionstandalone}{%
  \renewcommand{\thesection}{\arabic{section}}
  \renewcommand{\thesubsection}{\arabic{section}.\arabic{subsection}}
  \renewcommand*{\thedefinition}{\arabic{chapter}.\arabic{definition}}
  \renewcommand*{\theexample}{\arabic{chapter}.\arabic{example}}
}{
  \renewcommand{\thesection}{%
    \thechapter.\arabic{section}%
  }
  \renewcommand{\thesubsection}{
    \thesection.\arabic{subsection}
  }  %% TODO: Appendix
}

% Hook zum Ausführen nach dem Titel
\RequirePackage{apptools}
\gdef\osg@aftertitle{}
\newrobustcmd{\AfterTitle}[1]{
  \gappto{\osg@aftertitle}{#1}
}

\RequirePackage{osgbeamerref}      % (Entfernte) Referenzierung
\RequirePackage{osgbeamerlanguage} % Sprachbehandlung
\IfBool{osglegacy}{
  % Babel sollte nicht nochmal geladen werden, um einen Optionsclash zu
  % vermeiden
  \ExpandArgs{cc}\NewCommandCopy{osg@original@ver@babel.sty}{ver@babel.sty}
  % Behauptet, dass babel nie geladen wurde. \PreventPackageFromLoading
  % verhindert zwar Laden, aber kein *erneutes* Laden 
  \csundef{ver@babel.sty}
  % Allerdings muss 'csquotes' von 'babel' wissen, und zwar am Ende der Preamble 
  \AddToHook{begindocument/before}[osgbeamer/csq]{\csletcs{ver@babel.sty}{osg@original@ver@babel.sty}}
  \DeclareHookRule{begindocument/before}{osgbeamer/csq}{before}{csquotes}
  % Jetzt wird das Laden von babel durch den Nutzer verhindert
  \PreventPackageFromLoading*[
  \ClassWarning{\classname}{To avoid an option clash, repeated use of package
    'babel' is skipped. Please remove \string\usepackage[...]{babel}}
  ]{babel}
}{}

% Jetzt wird das OSG-Beamertheme geladen, das wiederum auf dem TUC-Theme aufbaut
\usetheme{osg}
\IfBool{osgbib}{%
  \RequirePackage{osgbeamerbib}      % Literaturreferenzen
}{}

%%%%%%%%%%%%%%%%%%%%%%%%%%%%%%%%%%%%%%%%%%%%%%%%%%%%%%%%%%%% 
% Fonts
%%%%%%%%%%%%%%%%%%%%%%%%%%%%%%%%%%%%%%%%%%%%%%%%%%%%%%%%%%%%
\RequirePackage{fontawesome5}
% LuaLaTeX kann Openface-Fonts, und liefert davon einige mit. Darunter
% auch einige Handschriften-Fonts,so dass eine Installation des
% nonfree Emerald-Packets nicht mehr notwendig ist.
% Allerdings sind die Abstände etwas anders. Es sollte daher bei alten
% Folien/Kapiteln eine Sichtprüfung vorgenommen werden.
\newfontfamily{\scriptfamily}{QTSanDiego}[Scale=1.15]
\DeclareTextFontCommand{\textscript}{\scriptfamily}
\IfBool{osglegacy}{
  \NewCSDeprecatedCommand{ECFJD}{}{\scriptfamily}[use \string\scriptfamily{} or emerald package ]
}{}

% Als Teletype-Schriftart wird in allen Modi DejaVuSansMono gewählt
\mode<presentation>{\setmonofont{DejaVuSansMono}[Scale=MatchLowercase]}
\mode<article>{\setmonofont{DejaVuSansMono}[Scale=MatchLowercase]}
%\newfontfamily\Bera{BeraSansMono}[Scale=0.85]
%\mode<presentation>{\setmonofont{lmmonolt10}[
  % Extension      = .otf,
  % UprightFont    = *-regular,
  % BoldFont       = *-bold,
  % ItalicFont     = *-oblique,
  % BoldItalicFont = *-boldoblique,
  % Scale=MatchLowercase]}
%\usepackage[scaled]{beramono}
%\setmonofont{BeraSansMono}[Scale=MatchLowercase]
%\renewcommand{\ttdefault}{fvm}
% Urls werden in tt-Schrift dargestellt.
\RequirePackage{url}
\def\UrlFont{\ttfamily\upshape}
\Urlmuskip = 0mu plus 1mu

%% aus dem pifont-Paket:
\newcommand{\Pifont}[1]{\fontfamily{#1}\fontencoding{U}\fontseries{m}\fontshape{n}\selectfont}
\newcommand{\Pisymbol}[2]{{\Pifont{#1}\char#2}}
\newcommand{\ding}{\Pisymbol{pzd}}
%%% Farbiger "daraus folgt"-Pfeil
\RequirePackage{xspace}
\newcommand<>{\follows}{\only#1{{\usebeamercolor[fg]{structure}\ding{225}}\xspace}}
\newcommand<>{\lfollows}{\only#1{{\usebeamercolor[fg]{structure}\rotatebox[origin=c]{180}{\ding{225}}\xspace}}}

%%%%%%%%%%%%%%%%%%%%%%%%%%%%%%%%%%%%%%%%%%%%%%%%%%%%%%%%%%%% 
% Fußnoten
%%%%%%%%%%%%%%%%%%%%%%%%%%%%%%%%%%%%%%%%%%%%%%%%%%%%%%%%%%%%
% \afootnote hat nur im Article eine Fußnote gesetzt.
% Das Macro ist nun abgekündigt.
\IfBool{osglegacy}{
  \mode<article>{
    \NewCSDeprecatedCommand{afootnote}{O{}m}{%
      \footnote[#1]{#2}%
    }[use \string\footnote<article>]
  }
  \mode<presentation>{
    \NewCSDeprecatedCommand{afootnote}{O{}m}{}[use \string\footnote<article>]
  }
}{
}

% In Beamer werden beim presentation-Mode nur die Fußnoten ausgeblendet, nicht
% die Fußnotenzeichen. Die OSG-Version stellt das IHMO konsistente
% Verhalten her.
%% Alter Kommentar:
% Außerdem sollen Fußnoten *immer* unten in einem Frame erscheinen,
% nicht am Ende von Blöcken. Die aktuelle Beamer-Variante scheint dies
% (entgegen der Doku) zwar per Default zu machen, zur Sicherheit wird
% aber stets die entsprechende Option übergeben. Ein evtl. doppeltes
% Vorkommen von 'frame' ist unkritisch, so dass das Argument nicht
% untersucht werden muss. 
\NewCommandCopy{\original@footnote}{\footnote}
\RenewDocumentCommand{\footnote}{D<>{all} O{} +m}{%
  \only<#1>{\original@footnote[#2]{#3}}%
}

%\RequirePackage{microtype}


%%%%%%%%%%%%%%%%%%%%%%%%%%%%%%%%%%%%%%%%%%%%%%%%%%%%%%%%%%%% 
% Farben
%%%%%%%%%%%%%%%%%%%%%%%%%%%%%%%%%%%%%%%%%%%%%%%%%%%%%%%%%%%%
\colorlet{infgreen}{tuccolor@if}
\colorlet{tucgreen}{tuccolor@tuc}
\colorlet{tucwarning}{tuccolor@warning}
\colorlet{tucinfo}{tuccolor@info}
\definecolor{lightgreen}{rgb}{0.824,0.878,0.804}
\definecolor{lightgreen2}{rgb}{0.844,0.918,0.834}
\definecolor{darkgreen}{rgb}{0,0.5,0}
\setbeamercolor{outline}{bg=yellow,fg=black}

%% TODO: Farben von Beamer übernehmen
% \mdfdefinestyle{block}{style=basicblock,backgroundcolor=tuccolor@if!10!white,linecolor=tuccolor@if!80!black,frametitlebackgroundcolor=tuccolor@if!90!white,
%   frametitlefont=\sffamily\color{white}}
% \mdfdefinestyle{alertblock}{style=basicblock,backgroundcolor=tuccolor@warning!20!white,linecolor=tuccolor@warning!80!black,frametitlebackgroundcolor=tuccolor@warning!80!white,
%   frametitlefont=\sffamily}
% \mdfdefinestyle{exampleblock}{style=basicblock,backgroundcolor=black!5!white,linecolor=darkgray,frametitlebackgroundcolor=black!65,
%   frametitlefont=\sffamily\color{white}}

% \IfToggle{osgscript}{
% %   \renewenvironment<>{block}[1]{%
% %     \begin{actionenv}#2\parshape0\osg@footnote@init%
% %     \tikzexternaldisable%
% %     \begin{mdframed}[style=block,frametitle=#1]\makeselectable{\\}
% %     }{\end{mdframed}\osg@footnote@do\end{actionenv}}
  
% % \renewenvironment<>{alertblock}[1]{%
% %   \begin{actionenv}#2\parshape0\osg@footnote@init%
% %     \tikzexternaldisable%
% %     \begin{mdframed}[style=alertblock,frametitle=#1]\makeselectable{\\}
% %     }{\end{mdframed}\osg@footnote@do\end{actionenv}
% % }

% % \renewenvironment<>{exampleblock}[1]{%
% %   \begin{actionenv}#2\parshape0\osg@footnote@init%
% %     \tikzexternaldisable%
% %     \begin{mdframed}[style=exampleblock,frametitle=#1]\makeselectable{\\}
% %     }{\end{mdframed}\osg@footnote@do\end{actionenv}
% % }

% \IfToggle{osgscript}{
%   \mdtheorem[style=block]{osgdefinition}{Definition}[chapter]
%   \mdtheorem[style=block]{osgtheorem}{Theorem}[chapter]
% }
% %    %%% TODO: Fußnotenmechanismus anpassen
% %    \renewenvironment{definition}[1][]{%
% %      \parshape0%
% %      \osg@footnote@init%
% %      \tikzexternaldisable%
% %      \begin{osgdefinition}[#1]\makeselectable{\\}
% %      }{\end{osgdefinition}\osg@footnote@do
% %    }
% %    \renewenvironment{theorem}[1][]{%
% %      \parshape0%
% %      \osg@footnote@init%
% %      \tikzexternaldisable%
% %      \begin{osgtheorem}[#1]\makeselectable{\\}
% %      }{\end{osgtheorem}\osg@footnote@do
% %    }
% %    \newcounter{example}[chapter]
% %    \renewenvironment{example}[1][]{%
% %      \refstepcounter{example}%
% %      \parshape0% 
% %      \osg@footnote@init%
% %      \ifstrempty{#1}{\def\extitle{}}{\def\extitle{: #1}}
% %      \begin{exampleblock}{Beispiel \theexample\extitle}\makeselectable{\\}
% %      }{\end{exampleblock}\osg@footnote@do
% %    }
%  }{
%    %\newtheorem{definition}{Definition}[chapter]
%    %\newtheorem{theorem}{Theorem}[chapter]
%    %\newtheorem{lemma}{Lemma}[chapter]
%    \def\th@osgexamplestyle{
%      \def\inserttheoremblockenv{exampleblock}
%    }
%    \def\th@osgblockstyle{ 
%      \def\inserttheoremblockenv{block}
%    }
%    %\theoremstyle{osgexamplestyle}
%    %\newtheorem{example}{Beispiel}[chapter]
%    %\setbeamertemplate{example begin}{}
%    %\setbeamertemplate{example end}{}
% }

\IfBool{osglegacy}{
  \NewCSDeprecatedCommand{tikzanchor}{m m}{%
    \tikz[baseline=(#2.base),remember picture]{\node[outer sep=0,inner sep=0pt](#2){#1};}\xspace%
  }[ use \string\markword]
}{}
%%%%%%%%%%%%%%%%%%%%%%%%%%%%%%%%%%%%%%%%%%%%%%%%%%%%%%%%%%%% 
% Textauszeichnungen
%%%%%%%%%%%%%%%%%%%%%%%%%%%%%%%%%%%%%%%%%%%%%%%%%%%%%%%%%%%%
\NewDocumentCommand{\sbf}{D<>{all} +m}{\structure<#1>{\textbf<#1>{#2}}}
\NewDocumentCommand{\newdef}{D<>{all} +m}{\structure<#1>{\textbf<#1>{#2}}}
\IfBool{osglegacy}{
  \NewCSDeprecatedCommand{ctiny}{D<>{presentation}}{\only<#1>{\tiny}}[use mode spezification]
  \NewCSDeprecatedCommand{cfootnotesize}{D<>{presentation}}{\only<#1>{\footnotesize}}[use mode spezification]
  \NewCSDeprecatedCommand{cscriptsize}{D<>{presentation}}{\only<#1>{\scriptsize}}[use mode spezification]
  \NewCSDeprecatedCommand{csmall}{D<>{presentation}}{\only<#1>{\small}}[use mode spezification]
  \NewCSDeprecatedCommand{cnormalsize}{D<>{presentation}}{\only<#1>{\normalsize}}[use mode spezification]
}{}
\RequirePackage{lua-ul}
\let\uline=\underLine
% \outline unterscheidet sich in den Modi: In der Presentation ist es
% Textmarker, im Article unterstrichen
%\mode<article>{
%   \newunderlinetype\beginUnderWavy[\number\dimexpr1ex]{\cleaders\hbox{%
%       \setlength\unitlength{.3ex}%
%       \begin{picture}(4,0)(0,1)
%           \thicklines
%           % \color{red}%
%           \qbezier(0,0)(1,1)(2,0)
%           \qbezier(2,0)(3,-1)(4,0)
%       \end{picture}%
%     }}
%   \NewDocumentCommand\outline{D<>{all}+m}{\alt<#1>{\underline{#2}}{#2}}
% }
%\mode<presentation>{
  \NewDocumentCommand\outline{D<>{all}+m}{\alt<#1>{\highLight{#2}}{\highLight[bg]{#2}}} % Die phantom \highLight[bg]-Version sorgt für
  % identische Abstände.
  % ToDo: \highLight anpassen, dass der Durchschuss nicht geändert wird.
%}
% Auch \stess (war: \grave) ist modespezifisch (ToDo: Soll das so bleiben?)
\mode<presentation>{\NewDocumentCommand\stress{D<>{all} m}{{\only<#1>{\color{structure}}{#2}}}}
\mode<article>{\NewDocumentCommand\stress{D<>{all} m}{\emph<#1>{#2}}}
% \grave wird jetzt durch amsmath als mathematischer Akzent definiert.
% Mit legacy-Option wird dies überschrieben
\IfBool{osglegacy}{
  \AddToHook{begindocument}{
    \RenewCSDeprecatedCommand{grave}{d<>m}{%
      \IfNoValueTF{#1}{%
        \textcolor{structure}{#2}%
      }{%
        \only<#1>{\textcolor{structure}}{#2}%
      }}[use \string\stress]
  }
}{}
\RequirePackage{calc}
\NewDocumentCommand\sourceref{sO{\ldeen{Quelle}{Source}} m O{style=sameline}}{%
  \IfBooleanTF{#1}{
    \begin{flushright}
        \fontsize{6pt}{2pt}\selectfont\textbf{#2:}\ \stress{#3}
    \end{flushright}
  }{
    \begin{description}[#4]\fontsize{6pt}{2pt}\selectfont
     \item[\textbf{#2:}] \stress{#3}
    \end{description}
  }
}
  %\par{\fontsize{6pt}{2pt}\selectfont\textbf{#1}:: \stress{#2}}}
%%%%%%%%%%%%%%%%%%%%%%%%%%%%%%%%%%%%%%%%%%%%%%%%%%%%%%%%%%%% 
% Columns
%%%%%%%%%%%%%%%%%%%%%%%%%%%%%%%%%%%%%%%%%%%%%%%%%%%%%%%%%%%%
%% Vereinfachtes Spaltenmanagement für zwei Spalten:
%% twocolumns nimmt zwei optionale Argumente (Reihenfolge egal):
%% - vertikale Ausrichtung {c|b|t|T}
%% - Anteil, den die erste Spalte einnimmt. 

\newlength{\osgleftcolwid}
\newlength{\osgrightcolwid}

\ExplSyntaxOn

\NewDocumentCommand{\IfNumberTF}{mmm}
 {%
  % regex from https://stackoverflow.com/a/23872060/923955
  % with modifications for expl3 and for adding a unit
  \regex_match:nnTF
   { \A [+\-]? ((\d+(\.\d*)?)|(\.\d+)) \Z} % regex
   { #1 } % test string
   { #2 } % true text
   { #3 } % false text
 }

\ExplSyntaxOff

\NewDocumentCommand{\twocolparam}{>{\SplitList{,}} m}{%
  \def\osgtcfcwid{0.5}%    Falls keine Länge angegeben wird
  \def\osgtcalignparam{c}% Falls kein Alignmentparameter angegeben wird
  \def\parseandsetarguments##1{%
    \IfNumberTF{##1}{%
      \def\osgtcfcwid{##1}%
    }{%
      \def\osgtcalignparam{##1}%
    }%
  }%
  \ProcessList{#1}{\parseandsetarguments}% W
}


\newenvironment<>{twocolumns}[1][t,0.5]{%
  \twocolparam{#1}% 
  \setlength{\osgleftcolwid}{\fpeval{\osgtcfcwid*0.975}\textwidth}%
  \setlength{\osgrightcolwid}{0.975\textwidth}%
  \addtolength{\osgrightcolwid}{-\osgtcfcwid\textwidth}%
  \alt#2{\def\nextcolumn{\column{\osgrightcolwid}}}{\def\nextcolumn{\par}}%
  \begin{columns}#2[\osgtcalignparam]\column{\osgleftcolwid}%
    }{%
  \end{columns}%\undef{\osgleftcolwid}\undef{\osgrightcolwid}
}

%%%%%%%%%%%%%%%%%%%%%%%%%%%%%%%%%%%%%%%%%%%%%%%%%%%%%%%%%%%% 
% Deko
%%%%%%%%%%%%%%%%%%%%%%%%%%%%%%%%%%%%%%%%%%%%%%%%%%%%%%%%%%%%
\RequirePackage{twemojis}
\newcommand{\makesmiley}[2]{\csgdef{osgsmiley#1}{\raisebox{-.2ex}{\texttwemoji{#2}}}}
\makesmiley{:-)}{1f60a} % happy 
\makesmiley{:-D}{1f600} % laugh
\makesmiley{:-|}{1f62c}% angry
\makesmiley{:-()}{1f615}% sad 
\makesmiley{:-P}{1f61b} % tongue
\makesmiley{;-)}{1f609} % wink 
\makesmiley{:-o}{1f62e} % astonished
\makesmiley{B-)}{1f60e} % cool
\newcommand{\smiley}[1]{\csuse{osgsmiley#1}}
\let\osgsmiley=\smiley

% Erstellt einen tikz-Anker (Name ist erster Parameter) um den
% gegebenen Text (2. Parameter), auf den später referenziert werden
% kann.
\newif\ifosgmarkmm%
\newcommand{\markword}[2]{%
  \relax\ifmmode% Ist die Markierung im Mathe-Modus?
  \osgmarkmmtrue\else\osgmarkmmfalse\fi%
  \ifdef{\tikzexternalize}{\tikzset{external/export next=false}}{}% deaktiviere tikz-Externalisierung
  \tikz[remember picture, baseline]\node[anchor=text,inner sep=0pt](#1){\ifosgmarkmm$#2$\else#2\fi};}

%%%%%%%%%%%%%%%%%%%%%%%%%%%%%%%%%%%%%%%%%%%%%%%%%%%%%%%%%%%% 
% Grafiken
%%%%%%%%%%%%%%%%%%%%%%%%%%%%%%%%%%%%%%%%%%%%%%%%%%%%%%%%%%%% 
\RequirePackage{pgf}
% Horizontal zentriertes Bild. Der optionale Parameter enthält die
% Breite (default: 0.9\textwidth). Es können zwei durch Komma
% getrennte Werte angegeben werden: Der erste gilt dann für
% den presentation-Mode, der zweite für den article-Mode.
\newcommand{\osg@centerpic}[2][]{
  \def\parsewidth##1,##2,{\dimdef{\pw}{##1}\dimdef\aw{##2}}
  \edef\x{\noexpand\in@{,}{#1}}\x
  \ifin@
    \parsewidth#1,\relax
  \else
    \dimdef\pw{#1}
    \dimdef\aw{#1}
   \fi
  \only<presentation>{\centerline{\includegraphics[width=\pw]{#2}}}
  \only<article>{\centerline{\includegraphics[width=\aw]{#2}}}
}
\NewDocumentCommand\centerpic{D<>{all} O{0.9\textwidth} m}{%
  \only<#1>{\osg@centerpic[#2]{#3}}
}
%%%%%%%%%%%%%%%%%%%%%%%%%%%%%%%%%%%%%%%%%%%%%%%%%%%%%%%%%%%% 
% Anhang
%%%%%%%%%%%%%%%%%%%%%%%%%%%%%%%%%%%%%%%%%%%%%%%%%%%%%%%%%%%%
% Anhang gibt es in zwei "Geschmacksrichtungen". Zum einen kann eine einzelne
% Vorlesung einen Anhang haben, zum anderen können ganze Vorlesungen zum
% Anhang erklärt werden, was i.d.R. nur im Zusammenhang mit Skripten sinnvoll
% ist.
% Um beide Varianten zu unterscheiden wird geprüft, ob die \appendix-Anweisung
% in der Preamble oder im Dokument steht.
\AtAppendix{%
  \ifx\@onlypreamble\@notprerr
   \renewcommand{\thesection}{\appendixname~\Alph{section}:}
   \renewcommand{\thesubsection}{\thesection.\arabic{subsection}}
    \setcounter{section}{0}
  \else
    \renewcommand{\thesection}{\thechapter.\arabic{section}}
   \renewcommand{\thesubsection}{\thesection.\arabic{subsection}}
  \fi
}


%%%%%%%%%%%%%%%%%%%%%%%%%%%%%%%%%%%%%%%%%%%%%%%%%%%%%%%%%%%% 
% Anmerkungen
%%%%%%%%%%%%%%%%%%%%%%%%%%%%%%%%%%%%%%%%%%%%%%%%%%%%%%%%%%%%
\ExplSyntaxOn
\NewDocumentEnvironment{specialitemize}{ O{} +b }
 {
  % do the setup
  \keys_set:nn { osg/itemize } { #1 }
  % split the contents at \item
  \seq_set_split:Nnn \l_osg_itemize_input_seq { \item } { #2 }
  % remove the first (empty) item
  \seq_pop_left:NN \l_osg_itemize_input_seq \l_tmpa_tl
  % issue the preamble
  \tl_use:N \l_osg_itemize_pre_tl
  % adorn the items
  \seq_set_map:NNn
    \l_osg_itemize_output_seq
    \l_osg_itemize_input_seq
    { \__osg_itemize_do:n { ##1 } }
  % output the items, separated by the chosen separator
  \seq_use:NV \l_osg_itemize_output_seq \l_osg_itemize_sep_tl
  % issue the postamble
  \tl_use:N \l_osg_itemize_post_tl
 }
 {}

\seq_new:N \l_osg_itemize_input_seq
\seq_new:N \l_osg_itemize_output_seq
\cs_generate_variant:Nn \seq_use:Nn { NV }

\keys_define:nn { osg/itemize }
 {
  pre    .tl_set:N  = \l_osg_itemize_pre_tl,
  post   .tl_set:N  = \l_osg_itemize_post_tl,
  sep    .tl_set:N  = \l_osg_itemize_sep_tl,
  action .code:n    = \cs_set_eq:NN \__osg_itemize_do:n #1,
  action .initial:n = \use:n,
}

%% Überflüssig wg. TeX3
% \ProvideDocumentCommand \NewEnvironmentCopy {mm} {
%     \expandafter \NewCommandCopy \csname#1\expandafter\endcsname \csname#2\endcsname
%     \expandafter \NewCommandCopy \csname end#1\expandafter\endcsname \csname end#2\endcsname
% }

% \ProvideDocumentCommand \RenewEnvironmentCopy {mm} {
%     \expandafter \RenewCommandCopy \csname#1\expandafter\endcsname \csname#2\endcsname
%     \expandafter \RenewCommandCopy \csname end#1\expandafter\endcsname \csname end#2\endcsname
% }
\ExplSyntaxOff

\NewDocumentCommand{\spacesaveitem}{m}{%
  \textbullet\ \detokenize{#1}\par%
}
\NewDocumentCommand{\quotepar}{m}{%
  ``#1''\par
}
%\RequirePackage[luatex,space=true]{accsupp}
\mode<article>{\def\Note#1{}}
\mode<presentation>{
  
   \let\osg@original@note=\note
  % \let\osg@original@beginitem=\beginitem

  \NewDocumentEnvironment{noteitemize}{}
  {
    \begin{specialitemize}[action=\spacesaveitem]
      }{
    \end{specialitemize}
  } 
  \NewEnvironmentCopy{originalitemize}{itemize}
  \RenewDocumentCommand{\note}{m}{
    \osg@original@note{
     \RenewEnvironmentCopy{itemize}{noteitemize}
     #1
     \RenewEnvironmentCopy{itemize}{originalitemize}
     }
   }
  % \def\Note#1{\beamer@inframenote[item]{#1}}
}

%%%%%%%%%%%%%%%%%%%%%%%%%%%%%%%%%%%%%%%%%%%%%%%%%%%%%%%%%%%% 
% Titeldaten
%%%%%%%%%%%%%%%%%%%%%%%%%%%%%%%%%%%%%%%%%%%%%%%%%%%%%%%%%%%%
\mode<article>{
  % Erlaube optionale Argumente für \author und \date
  \global\let\@osg@original@author=\author
  \global\let\@osg@original@title=\title
 
  \RenewDocumentCommand\author{o m}{
    \IfBlankTF{#1}{
      \gdef\osgshortauthor{#2}
    }{
      \gdef\osgshortauthor{#1}
    }    
    \gdef\osgauthor{#2}
    \gdef\@author{#2}
  }
  \RenewDocumentCommand\title{o m}{
    \IfBlankTF{#1}{
      \gdef\osgshorttitle{#2}
    }{
      \gdef\osgshorttitle{#1}
    }
    \gdef\osgtitle{#2}
    \gdef\@title{#2}
  }
  \gdef\lehrveranstaltung{\title}
  \RenewDocumentCommand\date{o m}{
    \IfBlankTF{#1}{
      \gdef\osgshortdate{#2}
    }{
      \gdef\osgshortdate{#1}
    }
    \gdef\@date{#2}
    \gdef\osgdate{#2}
  }
  \gdef\lehrveranstaltung{\title}
}
\mode<presentation>{
  \gdef\lehrveranstaltung{\subtitle}
}

\let\event=\lehrveranstaltung
\let\course=\lehrveranstaltung

\osg@titlematterhook

\let\osg@origina@Roman=\Roman
\renewcommand{\Roman}[1]{%
  \ifnum\value{#1}=0\relax N\else\osg@origina@Roman{#1}\fi}

\NewDocumentCommand\DebugFont{}{Family: \f@family\  with size: \f@size.}
% %%% Local Variables:
% %%% mode: latex
% %%% End:
%</osgbeamer>
